\documentclass[11pt,reqno,psamsfonts]{amsart}
\pdfoutput=1

\usepackage{amssymb}
\usepackage{amsthm}
\usepackage{amsmath}
\usepackage{latexsym}
\usepackage[T1]{fontenc}
\usepackage[utf8]{inputenc}
\usepackage[russian, french, 
english]{babel}
\newcommand\hmmax{0}
\newcommand\bmmax{0}
\usepackage{graphicx}
\usepackage{wrapfig}
\usepackage{mathtools}
\usepackage{amsbsy}
\usepackage[inline]{enumitem}
\usepackage{mathrsfs}
\usepackage{array}
\usepackage{multicol}
\usepackage{stmaryrd}
\usepackage{cancel}
\usepackage{lmodern}
\usepackage{mathabx}
\usepackage{upgreek}
\usepackage{titlesec}
\usepackage{titletoc}
\usepackage[spacing=true,kerning=true,babel=true,tracking=true]{microtype}
\usepackage[shortcuts]{extdash}
\usepackage[foot]{amsaddr}
\usepackage[left=1in,right=1in,top=1in,bottom=1in,bindingoffset=0cm]{geometry}
\usepackage{bm}
\usepackage{centernot}
\usepackage{mdframed}
\usepackage[hidelinks]{hyperref}
\usepackage{xspace}
%\usepackage{eucal,eufrak}
\usepackage[most]{tcolorbox}
\usepackage{framed}
\usepackage[labelfont=bf]{caption}
\captionsetup{font=small}
\usepackage{tikz}
\usetikzlibrary{math}
\usetikzlibrary{shapes,snakes}
\usetikzlibrary{arrows.meta}
\usetikzlibrary{decorations.pathmorphing}
\usetikzlibrary{patterns}
\usetikzlibrary{positioning,calc,intersections}
\usepackage{tikzsymbols}
\usepackage{float}

\DeclareFontFamily{U}{skulls}{}
\DeclareFontShape{U}{skulls}{m}{n}{ <-> skull }{}
\newcommand{\skull}{\text{\usefont{U}{skulls}{m}{n}\symbol{'101}}}

\usepackage[
    sortcites,
    backend=biber, style=alphabetic, sorting=nyt, maxnames=100,backref=true]{biblatex}
\addbibresource{main.bib}

\newtheoremstyle{bfnote}%
{}{}%
{\slshape}{}%
{\bfseries}{\bfseries.}%
{ }%
{\thmname{#1}\thmnumber{ #2}\thmnote{ \ep{\normalfont{}#3}}}

\theoremstyle{bfnote}
\newtheorem{theo}{Theorem}[section]
\newtheorem{theorem}[theo]{Theorem}
\newtheorem*{theo*}{Theorem}
\newtheorem{prop}[theo]{Proposition}
\newtheorem{lemma}[theo]{Lemma}
\newtheorem{claim}[theo]{Claim}
\newtheorem{subclaim}[theo]{Subclaim}
\newtheorem{corl}[theo]{Corollary}
\newtheorem{conj}[theo]{Conjecture}
\newtheorem{fact}[theo]{Fact}
\newtheorem{obs}[theo]{Observation}
\newtheorem*{corl*}{Corollary}

\theoremstyle{definition}
\newtheorem{defn}[theo]{Definition}
\newtheorem*{defn*}{Definition}
\newtheorem{exmp}[theo]{Example}
\newtheorem{exmps}[theo]{Examples}
\newtheorem{remk}[theo]{Remark}
\newtheorem{question}[theo]{Question}
\newtheorem{remks}[theo]{Remarks}
\newtheorem*{remks*}{Remarks}
\newtheorem*{exmp*}{Example}
\newtheorem{prob}[theo]{Problem}

\theoremstyle{remark}
\newtheorem*{ques*}{Question}
\newtheorem*{remk*}{Remark}

\newcommand*{\myproofname}{Proof}
\newenvironment{claimproof}[1][\myproofname]{\begin{proof}[#1]\renewcommand*{\qedsymbol}{\(\boxtimes\)}}{\end{proof}}

\makeatletter
\newcommand{\neutralize}[1]{\expandafter\let\csname c@#1\endcsname\count@}
\makeatother

\newenvironment{theocopy}[1]
{\renewcommand{\thetheo}{\ref{#1}}%
	\neutralize{theo}\phantomsection
	\begin{theo}}
	{\end{theo}}

\newenvironment{lemmacopy}[1]
{\renewcommand{\thetheo}{\ref{#1}}%
	\neutralize{theo}\phantomsection
	\begin{lemma}}
	{\end{lemma}}

    \newenvironment{corlcopy}[1]
{\renewcommand{\thetheo}{\ref{#1}}%
	\neutralize{theo}\phantomsection
	\begin{corl}}
	{\end{corl}}

\renewcommand{\qedsymbol}{$\blacksquare$}
\newcommand{\0}{\varnothing}
\newcommand{\set}[1]{\{#1\}}
\renewcommand{\C}{\mathbb{C}}
\newcommand{\Q}{\mathbb{Q}}
\renewcommand{\P}{\mathbb{P}}
\newcommand{\E}{\mathbb{E}}
\renewcommand{\epsilon}{\varepsilon}
\newcommand{\eps}{\epsilon}
\renewcommand{\phi}{\varphi}
\renewcommand{\theta}{\vartheta}
\renewcommand{\leq}{\leqslant}
\renewcommand{\geq}{\geqslant}
\newcommand{\defeq}{\coloneqq}
\newcommand{\im}{\mathsf{im}}
\newcommand{\bemph}[1]{{\normalfont#1}}
\newcommand{\ep}[1]{\bemph{(}#1\bemph{)}}
\newcommand{\eb}[1]{\bemph{[}#1\bemph{]}}
\newcommand{\red}[1]{{\color{red}#1}}
\newcommand{\blue}[1]{{\color{blue}#1}}
\newcommand{\todo}[1]{\red{\textbf{\upshape{[#1]}}}}
\newcommand{\emphd}[1]{{\fontseries{b}\selectfont\textsf{#1}}}
\newcommand{\acts}{\mathrel{\reflectbox{$\righttoleftarrow$}}}
\newcommand{\G}{\Gamma}
\newcommand{\pto}{\dashrightarrow}
\newcommand{\Free}{\mathsf{Free}}
\newcommand{\powerset}[1]{\mathcal{P}(#1)}
\newcommand{\symdif}{\mathbin{\triangle}}
\newcommand{\rest}[2]{{{#1}\vert_{#2}}}
\newcommand{\dist}{\mathsf{dist}}
\newcommand{\fs}[2]{{[{#1}]^{#2}}}
\newcommand{\finset}[1]{\fs{{#1}}{{<\infty}}}
\newcommand{\N}{\mathbb{N}}
\newcommand{\Z}{\mathbb{Z}}
\newcommand{\R}{\mathbb{R}}
\newcommand{\dom}{\mathsf{dom}}
\newcommand{\asi}{\mathsf{asi}}
\newcommand{\casi}{\asi_{\mathsf{c}}}
\newcommand{\asdim}{\mathsf{asdim}}
\newcommand{\supp}{\mathsf{supp}}
\newcommand{\Leb}{\mathcal{L}}
\newcommand{\Cay}{\mathsf{Cay}}
\newcommand{\Sch}{\mathsf{Sch}}
\newcommand{\Sub}{\mathsf{Sub}}
\newcommand{\In}{\mathsf{In}}
\newcommand{\Out}{\mathsf{Out}}
\newcommand{\iden}{\mathbf{1}}
\newcommand{\sr}{amply syndetic\xspace}
\newcommand{\Sr}{Amply syndetic\xspace}
\newcommand{\Sep}{\mathsf{Sep}}
\newcommand{\Synd}{\mathsf{Synd}}
\newcommand{\Min}{\mathsf{Min}}
\newcommand{\Vis}{\mathsf{Vis}}
\newcommand{\Ret}{\mathsf{Ret}}


\numberwithin{equation}{section}

\newenvironment{scproof}[1][]{\begin{proof}[\textsc{\upshape{Proof}}#1]}{\end{proof}}

\titleformat{\section}[block]{\large\bfseries\sffamily}{\thesection.}{1ex}{}
\titleformat{\subsection}[block]{\bfseries\sffamily}{\thesubsection.}{1ex}{}
\titleformat{\subsubsection}[runin]{\bfseries}{\bfseries\upshape\thesubsubsection.}{1ex}{}[.]

\titlespacing*{\section}{0pt}{*3}{*1}
\titlespacing*{\subsection}{0pt}{*3}{*1}
\titlespacing*{\subsubsection}{0pt}{*1.5}{*1}


 \titlecontents{section}
 [1.5em] %
 {\smallskip}
 {\bfseries\thecontentslabel\hspace{1.02em}}%\thecontentslabel
{\bfseries}%\hspace*{2.32em}}
 {\,\,\titlerule*[0.77pc]{}\bfseries\contentspage}
\titlecontents{subsection}
 [4em] %
 {\smallskip}
 {\thecontentslabel\hspace{1.02em}}%\thecontentslabel
{\hspace*{2.32em}}
 {\,\,\titlerule*[0.77pc]{.}\contentspage}


\renewcommand{\mkbibnamefirst}[1]{\textsc{#1}}
\renewcommand{\mkbibnamelast}[1]{\textsc{#1}}
\renewcommand{\mkbibnameprefix}[1]{\textsc{#1}}
\renewcommand{\mkbibnameaffix}[1]{\textsc{#1}}
\renewcommand{\finentrypunct}{}

\renewbibmacro{in:}{}

\renewbibmacro*{volume+number+eid}{%
	\printfield{volume}%
	%  \setunit*{\adddot}% DELETED
	\setunit*{\addnbspace}% NEW (optional); there's also \addnbthinspace
	\printfield{number}%
	\setunit{\addcomma\space}%
	\printfield{eid}}
%\DeclareFieldFormat[article]{number}{\mkbibparens{#1}}

\DeclareFieldFormat[article]{volume}{\textbf{#1}\space}
\DeclareFieldFormat[article]{number}{\mkbibparens{#1}}

\DeclareFieldFormat{journaltitle}{#1,}
\DeclareFieldFormat[thesis]{title}{\mkbibemph{#1}\addperiod}
\DeclareFieldFormat[article, unpublished, thesis]{title}{\mkbibemph{#1},}
\DeclareFieldFormat[book]{title}{\mkbibemph{#1}\addperiod}
\DeclareFieldFormat[unpublished]{howpublished}{#1, }

\DeclareFieldFormat{pages}{#1}

\DeclareFieldFormat[article]{series}{Ser.~#1\addcomma}

\renewcommand*{\bibfont}{\small}

\setlength{\footskip}{1.5\baselineskip}

\setlist{topsep=3pt,itemsep=3pt}

\DeclareRobustCommand{\cyrins}[1]{%
	\begingroup\fontfamily{cmr}%
	\foreignlanguage{russian}{#1}%
	\endgroup
}

\pagestyle{plain}


\title{\sffamily Flows with minimal subdynamics}
\date{}


\author{Anton~Bernshteyn}
\address{\normalfont (AB) Department of Mathematics, University of California, Los Angeles, CA, USA}
\email{bernshteyn@math.ucla.edu}

\author{Joshua~Frisch}
\address{\normalfont (JF) 
Department of Mathematics, University of California, San Diego, CA, USA}
\email{jfrisch@ucsd.edu}

\thanks{AB's research is partially supported by the NSF CAREER grant DMS-2528522 and the Sloan Research Fellowship (2025).
JF's research is partially supported by the NSF grant DMS-2348981.}

\begin{document}

%\vspace*{0pt}

\maketitle

\begin{abstract}
    Let $\G$ be a countably infinite discrete group. A $\G$-flow $X$ (i.e., a nonempty compact Hausdorff space equipped with a continuous action of $\G$) is called {$S$-minimal} for a subset $S \subseteq \G$
    if the partial orbit $S \cdot x$ is dense for every point
    $x \in X$. We show that for any countable family $(S_n)_{n \in \N}$ of infinite subsets of $\G$, there exists a free $\G$-flow $X$ that is $S_n$-minimal for all $n \in \N$; additionally, $X$ can be taken to be a subflow of $2^\G$. This vastly generalizes a result of Frisch, Seward, and Zucker,
    in which each $S_n$ is required to be a normal subgroup of $\G$. As a corollary, we show that for a given Polish $\G$-flow $X$, there exists a free $\G$-flow $Y$ disjoint from $X$ in the sense of Furstenberg if and only if $X$ has no wandering points. This completes a line of inquiry started by Glasner, Tsankov, Weiss, and Zucker. 
    As another application, we strengthen some of the results of Gao, Jackson, Krohne, and Seward on the structure of Borel complete sections. For example, we show that if $B$ is a Borel complete section in the free part of $2^\G$, then every union of sufficiently many shifts of $B$ contains an orbit (previously, this was only known for open sets $B$). Although our main results are purely dynamical, their proofs rely on recently developed machinery from descriptive set-theoretic combinatorics, namely the 
    asymptotic separation index introduced by Conley, Jackson, Marks, Seward, and Tucker-Drob and its links to the Lov\'asz Local Lemma.
\end{abstract}

\section{Introduction}\label{sec:intro}

    \subsection{Existence of flows and subshifts with minimal subdynamics}

    Throughout, $\G$ is an infinite
    countable discrete group.
    A \emphd{$\G$-flow} is a nonempty compact Hausdorff space $X$ equipped with a continuous action $\G \acts X$. A \emphd{subflow} of a $\G$-flow $X$ is a nonempty closed $\G$-invariant subset $Y \subseteq X$. A $\G$-flow $X$ is \emphd{minimal} if it has no subflows other than $X$ itself; equivalently, $X$ is minimal if the orbit $\G \cdot x$ of every point $x \in X$ is dense in $X$. A routine argument using Zorn's lemma shows that every $\G$-flow has a minimal subflow. That being said, finding explicit examples of minimal $\G$-flows with desirable properties is often a challenging task, leading to many intriguing open questions. In this paper we are interested in the following class of problems, which Seward, Zucker, and the second named author \cite{FSZ} called \emph{minimal subdynamics}:

    \begin{tcolorbox}[width=\linewidth, sharp corners=all, colback=white!95!black, boxrule=0pt,colframe=white] 
        Given a family $(\Delta_i)_{i \in I}$ of infinite subgroups of $\G$, when does there exist a free $\G$-flow $X$ such that the induced action $\Delta_i \acts X$ makes it a minimal $\Delta_i$-flow for all $i \in I$?
    \end{tcolorbox}

    Here we say that a group action $\G \acts X$ is \emphd{free} if the stabilizer of every point $x \in X$ is trivial, i.e., if $\gamma \cdot x \neq x$ for all $x \in X$ and $\gamma \in \G \setminus \set{\iden}$, where $\iden$ is the identity element of $\G$.

    The study of dynamical systems via their ``subdynamics,'' i.e., by restricting the action to a subgroup, 
    has long been a theme 
    in topological dynamics and ergodic theory. A classical example is the notion of a \emph{totally minimal} homeomorphism, i.e., a homeomorphism $T \colon X \to X$ such that for all nonzero $n \in \Z$, the action $\langle T^n \rangle \acts X$ of the group generated by $T^n$ is minimal. Totally minimal homeomorphisms were introduced by Hedlund in \cite{HedlundSturmian} (under the name ``powerfully minimal'').  
    This definition is extended to $\G$-flows for general $\G$ by saying that a $\G$-flow is \emph{totally minimal} if it is minimal as a $\Delta$-flow for all subgroups $\Delta \leq \G$ of finite index \cite[Defn.~2.27]{GottHed}. \emph{Totally ergodic} is an analogous concept in ergodic theory \cite[Defn.~10.3]{RSErgodic}. Various aspects of subdynamics for $\Z^d$-flows, such as the notion of \emph{expansive subdynamics} introduced 
    by Boyle and Lind in \cite{Expansive1}, have been extensively studied; see, e.g., \cite{Expansive2,Expansive3,Expansive4,Expansive5,Expansive6,Expansive7,Expansive8,Expansive9,Rosenthal}  for a small sample of the literature on the subject.
    
    Generalizing an earlier result of Zucker \cite{Zucker}, Seward, Zucker, and the second named author gave the following partial answer to the minimal subdynamics question:

    \begin{theo}[{JF--Seward--Zucker \cite[Thm.~0.3]{FSZ}}]\label{theo:FSZ}
        If $(\Delta_n)_{n \in \N}$ is a countable family of infinite {normal} subgroups of $\G$, then
        there is a free $\G$-flow that is minimal as a $\Delta_n$-flow for all $n \in \N$.
    \end{theo}

    The normality assumption on the subgroups $\Delta_n$ is quite restrictive, but it appears difficult to remove it with prior techniques.
    Indeed, even the following very modest special case of the general problem remained open until now:

    \begin{quote}
        {Let $\mathbb{F}_2$ be the free group of rank $2$. Is it true that for every non-identity element $\gamma$ of $\mathbb{F}_2$, % and let $\gamma$ be a non-identity element of $\mathbb{F}_2$. Does
        there exists a free $\mathbb{F}_2$-flow that is minimal as a $\langle \gamma \rangle$-flow?}
    \end{quote}

    In this paper we introduce a novel approach that combines tools from topological dynamics, descriptive set theory, and combinatorics and allows us to completely remove the normality assumption in Theorem~\ref{theo:FSZ}. Not only that,  %considerably broaden the scope of Theorem~\ref{theo:FSZ} by removing the assumption that each $\Delta_n$ is a \emph{normal} subgroup of $\G$. In fact,
    we do not even need to assume that each $\Delta_n$ is a \emph{subgroup} of $\G$. %Namely, we say that a $\G$-flow $X$ is \emphd{$S$-minimal} for a \emph{subset} $S \subseteq \G$ if the partial orbit $S \cdot x$
    %\[
    %    S \cdot x \,\defeq\, \set{\sigma \cdot x \,:\, \sigma \in S}
    %\]
    %is dense for every $x \in X$. Here is our main result:

    \begin{defn}[$S$-minimal $\G$-flows]\label{defn:S-minimal}
        We say that a $\G$-flow $X$ is \emphd{$S$-minimal} for a {subset} $S \subseteq \G$ if the partial orbit $S \cdot x$ is dense in $X$ for every point $x \in X$.
    \end{defn}

    %Here is our main result:

    \begin{tcolorbox}
    \begin{theo}\label{theo:main}
        If $(S_n)_{n \in \N}$ is a countable family of infinite subsets of $\G$, then there exists a free $\G$-flow %$X$
        that is $S_n$-minimal for all $n \in \N$.
        %
        %Let $(S_n)_{n \in \N}$ be a countable family of infinite subsets of $\G$.
        %
        %\begin{enumerate}[label=\ep{\normalfont\roman*}]
        %    \item There exists a free $\G$-flow $X$ that is $S_n$-minimal for all $n \in \N$.
        %
        %    \item\label{item:subshift} Furthermore, such $X$ can be chosen to be a subflow of $2^\G$.
        %\end{enumerate}
        %
        %, then there exists a free $\G$-flow $X$ that is $S_n$-minimal for all $n \in \N$. Furthermore, 
    \end{theo}
    \end{tcolorbox}

     %. It is interesting that most results in these papers obtain free subshifts that are ``large'' in various senses, whereas our Theorem~\ref{theo:main} yields a free subshift $X$ that is somehow ``small.''

    %Since $\G$ has only countably many finitely generated subgroups, Theorem~\ref{theo:main} yields a $\G$-flow that is  $\Delta$-minimal for every infinite finitely generated subgroup $\Delta \leq \G$. f $S \subseteq S' \subseteq \G$, then every $S$-minimal $\G$-flow is also $S'$-minimal. Note that  the following:

    %Recall that a group is called \emphd{locally finite} if all its finitely generated subgroups are finite.
    
    The following statements are immediate consequences of Theorem~\ref{theo:main}:

    \begin{tcolorbox}
    \begin{corl}\label{corl:fg_subgroups}
        There exists a free $\G$-flow %subshift $X \subset 2^\G$
        that is $\Delta$-minimal with respect to
        every non-locally finite subgroup $\Delta \leq \G$.
    \end{corl}
    \end{tcolorbox}
    \begin{scproof}
        %Recall that a group is {locally finite} if all its finitely generated subgroups are finite. 
        Note that if $S \subseteq S' \subseteq \G$, then every $S$-minimal $\G$-flow is also $S'$-minimal. Hence, it suffices to find a free $\G$-flow %subshift $X \subset 2^\G$
        that is $\Delta$-minimal for every infinite {finitely generated} subgroup $\Delta \leq \G$. Such a $\G$-flow %subshift
        exists by Theorem~\ref{theo:main} as $\G$ has only countably many finitely generated subgroups. %, we may apply Theorem~\ref{theo:main} to obtain a free $\G$-flow $X$ that is $\Delta$-minimal for every infinite finitely generated subgroup $\Delta \leq \G$. 
    \end{scproof}

    %Recall that a group is \emphd{locally finite} if all its finitely generated subgroups are finite. If $S \subseteq S' \subseteq \G$, then every $S$-minimal $\G$-flow is also $S'$-minimal. In particular, the $\G$-flow given by Corollary~\ref{corl:fg_subgroups} is $\Delta$-minimal for every non-locally finite subgroup $\Delta \leq \G$. Therefore, if $\G$ has no infinite locally finite subgroups, :

    %\begin{tcolorbox}
    %\begin{corl}\label{corl:all_subgroups}
    %    If $\G$ has no infinite locally finite subgroups, then there exists a free $\G$-flow that is $\Delta$-minimal for every infinite subgroup $\Delta \leq \G$.
    %\end{corl}
    %\end{tcolorbox}

    \begin{tcolorbox}
    \begin{corl}\label{corl:no_loc_fin}
        If $\G$ has at most countably many infinite locally finite subgroups, then there exists a free $\G$-flow %subshift $X \subset 2^\G$
        that is $\Delta$-minimal with respect to every infinite subgroup $\Delta \leq \G$.
    \end{corl}
    \end{tcolorbox}
    \begin{scproof}
        As in the proof of Corollary~\ref{corl:fg_subgroups}, we just need to find a free $\G$-flow %subshift $X \subset 2^\G$
        that is $\Delta$-minimal for every infinite subgroup $\Delta \leq \G$ that is either finitely generated or locally finite. %Immediate from Corollary~\ref{corl:fg_subgroups} \ep{and its proof}.
    \end{scproof}

    In particular, there exists a free $\mathbb{F}_2$-flow that is $\langle \gamma \rangle$-minimal for all non-identity elements $\gamma \in \mathbb{F}_2$ simultaneously. In general, Corollaries~\ref{corl:fg_subgroups} and \ref{corl:no_loc_fin} are new %result of Corollary~\ref{corl:no_loc_fin} is new
    for most non-Abelian groups $\G$. %of Corollary~\ref{corl:no_loc_fin} is new for most groups $\G$. %, for example for the non-Abelian free groups. %$\mathbb{F}_n$. %(the only previously known cases were the groups $\G$ all of whose infinite subgroups are normal, which are covered by Theorem~\ref{theo:FSZ}).
    %Note also that the $\G$-flows in Corollaries~\ref{corl:fg_subgroups} and \ref{corl:no_loc_fin} can be chosen to be subflows of $2^\G$.
    Notice that a nontrivial $\G$-flow cannot be $S$-minimal for all infinite subsets $S \subseteq \G$ at once; in that sense, the restriction to subgroups in Corollary~\ref{corl:no_loc_fin} is essential.

    \begin{remk}
        %We remark that
        The conclusion of Corollary~\ref{corl:no_loc_fin} may fail if $\G$ has uncountably many infinite locally finite subgroups. %Indeed, % cannot in general be extended to include \emph{all} (possibly locally finite) infinite subgroups $\Delta \leq \G$. % may fail when $\G$ does have an infinite locally finite subgroup.
    For example,
    Seward, Zucker, and the second named author showed that  it fails for %$\G \cong \bigoplus_{n \in \N} (\Z/2\Z)$---and hence also if
    $\G = \bigoplus_{n \in \N} (\Z/2\Z)$ %, every $\G$-flow on an infinite underlying space fails to be $\Delta$-minimal for some infinite subgroup $\Delta \leq \G$
    \cite[Thm.~2.5]{FSZ}. %It remains an interesting open problem to characterize precisely the class of groups $\G$ for which  %, $\G$ has an infinite subgroup for which the flow is not minimal \todo{cite}. %It remains an interesting open problem to characterize the class of groups $\G$ that admit a free $\G$-flow which is $\Delta$-minimal for \emph{all} infinite subgroups $\Delta \leq \G$. %for which the conclusion of Corollary~\ref{corl:all_subgroups} holds.
    Exactly characterizing the class of groups $\G$ that satisfy % for which
    the conclusion of Corollary~\ref{corl:no_loc_fin} %. %holds.
    %admit a free $\G$-flow which is $\Delta$-minimal with respect to every infinite subgroup $\Delta \leq \G$
    remains an interesting open problem.
    \end{remk}

    %\subsection{A generalization: $\mathcal{F}$-minimal $\G$-flows}

    %\subsection{A refinement: syndetic sets and $\mathcal{F}$-minimal flows}

    %For some applications, Theorem~\ref{theo:main} %can be further generalized
    %needs to
    %Let us now we describe some refinements of Theorem~\ref{theo:main}. First, we note that Theorem~\ref{theo:main} can be generalized using the following definitions:
    We can generalize Theorem~\ref{theo:main} further using the following definitions:

    \begin{defn}[Syndetic sets]
        Let $\G \acts X$ be an action of $\G$. A set $U \subseteq X$ is \emphd{$F$-syndetic} for a finite subset $F \subset \G$ if $F^{-1} \cdot U = X$, i.e., if for each $x \in X$, there is $\sigma \in F$ such that $\sigma \cdot x \in U$.
    \end{defn}

    \begin{defn}[$\mathcal{F}$-minimal $\G$-flows]\label{defn:F-minimal}
        Let $\mathcal{F}$ be a family of finite subsets of $\G$. We say that a $\G$-flow $X$ is \emphd{$\mathcal{F}$-minimal} if for every nonempty open set $U \subseteq X$, there is $F \in \mathcal{F}$ such that $U$ is $F$-syndetic. %$F^{-1} \cdot U = X$ \ep{that is, for each $x \in X$, there is $\sigma \in F$ such that $\sigma \cdot x \in U$}.
    \end{defn}

    Definition~\ref{defn:F-minimal} extends Definition~\ref{defn:S-minimal} thanks to the following well-known observation:

    \begin{prop}\label{prop:synd1}
        A $\G$-flow $X$ is $S$-minimal for a subset $S \subseteq \G$ if and only if it is $\finset{S}$-minimal, where $\finset{S}$ is the family of all finite subsets of $S$. %every nonempty open subset $U \subseteq X$ is $S$-syndetic.
    \end{prop}
    \begin{scproof}
        Clearly, if $X$ is $\finset{S}$-minimal, the partial orbit
        $S \cdot x$ is dense for all $x \in X$. Conversely, suppose $X$ is $S$-minimal and let $U \subseteq X$ be nonempty and open. Then we have  %Since $X$ is $S$-minimal, we have
        %$(S \cdot x) \cap U \neq \0$ for all $x \in X$, i.e.,
        $S^{-1} \cdot U = X$, and the compactness of $X$ yields a finite set $F \in \finset{S}$ such that $F^{-1}\cdot U = X$. %, as desired.
    \end{scproof}

    A family of finite sets $\mathcal{F}$ is called \emphd{unbounded} if $\sup_{F \in \mathcal{F}} |F| = \infty$. Theorem~\ref{theo:main} %and Corollary~\ref{corl:subshift} 
    remains valid if we replace the sets $S_n$ and the corresponding families $\finset{S_n}$ by {arbitrary} unbounded families $\mathcal{F}_n$: % of finite subsets of $\G$:

    \begin{tcolorbox}
    \begin{theo}\label{theo:main_prime}
        If $(\mathcal{F}_n)_{n \in \N}$ is a sequence of unbounded families of finite subsets of $\G$, then there exists a free $\G$-flow %$X$
        that is $\mathcal{F}_n$-minimal for all $n \in \N$.
    \end{theo}
    \end{tcolorbox}

    %\begin{tcolorbox}
    %\begin{corl}\label{corl:subshift_prime}
    %    If $(\mathcal{F}_n)_{n \in \N}$ is a sequence of unbounded families of finite subsets of $\G$, then there exists a free subshift $Z \subset 2^\G$
     %   that is $\mathcal{F}_n$-minimal for all $n \in \N$.
    %\end{corl}
    %\end{tcolorbox}

    %In the remainder of the paper, 
    From now on, we shall treat Theorem~\ref{theo:main_prime} as our main result (with Theorem~\ref{theo:main} as a special case), and most of the paper will be devoted to its proof. Although the concept of $\mathcal{F}$-minimality may appear somewhat artificial, it provides a natural framework for our proof techniques. Additionally, for some of the %in \S\ref{subsec:sections}, we shall present some
    applications we present here, %where %that require
    the general setting of $\mathcal{F}$-minimality is essential. % in \S\todo{ref}. %hile the concept of $\mathcal{F}$-minimality may seem somewhat artificial, it provides a natural setting for our arguments and has some compelling applications that we shall describe in the next subsection. % in descriptive set theory. %; see \S\todo{ref}. %the above results have natural applications in descriptive set theory.

    %\subsection{Subshifts with minimal subdynamics and their applications}\label{subsec:subshift_statements}

    
    %It is worth noting that the $\G$-flow in Theorem~\ref{theo:main} may be chosen to have a special form, namely that of a \emph{subshift}. 
    %
    %
    %In Theorem~\ref{theo:main}\ref{item:subshift}, $2^\G$ is the space of all maps $\G \to \set{0,1}$ equipped with the {shift action} $\G \acts 2^\G$, which is defined as follows. % a \emph{Bernoulli shift}, which is a $\G$-flow defined as follows.
    %Another corollary of Theorem~\ref{theo:main} allows us to choose 
    
    %It is possible to be more specific about the class of $\G$-flows to which %from which
    %the $\G$-flow %produced by
    %in Theorem~\ref{theo:main_prime} belongs. %is drawn. %type of a $\G$-flow produced by Theorem~\ref{theo:main_prime}.
    Our proof of Theorem~\ref{theo:main_prime} naturally yields a $\G$-flow whose underlying space is Polish---indeed, it is homeomorphic to the Cantor space. With a little extra work, we are able to give it a particularly nice special form, namely that of a \emph{subshift}. %First, we need a few definitions. 
    We use the notation %In what follows, we write
    $\N \defeq \set{0,1,2,\ldots}$ and $\N^+\defeq\set{1,2,3,\ldots}$. %to denote the sets of all nonnegative and positive integers, respectively.
    Each natural number
    $k \in \N$ is identified with the $k$-element set $\set{i \in \N \,:\, i < k}$ and given the discrete topology. For $k \in \N^+$, the \emphd{shift action} $\G \acts k^\G$ %is the action
    of $\G$ on the product space $k^\G$ is defined by
    \[
        (\gamma \cdot x)(\delta) \,\defeq\, x(\delta \gamma) \quad \text{for all $x \colon \G \to k$ and $\gamma$, $\delta \in \G$}.
    \]
    This makes $k^\G$ a %zero-dimensional Polish
    $\G$-flow, called a \emphd{Bernoulli shift}, or simply a \emphd{shift}. Subflows of $k^\G$ are called \emphd{subshifts} and are important examples of dynamical systems %that are
    studied in {symbolic dynamics} \cite{subshifts1,subshifts2}. 

    \begin{tcolorbox}
    \begin{corl}\label{corl:subshift}\label{corl:subshift_prime}
        If $(\mathcal{F}_n)_{n \in \N}$ is a sequence of unbounded families of finite subsets of $\G$, then there exists a free subflow of $2^\G$
        that is $\mathcal{F}_n$-minimal for all $n \in \N$.%If $(S_n)_{n \in \N}$ is a countable family of infinite subsets of $\G$, then there exists a free subshift $Z \subset 2^\G$
        %that is $S_n$-minimal for all $n \in \N$.
        %
        %Let $(S_n)_{n \in \N}$ be a countable family of infinite subsets of $\G$.
        %
        %\begin{enumerate}[label=\ep{\normalfont\roman*}]
        %    \item There exists a free $\G$-flow $X$ that is $S_n$-minimal for all $n \in \N$.
        %
        %    \item\label{item:subshift} Furthermore, such $X$ can be chosen to be a subflow of $2^\G$.
        %\end{enumerate}
        %
        %, then there exists a free $\G$-flow $X$ that is $S_n$-minimal for all $n \in \N$. Furthermore, 
    \end{corl}
    \end{tcolorbox}

    In particular, for any countable family $(S_n)_{n \in \N}$ of infinite subsets of $\G$, there exists a free subflow of $2^\G$ %subshift $X \subset 2^\G$
    that is $S_n$-minimal for all $n \in \N$.    
    To derive Corollary~\ref{corl:subshift} from Theorem~\ref{theo:main_prime}, we use a result of Seward and Tucker-Drob \cite{ST-D}, which yields a free subshift $Y \subset 2^\G$ with a certain ``universality'' property, and then find the desired $\G$-flow among the subflows of $Y$. %must have a subflow $X \subseteq Y$ that is $S_n$-minimal for all $n \in \N$.
    %The details of this argument are given in \S\ref{sec:subshift}.
    See \S\ref{sec:subshift} for details.
    
    We should note that the existence of a free subshift for an arbitrary countable group $\G$ is already a highly nontrivial fact. It was established in full generality by Gao, Jackson, and Seward \cite{GJS1,GJS2} after partial results in various special cases due to Dranishnikov and Schroeder \cite{DS} and Glasner and Uspenskij \cite{GU} (the $\G = \Z$ case appears implicitly as far back as the work of Thue \cite{Thue}, Morse \cite{Morse}, and Morse--Hedlund \cite{MorseHedlund}). %See also %the subsequent paper
    %\cite{ABT} by
    Aubrun, Barbieri, and Thomass\'{e} \cite{ABT} subsequently found a simpler proof for general $\G$ using the probabilistic method. Since then, several works have appeared that construct free subshifts with various additional properties \cite{ST-D,LLLDyn1,LLLDyn4}, and Corollary~\ref{corl:subshift} continues this trend.


    \subsection{Applications to constructing disjoint flows}

    It turns out that our results have close ties to the theory of {disjointness} for $\G$-flows, introduced by Furstenberg in his seminal paper \cite{Furstenberg}. 

    \begin{defn}[Joinings and disjointness]
        Let $X$ and $Y$ be $\G$-flows. We view $X \times Y$ as a $\G$-flow with the diagonal action of $\G$. A \emphd{joining} of $X$ and $Y$ is a subflow $Z \subseteq X \times Y$ that projects onto $X$ and $Y$. The $\G$-flows $X$, $Y$ are \emphd{disjoint}, in symbols $X \perp Y$, if their only joining is $X \times Y$ itself.
    \end{defn}

    An easy observation is that if $X \perp Y$, then at least one of $X$, $Y$ is minimal. Also, if both $X$ and $Y$ are minimal, then $X \perp Y$ if and only if $X \times Y$ is also minimal.

    Every $\G$-flow is disjoint from the trivial action of $\G$ on a single-point space, but this is hardly an interesting instance of the disjointness relation. Seeking more illuminating examples, we ask: % Here we are interested in the following basic question:

    \begin{tcolorbox}[width=\linewidth, sharp corners=all, colback=white!95!black, boxrule=0pt,colframe=white] 
        For which $\G$-flows $X$ does there exist some free $\G$-flow $Y$ such that $X \perp Y$?
    \end{tcolorbox}

    We shall specifically focus on {Polish} (equivalently, metrizable) $\G$-flows $X$. Some important results related to this question were obtained by Glasner, Tsankov, Weiss, and Zucker in their breakthrough paper \cite{BernoulliDisjointness}. Generalizing a theorem of Furstenberg \cite{Furstenberg}, they showed that $2^\G$ and, more generally, all so-called \emph{strongly irreducible subshifts} are disjoint from every minimal $\G$-flow \cite[Thm.~6.2]{BernoulliDisjointness}. They also showed that if a Polish $\G$-flow $X$ is minimal, then it is disjoint from some free minimal Polish $\G$-flow \cite[Thm.~1.2(i)]{BernoulliDisjointness}. The assumption that $X$ is Polish cannot be removed here, as the \emph{universal minimal flow} $\mathcal{M}(\G)$ is a non-metrizable counterexample. 
    
    Using our results on minimal subdynamics, we give a complete answer to the above question (for Polish $\G$-flows). We need the following standard definition:

    \begin{defn}[Wandering sets and points]
        A subset $U \subseteq X$ of a $\G$-flow $X$ is \emphd{wandering} if $U \cap (\gamma \cdot U) = \0$ for all but finitely many group elements $\gamma \in \G$, and a point $x \in X$ is \emphd{wandering} if it has a wandering neighborhood.
    \end{defn}

    %Let us make a few remarks about this definition.
    By definition, the set of all wandering points in a $\G$-flow $X$ is open and $\G$-invariant. It is clear that if $x \in X$ is a wandering point, then the stabilizer of $x$ is finite. It is also not hard to see that if the stabilizer of a wandering point $x$ is trivial (for instance, if $\G$ is torsion-free), then $x$ has a neighborhood $U$ such that %$U \cap (\gamma \cdot U) = \0$ for all non-identity elements $\gamma \in \G$, i.e.,
    the sets $(\gamma \cdot U)_{\gamma \in \G}$ are pairwise disjoint. %Minimal $\G$-flows have no wandering points (see, e.g., \cite[Lem.~2.1]{LLLDyn3}).

    We show that the existence of wandering points is the only obstruction to being disjoint from some free $\G$-flow:

    \begin{tcolorbox}
    \begin{theo}\label{theo:disjoint_one}
        The following statements are equivalent for a Polish $\G$-flow $X$:
        \begin{enumerate}[label=\ep{\normalfont\arabic*}]
            \item $X$ is disjoint from some infinite $\G$-flow,
            \item $X$ is disjoint from some free minimal subflow of $2^\G$,
            \item $X$ has no wandering points.
        \end{enumerate}
    \end{theo}
    \end{tcolorbox}

    Furthermore, we can extend this result to countable families of Polish $\G$-flows:

    \begin{tcolorbox}
    \begin{theo}\label{theo:disjoint_many}
        If $(X_n)_{n \in \N}$ is a countable family of Polish $\G$-flows with no wandering points, then there exists a free minimal subflow of $2^\G$ that is disjoint from $X_n$ for all $n \in \N$.
    \end{theo}
    \end{tcolorbox}

    %Since a minimal $\G$-flow has no wandering points (see, e.g., \cite[Lem.~2.1]{LLLDyn3}), Theorem~\ref{theo:disjoint_one} in particular implies the result of Glasner \emph{et al.} mentioned above that every minimal Polish $\G$-flow is disjoint from some other minimal Polish $\G$-flow \cite[Thm.~1.2(i)]{BernoulliDisjointness}.
    These results are established in \S\ref{sec:disjoint} by reducing the disjointness relation to an instance of the minimal subdynamics problem (see Proposition~\ref{prop:characterize}). %As a step in the proof, we show that if $X$ is a Polish $\G$-flow with no wandering points, then there is a family $(S_n)_{n \in \N}$ of infinite subsets of $\G$ such that $X \perp Y$ whenever $Y$ is $S_n$-minimal for all $n \in \N$.
    It is possible to deduce Theorem~\ref{theo:disjoint_many} from Theorem~\ref{theo:disjoint_one} by letting $X$ be the one-point compactification of %the disjoint union
    $\bigsqcup_{n \in \N} X_n$, but we give a direct proof instead. Minimal $\G$-flows have no wandering points (see, e.g., \cite[Lem.~2.1]{LLLDyn3}), so Theorem~\ref{theo:disjoint_one} includes the aforementioned result \cite[Thm.~1.2(i)]{BernoulliDisjointness} of Glasner \emph{et al.} as a special case.  %We prove Theorems~\ref{theo:disjoint_one} and \ref{theo:disjoint_many} in \S\ref{sec:disjoint}. %These results are proved in \S\ref{sec:disjoint}. In brief,
    %As a key step in the proof of Theorem~\ref{theo:disjoint_many}, we show that if $X$ is a Polish $\G$-flow with no wandering points, then there is a family $(S_n)_{n \in \N}$ of infinite subsets of $\G$ such that $X \perp Y$ whenever $Y$ is $S_n$-minimal for all $n \in \N$.
    
    
    One of the reasons these facts have not been discovered earlier is their % To highlight the 
    tight connection to the minimal subdynamics problem. %, between disjointness and minimal subdynamics,
    Indeed, consider the following example. %let us mention that Theorem~\ref{theo:disjoint_many} can be used to easily derive the subgroup case of Theorem~\ref{theo:main} (and hence also its Corollaries~\ref{corl:fg_subgroups} and \ref{corl:no_loc_fin}). Indeed,
    Let $\Delta \leq \G$ be a subgroup. Then $\G$ naturally acts on the set $\G/\Delta$ of the left cosets of $\Delta$. We endow $\G/\Delta$ with the discrete topology and let $(\G/\Delta)^*$ be the one-point extension of $\G/\Delta$, where the action of $\G$ is extended to the point at infinity by making it fixed. (Note that if the index of $\Delta$ in $\G$ is finite, then the point at infinity is isolated.) %compactification of $\G/\Delta$ obtained by adjoining at most one fixed point at infinity. 
    %
    %then compactify $\G/\Delta$ as follows. If $|\G \colon \Delta| < \infty$, then $\G/\Delta$ is already compact, so we let $(\G/\Delta)^* \defeq \G/\Delta$. On the other hand, if $|\G \colon \Delta| = \infty$, then we let $(\G/\Delta)^*$ be the one-point compactification of $\G/\Delta$, where the action of $\G$ on $\G/\Delta$ is extended to $(\G/\Delta)^*$ by making the point at infinity fixed.
    The space $(\G/\Delta)^*$ is a (countable) Polish $\G$-flow, and we have the following:

    \begin{prop}\label{prop:cosets}
        %Let $\Delta \leq \G$ be a subgroup. A minimal $\G$-flow $X$ is disjoint from $(\G/\Delta)^*$ if and only if $X$ is $\Delta$-minimal
        %
        The following statements are equivalent for a $\G$-flow $X$ and a subgroup $\Delta \leq \G$:
        \begin{enumerate}[label=\ep{\normalfont\arabic*}]
            \item $X \perp (\G/\Delta)^*$,
            \item $X$ is $\Delta$-minimal.
        \end{enumerate}
    \end{prop}

    See \S\ref{sec:disjoint} for the proof. If $\Delta$ is infinite, then there are no wandering points in $(\G/\Delta)^*$ (because the stabilizer of every point is infinite). Therefore, Theorem~\ref{theo:disjoint_many} implies the subgroup case of Theorem~\ref{theo:main}, i.e., it shows that for any countable family $(\Delta_n)_{n \in \N}$ of infinite subgroups of $\G$, there exists a free $\G$-flow $X$ that is $\Delta_n$-minimal for all $n \in \N$. In particular, Theorem~\ref{theo:disjoint_many} is already sufficient to derive Corollaries~\ref{corl:fg_subgroups} and \ref{corl:no_loc_fin}.

    

    \subsection{Applications to Borel complete sections}\label{subsec:sections}

    As mentioned previously, the existence of free subshifts for arbitrary $\G$ was established by Gao, Jackson, and Seward \cite{GJS1,GJS2}. Their work 
%
 %    Gao, Jackson, and Seward's work in \cite{GJS1,GJS2}
 was motivated by the realization that free subshifts offer deep insight into the topological and descriptive set-theoretic properties of %interplay between
 the shift action $\G \acts 2^\G$ %and the topology and the
 %Borel structure on $2^\G$
 by putting significant constraints on the geometry of complete sections in $\Free(2^\G)$ (see below for the definition). This idea was applied with great success in subsequent contributions by
    %
    %As convincingly demonstrated by
    Gao, Jackson, Krohne, and Seward \cite{GJKS_Borel,GJKS_continuous}. %, with a focus on $\Z^n$-actions. %, free subshifts % the existence of free subshifts is a powerful tool for understanding the interplay between
    %are extremely useful for relating the shift action $\G \acts 2^\G$ %in relation to %
    %to %and both %its interactions with
    %the topology \cite{GJKS_continuous} and the Borel structure \cite{GJKS_Borel} %$\sigma$-algebra
    %on $2^\G$.  
    %
    %
    %from the point of view of descriptive set theory
    %\cite{GJKS_continuous,GJKS_Borel}.
    %Some of the results in these papers may be improved by using free subshifts with minimal subdynamics. We describe one particularly striking example here.
    It turns out that some of the %To illustrate the power of our results, we now describe how some of
    main results of % describe some theorems from %proved in 
    \cite{GJKS_Borel}
    %with the help of a free subshift
    %that
    can be significantly improved with the help of free subshifts with minimal subdynamics, %by using
    %Corollary~\ref{corl:subshift_prime}. %free subshifts with minimal subdynamics.
    %As we explain below, these applications
    specifically using the notion of $\mathcal{F}$-minimality (Definition~\ref{defn:F-minimal}).
    
    %For example, the following result is proved in \cite{GJKS_Borel} with the help of a free subshift. 
    Let $\Free(2^\G)$ denote the \emphd{free part} of $2^\G$, i.e., the set of all points $x \in 2^\G$ %such that $\gamma \cdot x \neq x$ for all $\gamma \in \G \setminus \set{\iden}$. %
    with trivial stabilizers.
    Then $\Free(2^\G)$ is a shift-invariant dense $G_\delta$ subset of $2^\G$ % that is invariant under the shift action $\G \acts 2^\G$
    (see, e.g.,
    \cite[Lem.~2.3]{BernoulliDisjointness} for a proof of density), and it is the largest subspace of $2^\G$ on which $\G$ acts freely. % the shift action is free. % the action of $\G$ on $\Free(2^\G)$ is, by definition, free.
    A \emphd{complete section} in $\Free(2^\G)$ is a subset $B \subseteq \Free(2^\G)$ that meets every $\G$-orbit, i.e.,
    such that
    $\G \cdot B = \Free(2^\G)$. %that meets every $\G$-orbit.
    The following is a typical fact that can be proved using Gao \emph{et al.}'s methods:

    \begin{theo}[{Gao--Jackson--Krohne--Seward \cite[Thm.~1.1]{GJKS_Borel}}]\label{theo:GJKS1}
        Suppose $(B_n)_{n \in \N}$ is a sequence of Borel complete sections in $\Free(2^\G)$ and $(F_n)_{n \in \N}$ is a sequence of finite subsets of $\G$ such that every finite subset $F \subset \G$ is contained in $F_n$ for some $n \in \N$. %for all finite $F \subset \G$, there is $n \in \N$ with $F \subseteq F_n$.
        Then there exists a point $x \in \Free(2^\G)$ satisfying $x \in F_n \cdot B_n$ %$(F_n \cdot x) \cap B_n \neq \0$
        for infinitely many $n \in \N$.
    \end{theo}

    Working in a free $\set{F_n}_{n \in \N}$-minimal subshift, we are able to generalize Theorem~\ref{theo:GJKS1} by removing all constraints on the shape of the sets $F_n$---we only need $|F_n|$ to be unbounded, so that Corollary~\ref{corl:subshift} may be applied to the family $\set{F_n}_{n \in \N}$:

    \begin{tcolorbox}
    \begin{theo}\label{theo:DST1}
        Suppose $(B_n)_{n \in \N}$ is a sequence of Borel complete sections in $\Free(2^\G)$ and $(F_n)_{n \in \N}$ is a sequence of finite subsets of $\G$ such that $\sup_{n \in \N}|F_n| = \infty$. Then there exists a point $x \in \Free(2^\G)$ satisfying $x \in F_n \cdot B_n$ for infinitely many $n \in \N$.
    \end{theo}
    \end{tcolorbox}

    
    Here is another example. Following the terminology of \cite{LLLDyn3}, we say that a set $A \subseteq \Free(2^\G)$ \emphd{traps} a point $x \in \Free(2^\G)$ (or that $x$ is \emphd{trapped} in $A$) if $\G \cdot x \subseteq A$. Gao \emph{et al.}~showed that for every Borel complete section $B \subseteq \Free(2^\G)$, some point is trapped in the union of finitely many shifts of $B$:

    \begin{theo}[{Gao--Jackson--Krohne--Seward \cite[Thm.~1.2]{GJKS_Borel}}]\label{theo:GJKS}
        If $B \subseteq \Free(2^\G)$ is a Borel complete section, then there exists a finite set $F \subset \G$ such that the set $F \cdot B$ traps a point.
    \end{theo}

    An interesting question is what additional properties the set $F$ in Theorem~\ref{theo:GJKS} may have. %can be chosen to have some additional desirable properties.
    For example, Gao \emph{et al.}~showed that for the group $\G = \Z^n$ with $n \in \N^+$, it is possible to pick $F$ so that $\|\gamma\|_1$ is odd for all $\gamma \in F$ \cite[Thm.~1.3]{GJKS_Borel}. This result is obtained with an ad hoc construction specifically tailored to the group $\Z^n$ and the ``odd $1$-norm'' property, %. As a consequence, it is unclear how far such results can be generalized. %
    and it is unclear how far it can be generalized. 
    %Thanks to
    Corollary~\ref{corl:subshift_prime} %we are able to %
    allows us to eliminate the need for such ad hoc arguments and
    completely settle the problem. %in full generality.
    Namely, we show that \emph{any} large enough finite set $F \subset \G$ works: % as long as it is large enough:

    \begin{tcolorbox}
    \begin{theo}\label{theo:DST}
        If $B \subseteq \Free(2^\G)$ is a Borel complete section, then there exists $n \in \N$ such that for every finite set $F \subset \G$ of size at least $n$, the set $F \cdot B$ traps a point.
    \end{theo}
    \end{tcolorbox}

    We argue by contradiction: assuming Theorem~\ref{theo:DST} fails for a Borel complete section $B \subseteq \Free(2^\G)$, we may pick sets $F_n \subset \G$ with $|F_n| = n$ so that $F_n \cdot B$ does not trap a point. We then finish the proof by working in a free $\set{F_n^{-1}}_{n \in \N}$-minimal subshift. The details are given in \S\ref{subsec:subshift_appl}.

    A version of Theorem~\ref{theo:DST} is proved in \cite[Thm.~1.2]{LLLDyn3} for \emph{open} sets $B \subseteq \Free(2^\G)$ (the proof in the open case is much simpler). We find it quite surprising that such a strong property is in fact shared by all \emph{Borel} complete sections. (Actually, our proofs of Theorems~\ref{theo:DST1} and \ref{theo:DST} do not even need the complete sections to be Borel---it is enough to assume that their intersections with every free subshift $X \subset 2^\G$ are Baire-measurable in $X$.) 

    

    %\medskip

    
    

    %\subsection{Proof overview}

    %\subsubsection{The sources of our approach}

    \subsection{What goes into the proof of Theorem~\ref{theo:main_prime}}\label{subsec:tools}

    %Let us now say a few words about the proof of Theorem~\ref{theo:main}. 
    %We now present an outline of the proof of Theorem~\ref{theo:main}. Our arguments draw
    Let us now say a few words about the proof of Theorem~\ref{theo:main_prime}. Although it is a result that concerns some of the most basic concepts in topological dynamics, its proof relies on a range of distinctly modern techniques. Moreover, these techniques come not only from topological dynamics, but also from descriptive set theory and combinatorics. %More specifically, it % Our proof of Theorem~\ref{theo:main}
    %draws upon three recent streams of ideas, which we now briefly describe. %before moving on to a more detailed discussion of the proof.
    Here we describe the main streams of ideas that our argument draws upon and highlight some of the key novelties of our approach. 
    This high-level discussion will be followed in \S\ref{sec:outline} by a detailed outline of the proof of Theorem~\ref{theo:main_prime} that explains how these ideas fit together.

    %\medskip


    \subsubsection*{Existence via genericity} Instead of constructing an explicit example of a $\G$-flow fulfilling the requirements of Theorem~\ref{theo:main_prime}, we will show that the theorem is witnessed by a ``typical'' $\G$-flow from a certain class. More precisely, we will prove that $\G$-flows with the desired properties form a dense $G_\delta$ subset in that class, and hence they are \emph{generic} in the Baire category sense.
    
    The study of generic dynamical systems is an old subject, dating back at least to the seminar work of Oxtoby and Ulam \cite{OxtobyUlam}, Halmos \cite{Halmos1,Halmos2}, and Rokhlin \cite{Rokhlin}. See \cite{Gen1,Gen2,Gen3,Gen4,Gen5,Gen6,Gen7,Gen8,Gen9} for a selection of works on genericity properties in topological dynamics from the past 25 years. % are \emph{generic} in that class, where the word ``generic'' is used in the Baire category sense. That is, we call a property of elements of some Polish space \emphd{generic} if it is satisfied by a comeager set of points.
        
    %The immediate context for our work is given by the papers \cite{FTVF,FSZ}.
    
    We borrow the idea of using genericity to obtain minimal $\G$-flows with favorable properties from the papers \cite{FTVF,FSZ}. In \cite{FTVF}, Tamuz, Vahidi Ferdowsi, and the second named author gave a genericity argument showing that every group $\G$ with a nontrivial ICC quotient admits a nontrivial proximal minimal $\G$-flow. The proof of Theorem~\ref{theo:FSZ} given in \cite{FSZ} similarly relies on genericity. The approach employed in \cite{FTVF,FSZ}, which originates in the paper \cite{FT} by Tamuz and the second named author, focuses on a specific class of $\G$-flows, namely the (closure of the) space of {strongly irreducible subshifts}. By contrast, we consider a different class that is defined specifically for the purpose of proving Theorem~\ref{theo:main_prime} (though we expect it to have other uses as well). To identify this class, we rely on the existence of a certain \emph{non-compact} action of $\G$; this is one of the central innovations of our paper. This part of the argument is described in %more detail in
    \S\ref{subsubsec:ample}.
    
    %These papers employ an approach originating in the paper \cite{FT} by Tamuz and the second named author  genericity arguments to show the existence of minimal $\G$-flows with additional properties. Namely, in \cite{FTVF}, Tamuz, Vahidi Ferdowsi, and the second named author  This style of argument was introduced by Tamuz and the second named author in \cite{FT} and made explicit by Tamuz, Vahidi Ferdowsi, and the second named author in \cite{FTVF}. It was also employed in the proof of Theorem~\ref{theo:FSZ} given in \cite{FSZ}. The papers \cite{FTVF,FSZ} work in the same class of $\G$-flows, namely the (closure of the) space of so-called \emph{strongly irreducible subshifts}. By contrast, we use a different class that is defined specifically for the purpose of proving Theorem~\ref{theo:main}. %In order to argue that the class of subshifts we consider is nonempty, we rely on a result of Seward and Tucker-Drob \todo{cite} (or, to be precise, on the simplified proof of their result due to the first named author \todo{cite}), which yields a free subshift $Y \subset 2^\G$ with a certain ``Borel universality'' property.
    %    

    \subsubsection*{Descriptive combinatorics and large-scale geometry} Our methods are strongly informed by the area called \emph{descriptive combinatorics}. This is a subject that fuses combinatorics and descriptive set theory by studying definability properties of combinatorial constructions on Polish spaces and viewing descriptive set-theoretic problems through a combinatorial lens. The systematic study of interactions between descriptive set theory and combinatorics was launched by Kechris, Solecki, and Todorcevic in their seminal paper \cite{KST} and has since developed into a rich subject with many connections to other fields, including dynamical systems. For an overview of this area, see the survey \cite{KechrisMarks} by Kechris and Marks, the introductory article \cite{Pikh_survey} by Pikhurko, and %as well as
    the note \cite{Notices} by the first named author. %covering some more recent developments.

    One of the most exciting trends in descriptive combinatorics of the past few years has been the use of concepts inspired by \emph{large-scale geometry}, pioneered by Conley, Jackson, Marks, Seward, and Tucker-Drob in their landmark paper \cite{CJMSTD}. A fundamental notion in large-scale geometry is the \emph{asymptotic dimension} of a metric space, introduced by Gromov \cite[\S1.E]{gromov1993asymptotic}. In \cite{CJMSTD}, Conley \emph{et al.}~developed a {Borel} version of asymptotic dimension and presented an array of impressive applications. They also defined a closely related parameter, the \emph{asymptotic separation index}, which turned out to be extremely useful in descriptive combinatorics \cite{ASIalgorithms,BW,BWKonig,CJMSTD,WeilacherFinDim,Grids2}. As we explain in \S\S\ref{subsubsec:example} and \ref{subsubsec:asi}, a key ingredient in our proof of Theorem~\ref{theo:main_prime} is a {continuous} variant of asymptotic separation index.
    
    Although asymptotic separation index has been associated with group actions since its initial appearance in \cite{CJMSTD}, the present work is the first to apply this notion to solve a problem of a purely dynamical nature, and we hope it will become a part of the topological dynamics toolkit.

    \subsubsection*{Definable versions of the Lov\'asz Local Lemma} The \emph{Lov\'asz Local Lemma}, or the \emph{LLL} for short, is a powerful probabilistic tool due to Erd\H{o}s and Lov\'asz \cite{EL} which has a plethora of applications throughout combinatorics; see \cite{AS,MolloyReed} for many examples. Furthermore, the LLL has recently been employed to address a number of problems in ergodic theory and topological dynamics; see, e.g., \cite{ABT,Elek,Ber_cont,LLLDyn1,LLLDyn2,LLLDyn3,LLLDyn4}. %This work presents %yet
    %another addition to this list.
    
    The LLL is an existence result, and it is particularly well-suited for showing that a given structure $X$ admits a mapping $f \colon X \to k$ for some $k \in \N^+$ satisfying a specified set of constraints. Roughly speaking, in order for the LLL to apply in this context, two requirements must be met: First, a random mapping $f \colon X \to k$ should be ``likely'' to fulfill each individual constraint; second, the constraints must not interact with each other ``too much.'' (See \S\ref{subsec:LLL} for the precise statement.)

    %A matter of major interest is to determine under which assumptions the LLL or an LLL-like statement can be used to derive ``constructive'' results. For instance,

    A major research direction concerns 
    %There has been significant recent work on
    versions of the LLL that are ``constructive'' in various senses: algorithmic \cite{Beck,BGRDeterministicLLL,BMUDeterministicLLL,FG,MoserTardos}, computable \cite{RSh}, Borel/measurable \cite{CGMPT,BerDist,BerMeas,BerYuBorelLLL}, and continuous \cite{BerDist,Ber_cont}. % versions of the LLL.
    Most relevantly for the present work, Weilacher and the first named author proved a {Borel} version of the LLL under a finite asymptotic separation index assumption \cite{BW}. We employ a slight variation of their argument to establish a version of the LLL that yields \emph{continuous} maps $f \colon X \to k$ satisfying the constraints. %Under suitable assumptions, this result can be employed to obtain a \emph{continuous} map $f \colon X \to k$ satisfying the given constraints.
    We discuss %our continuous version of the LLL and
    the way our continuous version of the LLL becomes used in the proof of Theorem~\ref{theo:main_prime} in \S\ref{subsubsec:combi+LLL}.

    %\medskip

    %\noindent %We now turn to a detailed outline of the proof of Theorem~\ref{theo:main}, showing how the above ideas fit together.
    %After this high-level overview, we are ready to show how the above ideas fit together and give a detailed outline of the proof of Theorem~\ref{theo:main}. % in more concrete detail and show how the above ideas fit together.
    %In the next section, \S\ref{sec:outline}, we present a detailed outline of the proof of Theorem~\ref{theo:main} that explains how the above ideas fit together and reduces the task of proving Theorem~\ref{theo:main} to verifying several lemmas, which are proved in the subsequent sections. 

    \subsubsection*{Main innovations} %Let us now briefly highlight the main novel
    %features of our method.
    %
    As mentioned above, the LLL has already been applied in topological dynamics in the past. However, in most previous applications, such as the ones in \cite{Elek,ABT,LLLDyn1,LLLDyn3}, it is the ``classical,'' pure existence version of the LLL that is used. By contrast, our proof of Theorem~\ref{theo:main_prime} crucially relies on a continuous variant of the LLL, thus making significant use of recent developments in descriptive combinatorics.
    
        While our ultimate goal is to find a certain continuous action of $\G$ on a \emph{compact} space, with the compactness requirement being the main challenge, our argument proceeds by first constructing an action $\G \acts X$ on a \emph{non-compact} Polish space $X$ with a technical property of being \emph{\sr}. % (Definition~\ref{defn:amply_syndetic}).
        The desired $\G$-flow is then obtained by considering spaces of the form $\overline{\rho(X)}$, where $\rho \colon X \to Y$ is a continuous $\G$-equivariant map from $X$ to a $\G$-flow $Y$. %and the line indicates closure.
        This somewhat counterintuitive approach separates the topological aspect of the problem---namely, compactness---from its combinatorial core, thus reducing the task to finding a Polish $\G$-space $X$ with sufficiently ``rich'' combinatorics.

        To find such a space $X$, we use asymptotic separation index. Specifically, we consider the space $\Sep(s)$ of all ``witnesses'' to the asymptotic separation index being at most $s$ and argue---using the LLL---that its free part is amply syndetic. Our work suggests that %Although
        the space $\Sep(s)$, which is implicit in the study of asymptotic separation index, %our work suggests that its topological and dynamical properties are worthy of being
        is worthy of being explored in its own right, and that its topological and dynamical properties % is interesting in its own right and that its topological and dynamical properties are worthy of further investigation. %
         may find further use elsewhere.

         %The necessary combinatorial properties of the space $\Sep(s)$ are proved with the help of the LLL. As mentioned before, the LLL has been applied in topological dynamics in the past. However, in most previous applications, such as the ones in \cite{Elek,ABT,LLLDyn1,LLLDyn3}, it is the ``classical,'' pure existence version of the LLL that is used. By contrast, our proof of Theorem~\ref{theo:main} crucially relies on a continuous variant of the LLL, thus making significant use of recent developments in descriptive combinatorics.

    \section{Outline of the proof of Theorem~\ref{theo:main_prime}}\label{sec:outline}

    \subsection{\Sr $\G$-spaces}\label{subsubsec:ample}

    %\todo{AB: thinking about the best presentation} 
    
    
    
    A \emphd{$\G$-space} is a topological space $X$ equipped with a continuous action $\G \acts X$. As mentioned in \S\ref{subsec:tools}, the heart of our argument is a construction of a (non-compact) Polish $\G$-space $X$ % continuous action $\G \acts X$ on a (non-compact) Polish space $X$
    satisfying the following slightly technical condition: %property:

    \begin{defn}[\Sr $\G$-spaces]\label{defn:amply_syndetic}
        An \emphd{\sr $\G$-space} is a %nonempty
        Polish $\G$-space $X$ %equipped with a %free
        %continuous action $\G \acts X$ of $\G$
        with the following property: %such that,
        for every finite tuple $U_1$, \ldots, $U_k$ of nonempty open subsets of $X$, there is an integer $n \in \N$ such that for every finite set $F \subset \G$ of size at least $n$, there exists a continuous $\G$-equivariant map $\pi \colon X \to X$ for which the sets $\pi^{-1}(U_1)$, \ldots, $\pi^{-1}(U_k)$ are $F$-syndetic. %, that is, we have $\pi(X) \subseteq \bigcap_{i = 1}^k (F^{-1} \cdot U_i)$. %satisfying
            %\[
            %    F^{-1} \cdot \pi^{-1}\left(U_1\right) \,=\, \cdots \,=\, F^{-1} \cdot   \pi^{-1}\left(U_k\right) \,=\, X.
            %\]
        %are $F$-syndetic.
    \end{defn}

    \begin{tcolorbox}
    \begin{theo}\label{theo:rich}
        There exists a nonempty, free, zero-dimensional %\footnote{Recall that a topological space is \emph{zero-dimensional} if it has a basis consisting of clopen sets.}  
        \sr $\G$-space.
    \end{theo}
    \end{tcolorbox}

    %Theorem~\ref{theo:rich} is proved with a simple explicit example, presented in Theorem~\ref{theo:example}. %of a zero-dimensional \sr $\G$-space. % in \S\ref{sec:construction}, specifically in Theorem~\ref{theo:example}.
    %However, before describing it, let us explain how Theorem~\ref{theo:rich} is used to derive Theorem~\ref{theo:main}.  %  For now, let us assume Theorem~\ref{theo:rich} holds and explain how %, but first let us explain how Theorem~\ref{theo:main} is derived from Theorem~\ref{theo:rich}. %It is inspired by ideas from descriptive set theory, namely the concept of \emph{asymptotic separation index} introduced in \cite{CJMSTD} by Conley, Jackson, Marks, Seward, and Tucker-Drob. To prove that the $\G$-space we define is indeed \sr, we rely on a tool from probabilistic combinatorics, namely the \emph{Lov\'asz Local Lemma} \cite{EL}, whose relationship to asymptotic separation index was explored by Bernshteyn and Weilacher in \cite{BW}. See \S\ref{sec:construction} for a further discussion.

    %To clarify, here we call $X$ free if the action $\G \acts X$ is free.
    Recall that a topological space is \emphd{zero-dimensional} if it has a basis consisting of clopen sets. For example, the space $k^\G$ for $k \in \N^+$ and all its subspaces are zero-dimensional.
    
    Before describing the proof of Theorem~\ref{theo:rich} (which forms the bulk of the paper), let us explain how it is used to derive Theorem~\ref{theo:main_prime}. 
    %
    %We now outline the derivation of Theorem~\ref{theo:main} from Theorem~\ref{theo:rich}. %In fact, we will prove a significantly stronger statement.
    %We will show that Theorem~\ref{theo:main} is witnessed by a ``typical'' $\G$-flow from a certain class. More precisely, we will prove that $\G$-flows with the desired properties form a dense $G_\delta$ subset in that class, and hence they are \emph{generic} in the Baire category sense. % are \emph{generic} in that class, where the word ``generic'' is used in the Baire category sense. That is, we call a property of elements of some Polish space \emphd{generic} if it is satisfied by a comeager set of points.
    %This style of argument was introduced by Frisch and Tamuz in \cite{FT} and made explicit by Frisch, Tamuz, and Vahidi Ferdowsi \cite{FTVF}. It was also employed by Frisch, Seward, and Zucker \cite{FSZ} in their proof of Theorem~\ref{theo:FSZ}. The papers \cite{FTVF,FSZ} work in the space of so-called \emph{strongly irreducible subshifts}, while our proof of Theorem~\ref{theo:main} relies on a different class of $\G$-flows defined with the help of an \sr $\G$-space.
    %
    %
    %Let $X$ be an \sr $\G$-space. %, which exists by Theorem~\ref{theo:rich}. %It is not hard to see that there exists a free Polish $\G$-flow $Y$ that admits a continuous $\G$-equivariant map $\rho\colon X \to Y$ \ep{Proposition~\todo{ref}}.
    Take an arbitrary Polish $\G$-flow $Y$. (To be clear, since $Y$ is a $\G$-flow, it is a \emph{compact} Polish space.) Let $\Sub(Y)$ be the set of all subflows of $Y$. %Note that, since $Y$ is free, every $\G$-flow in $\Sub(Y)$ is free as well.
    We equip $\Sub(Y)$ with the \emphd{Vietoris topology}. This topology has a basis consisting of all sets of the form \label{page:Vietoris}
    \[
        \llbracket U_0; U_1, \ldots, U_k\rrbracket \,\defeq\, \set{Z \in \Sub(Y) \,:\, Z \subseteq U_0, \, Z \cap U_1 \neq \0, \, \ldots, \, Z \cap U_k \neq \0},
    \]
    where $k \in \N$ and $U_0$, \ldots, $U_k$ are open subsets of $Y$. It is a standard fact that the Vietoris topology makes $\Sub(Y)$ a compact Polish space \cite[\S4.F]{KechrisDST}. %(see \S\todo{ref} for the definition), which makes $\Sub(Y)$ a compact Polish space in its own right \todo{cite}.
    Next we fix an arbitrary \sr $\G$-space $X$ and consider the following closed subset of $\Sub(Y)$:
    \[
        \Sub_X(Y) \,\defeq\, \overline{\big\{\overline{\rho(X)} \,:\, \text{$\rho \colon X \to Y$ is a continuous $\G$-equivariant map}\big\}}.
    \]
    Here horizontal lines indicate topological closure (either in $Y$ or in $\Sub(Y)$). %Note that $\Sub_X(Y)$ is a nonempty compact Polish space.
    %We claim that a generic member of $\Sub_X(Y)$ is a witness to Theorem~\ref{theo:main}:
    We claim that for any unbounded family $\mathcal{F}$ of finite subsets of $\G$, % infinite set $S \subseteq \G$,
    a generic member of $\Sub_X(Y)$ is $\mathcal{F}$-minimal:

    \begin{tcolorbox}
    \begin{theo}\label{theo:generic}
        Let $X$ be an \sr $\G$-space and let $Y$ be a Polish $\G$-flow. For each %infinite set $S \subseteq \G$
        unbounded family $\mathcal{F}$ of finite subsets of $\G$, the set \[\set{Z \in \Sub_X(Y) \,:\, \text{$Z$ is $\mathcal{F}$-minimal}}\] is dense and $G_\delta$ in $\Sub_X(Y)$.
    \end{theo}
    \end{tcolorbox}

    With this, Theorem~\ref{theo:main_prime} easily follows:

    \begin{scproof}[ of Theorem~\ref{theo:main_prime}]
        Fix a sequence $(\mathcal{F}_n)_{n \in \N}$ of unbounded families of finite subsets of $\G$ %countable family $(S_n)_{n \in \N}$ of infinite subsets of $\G$
        and let $X$ be a nonempty, free, zero-dimensional \sr $\G$-space given by Theorem~\ref{theo:rich}. We can find a free Polish $\G$-flow $Y$ that admits a continuous $\G$-equivariant map $\rho \colon X \to Y$, and hence $\Sub_X(Y) \neq \0$ (this is the only place where we use that $X$ is zero-dimensional; see Proposition~\ref{prop:map_to_free_flow}). Every subflow of $Y$ is free, so it suffices to argue that there exists a subflow $Z \in \Sub_X(Y)$ that is $\mathcal{F}_n$-minimal for all $n \in \N$. To this end, we note that, by Theorem~\ref{theo:generic}, \[\set{Z \in \Sub_X(Y) \,:\, \text{$Z$ is $\mathcal{F}_n$-minimal}}\] is a dense $G_\delta$ subset of $\Sub_X(Y)$ for each $n \in \N$. By the Baire category theorem, the intersection of all these sets is also dense $G_\delta$, so, in particular, it is nonempty, as desired.
    \end{scproof}

    Theorem~\ref{theo:generic} is proved in \S\ref{sec:generic}. We now turn our attention to the proof of Theorem~\ref{theo:rich}, i.e., the existence of a nonempty, free, zero-dimensional \sr $\G$-space.

    \subsection{Separators and examples of \sr $\G$-spaces}\label{subsubsec:example}

    In this subsection we describe a family of simple concrete examples of nonempty, free, zero\=/dimensional \sr $\G$-spaces. %We prove Theorem~\ref{theo:rich} by giving concrete examples of nonempty, free, zero-dimensional \sr $\G$-spaces, which we shall now describe.
    
    
    We begin with some graph-theoretic terminology. Let $G$ be a (simple, undirected) graph. For a set $U \subseteq V(G)$ of vertices, we let $G[U]$ be the subgraph of $G$ \emphd{induced} by $U$, i.e., the graph with vertex set $U$ whose adjacency relation is inherited from $G$. We say that a set $U \subseteq V(G)$ is \emphd{$G$-finite} if every connected component of the graph $G[U]$ is finite.

     Now let $\G \acts X$ be an action of $\G$. Given a set $\Phi \subseteq \G$, the \emphd{Schreier graph} $\Sch(X, \Phi)$ of this action is the graph with vertex set $X$ and edge set
    \[
        E(\Sch(X,\Phi)) \,\defeq\, \big\{\set{x, \sigma \cdot x} \,:\, x \in X \text{ and } \sigma \in \Phi \text{ such that } \sigma \cdot x \neq x\big\}.
    \]
    In the case when $X = \G$ equipped with the left multiplication action $\G \acts \G$, we write \[\Cay(\G, \Phi) \,\defeq\, \Sch(\G, \Phi)\] and call $\Cay(\G,\Phi)$ the \emphd{Cayley graph} of $\G$ corresponding to $\Phi$. Note that if an action $\G \acts X$ is free, then the Schreier graph $\Sch(X, \Phi)$ is obtained by placing a copy of $\Cay(\G,\Phi)$ onto each orbit of $\G$ in $X$. For brevity, we say that a subset $U \subseteq X$ is \emphd{$\Phi$-finite} if it is $\Sch(X,\Phi)$-finite.

    \begin{defn}[Separators]\label{defn:sep}
        Let $s \in \N^+$ and let $\Phi \subset \G$ be a finite subset. An \emphd{$(s,\Phi)$-separator} is a mapping $x \colon \G \to (s+1)$ such that %the sets $x^{-1}(0)$, \ldots, $x^{-1}(s)$ are 
        for all $0 \leq i \leq s$, the set $x^{-1}(s) \subseteq \G$ is
        $\Phi$-finite \ep{with respect to the left multiplication action $\G \acts \G$}. The set of all $(s,\Phi)$-separators is denoted by $\Sep(s,\Phi)$.
    \end{defn}

    \begin{figure}[t]
			\centering
			\begin{tikzpicture}
                \draw[fill=blue!15,draw=white] (-2.9,-2.9) rectangle (3.9,2.9);
                \draw[rounded corners,fill=red!15,draw=lightgray,dashed] (-2-0.25,-2.5-0.25) -- (-2-0.25,2.5+0.25) -- (3+0.25,2.5+0.25) -- (3+0.25,-2.5-0.25) -- cycle;
                \draw[rounded corners,fill=blue!15,draw=lightgray,dashed] (-1.5-0.25, -1-0.25) -- (-1.5-0.25,1.5+0.25) -- (1+0.25,1.5+0.25) -- (1+0.25,1+0.25) -- (2.5+0.25,1+0.25) -- (2.5+0.25,-1.5-0.25) -- (1.5-0.25, -1.5-0.25) -- (0-0.25,-1.5-0.25) -- (0-0.25, -1-0.25) -- cycle;
                \draw[rounded corners,fill=red!15,draw=lightgray,dashed] (-1-0.25,0-0.25) -- (0+0.25,0-0.25) -- (0+0.25,1+0.25) -- (0-0.25, 1+0.25) -- (0-0.25,0.5+0.25) -- (-0.5-0.25,0.5+0.25) -- (-0.5-0.25,0+0.25) -- (-1-0.25,0+0.25) -- cycle;
                \draw[rounded corners,fill=red!15,draw=lightgray,dashed] (1-0.25,0.5+0.25) -- (2+0.25,0.5+0.25) -- (2+0.25,-0.5-0.25) -- (1-0.25,-0.5-0.25) -- cycle;
                %\draw[rounded corners,fill=blue!15,draw=white] (-1.5-0.25,-2+0.25) -- (-0.5+0.25,-2+0.25) -- (-0.5+0.25,-2.5+0.25) -- (2+0.25,-2.5+0.25) -- (2+0.25,-3-0.25) -- (-1.5-0.25,-3-0.25) -- cycle;
                \draw[rounded corners,fill=blue!15,draw=lightgray,dashed] (1.5-0.25,2-0.25) -- (1.5-0.25,2+0.25) -- (2.5+0.25,2+0.25) -- (2.5+0.25,2-0.25) -- cycle;
                %\draw[rounded corners,fill=red!15,draw=white] (-1-0.25,-2.5-0.25) -- (-1-0.25,-2.5+0.25) -- (-1+0.25,-2.5+0.25) -- (-1+0.25,-2.5-0.25) -- cycle;
                \draw[rounded corners,fill=blue!15,draw=lightgray,dashed] (1.5-0.25,0-0.25) -- (1.5-0.25,0+0.25) -- (1.5+0.25,0+0.25) -- (1.5+0.25,0-0.25) -- cycle;
                \draw[rounded corners,fill=blue!15,draw=lightgray,dashed] (-1-0.25,-2-0.25) -- (-1-0.25,-2+0.25) -- (-1+0.25,-2+0.25) -- (-1+0.25,-2-0.25) -- cycle;
   
                \foreach \i in {-5,...,7}
                \draw[gray] (\i/2,-2.9) -- (\i/2,2.9);

                \foreach \j in {-5,...,5}
                \draw[gray] (-2.9,\j/2) -- (3.9,\j/2);

                %\draw[rounded corners,draw=red,very thick] (-2-0.25,-3.5-0.25) -- (-2-0.25,2.5+0.25) -- (3+0.25,2.5+0.25) -- (3+0.25,-3.5-0.25) -- cycle;
                %\draw[rounded corners,draw=blue,very thick] (-1.5-0.25, -1-0.25) -- (-1.5-0.25,1.5+0.25) -- (1+0.25,1.5+0.25) -- (1+0.25,1+0.25) -- (2.5+0.25,1+0.25) -- (2.5+0.25,-1.5-0.25) -- (1.5-0.25, -1.5-0.25) -- (0-0.25,-1.5-0.25) -- (0-0.25, -1-0.25) -- cycle;
                %\draw[rounded corners,draw=red,very thick] (-1-0.25,0-0.25) -- (0+0.25,0-0.25) -- (0+0.25,1+0.25) -- (0-0.25, 1+0.25) -- (0-0.25,0.5+0.25) -- (-0.5-0.25,0.5+0.25) -- (-0.5-0.25,0+0.25) -- (-1-0.25,0+0.25) -- cycle;
                %\draw[rounded corners,draw=red,very thick] (1-0.25,0.5+0.25) -- (2+0.25,0.5+0.25) -- (2+0.25,-0.5-0.25) -- (1-0.25,-0.5-0.25) -- cycle;
                %\draw[rounded corners,draw=blue,very thick] (-1.5-0.25,-2+0.25) -- (-0.5+0.25,-2+0.25) -- (-0.5+0.25,-2.5+0.25) -- (2+0.25,-2.5+0.25) -- (2+0.25,-3-0.25) -- (-1.5-0.25,-3-0.25) -- cycle;
                %\draw[rounded corners,draw=blue,very thick] (1.5-0.25,2-0.25) -- (1.5-0.25,2+0.25) -- (2.5+0.25,2+0.25) -- (2.5+0.25,2-0.25) -- cycle;
                %\draw[rounded corners,draw=red,very thick] (-1-0.25,-2.5-0.25) -- (-1-0.25,-2.5+0.25) -- (-1+0.25,-2.5+0.25) -- (-1+0.25,-2.5-0.25) -- cycle;

                \filldraw[blue] (-2.5,2.5) circle (2pt);
                \filldraw[red] (-2,2.5) circle (2pt);
                \filldraw[red] (-1.5,2.5) circle (2pt);
                \filldraw[red] (-1,2.5) circle (2pt);
                \filldraw[red] (-.5,2.5) circle (2pt);
                \filldraw[red] (0,2.5) circle (2pt);
                \filldraw[red] (.5,2.5) circle (2pt);
                \filldraw[red] (1,2.5) circle (2pt);
                \filldraw[red] (1.5,2.5) circle (2pt);
                \filldraw[red] (2,2.5) circle (2pt);
                \filldraw[red] (2.5,2.5) circle (2pt);
                \filldraw[red] (3,2.5) circle (2pt);
                \filldraw[blue] (3.5,2.5) circle (2pt);

                \filldraw[blue] (-2.5,2) circle (2pt);
                \filldraw[red] (-2,2) circle (2pt);
                \filldraw[red] (-1.5,2) circle (2pt);
                \filldraw[red] (-1,2) circle (2pt);
                \filldraw[red] (-.5,2) circle (2pt);
                \filldraw[red] (0,2) circle (2pt);
                \filldraw[red] (.5,2) circle (2pt);
                \filldraw[red] (1,2) circle (2pt);
                \filldraw[blue] (1.5,2) circle (2pt);
                \filldraw[blue] (2,2) circle (2pt);
                \filldraw[blue] (2.5,2) circle (2pt);
                \filldraw[red] (3,2) circle (2pt);
                \filldraw[blue] (3.5,2) circle (2pt);

                \filldraw[blue] (-2.5,1.5) circle (2pt);
                \filldraw[red] (-2,1.5) circle (2pt);
                \filldraw[blue] (-1.5,1.5) circle (2pt);
                \filldraw[blue] (-1,1.5) circle (2pt);
                \filldraw[blue] (-.5,1.5) circle (2pt);
                \filldraw[blue] (0,1.5) circle (2pt);
                \filldraw[blue] (.5,1.5) circle (2pt);
                \filldraw[blue] (1,1.5) circle (2pt);
                \filldraw[red] (1.5,1.5) circle (2pt);
                \filldraw[red] (2,1.5) circle (2pt);
                \filldraw[red] (2.5,1.5) circle (2pt);
                \filldraw[red] (3,1.5) circle (2pt);
                \filldraw[blue] (3.5,1.5) circle (2pt);

                \filldraw[blue] (-2.5,1) circle (2pt);
                \filldraw[red] (-2,1) circle (2pt);
                \filldraw[blue] (-1.5,1) circle (2pt);
                \filldraw[blue] (-1,1) circle (2pt);
                \filldraw[blue] (-.5,1) circle (2pt);
                \filldraw[red] (0,1) circle (2pt);
                \filldraw[blue] (.5,1) circle (2pt);
                \filldraw[blue] (1,1) circle (2pt);
                \filldraw[blue] (1.5,1) circle (2pt);
                \filldraw[blue] (2,1) circle (2pt);
                \filldraw[blue] (2.5,1) circle (2pt);
                \filldraw[red] (3,1) circle (2pt);
                \filldraw[blue] (3.5,1) circle (2pt);

                \filldraw[blue] (-2.5,.5) circle (2pt);
                \filldraw[red] (-2,.5) circle (2pt);
                \filldraw[blue] (-1.5,.5) circle (2pt);
                \filldraw[blue] (-1,.5) circle (2pt);
                \filldraw[red] (-.5,.5) circle (2pt);
                \filldraw[red] (0,.5) circle (2pt);
                \filldraw[blue] (.5,.5) circle (2pt);
                \filldraw[red] (1,.5) circle (2pt);
                \filldraw[red] (1.5,.5) circle (2pt);
                \filldraw[red] (2,.5) circle (2pt);
                \filldraw[blue] (2.5,.5) circle (2pt);
                \filldraw[red] (3,.5) circle (2pt);
                \filldraw[blue] (3.5,.5) circle (2pt);

                \filldraw[blue] (-2.5,0) circle (2pt);
                \filldraw[red] (-2,0) circle (2pt);
                \filldraw[blue] (-1.5,0) circle (2pt);
                \filldraw[red] (-1,0) circle (2pt);
                \filldraw[red] (-.5,0) circle (2pt);
                \filldraw[red] (0,0) circle (2pt);
                \filldraw[blue] (.5,0) circle (2pt);
                \filldraw[red] (1,0) circle (2pt);
                \filldraw[blue] (1.5,0) circle (2pt);
                \filldraw[red] (2,0) circle (2pt);
                \filldraw[blue] (2.5,0) circle (2pt);
                \filldraw[red] (3,0) circle (2pt);
                \filldraw[blue] (3.5,0) circle (2pt);

                \filldraw[blue] (-2.5,-.5) circle (2pt);
                \filldraw[red] (-2,-.5) circle (2pt);
                \filldraw[blue] (-1.5,-.5) circle (2pt);
                \filldraw[blue] (-1,-.5) circle (2pt);
                \filldraw[blue] (-.5,-.5) circle (2pt);
                \filldraw[blue] (0,-.5) circle (2pt);
                \filldraw[blue] (.5,-.5) circle (2pt);
                \filldraw[red] (1,-.5) circle (2pt);
                \filldraw[red] (1.5,-.5) circle (2pt);
                \filldraw[red] (2,-.5) circle (2pt);
                \filldraw[blue] (2.5,-.5) circle (2pt);
                \filldraw[red] (3,-.5) circle (2pt);
                \filldraw[blue] (3.5,-.5) circle (2pt);

                \filldraw[blue] (-2.5,-1) circle (2pt);
                \filldraw[red] (-2,-1) circle (2pt);
                \filldraw[blue] (-1.5,-1) circle (2pt);
                \filldraw[blue] (-1,-1) circle (2pt);
                \filldraw[blue] (-.5,-1) circle (2pt);
                \filldraw[blue] (0,-1) circle (2pt);
                \filldraw[blue] (.5,-1) circle (2pt);
                \filldraw[blue] (1,-1) circle (2pt);
                \filldraw[blue] (1.5,-1) circle (2pt);
                \filldraw[blue] (2,-1) circle (2pt);
                \filldraw[blue] (2.5,-1) circle (2pt);
                \filldraw[red] (3,-1) circle (2pt);
                \filldraw[blue] (3.5,-1) circle (2pt);

                \filldraw[blue] (-2.5,-1.5) circle (2pt);
                \filldraw[red] (-2,-1.5) circle (2pt);
                \filldraw[red] (-1.5,-1.5) circle (2pt);
                \filldraw[red] (-1,-1.5) circle (2pt);
                \filldraw[red] (-.5,-1.5) circle (2pt);
                \filldraw[blue] (0,-1.5) circle (2pt);
                \filldraw[blue] (.5,-1.5) circle (2pt);
                \filldraw[blue] (1,-1.5) circle (2pt);
                \filldraw[blue] (1.5,-1.5) circle (2pt);
                \filldraw[blue] (2,-1.5) circle (2pt);
                \filldraw[blue] (2.5,-1.5) circle (2pt);
                \filldraw[red] (3,-1.5) circle (2pt);
                \filldraw[blue] (3.5,-1.5) circle (2pt);

                \filldraw[blue] (-2.5,-2) circle (2pt);
                \filldraw[red] (-2,-2) circle (2pt);
                \filldraw[red] (-1.5,-2) circle (2pt);
                \filldraw[blue] (-1,-2) circle (2pt);
                \filldraw[red] (-.5,-2) circle (2pt);
                \filldraw[red] (0,-2) circle (2pt);
                \filldraw[red] (.5,-2) circle (2pt);
                \filldraw[red] (1,-2) circle (2pt);
                \filldraw[red] (1.5,-2) circle (2pt);
                \filldraw[red] (2,-2) circle (2pt);
                \filldraw[red] (2.5,-2) circle (2pt);
                \filldraw[red] (3,-2) circle (2pt);
                \filldraw[blue] (3.5,-2) circle (2pt);

                \filldraw[blue] (-2.5,-2.5) circle (2pt);
                \filldraw[red] (-2,-2.5) circle (2pt);
                \filldraw[red] (-1.5,-2.5) circle (2pt);
                \filldraw[red] (-1,-2.5) circle (2pt);
                \filldraw[red] (-.5,-2.5) circle (2pt);
                \filldraw[red] (0,-2.5) circle (2pt);
                \filldraw[red] (.5,-2.5) circle (2pt);
                \filldraw[red] (1,-2.5) circle (2pt);
                \filldraw[red] (1.5,-2.5) circle (2pt);
                \filldraw[red] (2,-2.5) circle (2pt);
                \filldraw[red] (2.5,-2.5) circle (2pt);
                \filldraw[red] (3,-2.5) circle (2pt);
                \filldraw[blue] (3.5,-2.5) circle (2pt);
			\end{tikzpicture}
        \caption{A piece of a $(1, \Phi)$-separator $x \colon \Z^2 \to 2$ in the group $\Z^2$, where $\Phi \defeq \set{(1,0),(0,1)}$ is the standard generating set. Here each element $\gamma \in \Z^2$ is labeled red or blue depending on the value $x(\gamma)$, splitting the Cayley graph $\Cay(\Z^2, \Phi)$ into finite monochromatic components.  %A fragment of a $(1, \set{(1,0),(0, 1)})$-separator in $\Z^2$.
        }\label{fig:separator}
	\end{figure}

    Definition~\ref{defn:sep} is illustrated in Figure~\ref{fig:separator}. It is easy to see that for $s \in \N^+$ and finite $\Phi \subset \G$, $\Sep(s,\Phi)$ is a shift-invariant dense $G_\delta$ subset of $(s+1)^\G$. In particular, $\Sep(s,\Phi)$ is a nonempty zero-dimensional Polish space (in the subspace topology) equipped with the shift action $\G \acts \Sep(s,\Phi)$. 

    \begin{defn}[Asymptotic separators]
    For the remainder of the paper, we fix an enumeration $(\Phi_n)_{n \in \N}$ of the finite subsets of $\G$ in which every set appears infinitely often. Given $s \in \N^+$, we let
    \[
        \Sep(s) \,\defeq\, \prod_{n \in \N} \Sep(s, \Phi_n).
    \]
    The space $\Sep(s)$ is equipped with the product topology and the diagonal action of $\G$. We call the elements of $\Sep(s)$ \emphd{asymptotic $s$-separators}.
    \end{defn}
    
    Note that the action $\G \acts \Sep(s)$ is often free:

    \begin{prop}\label{prop:torsion-free}
        If $\G$ is torsion-free, then for all $s \in \N^+$, the action $\G \acts \Sep(s)$ is free.
    \end{prop}
    \begin{scproof}
        Take any $\gamma \in \G \setminus \set{\iden}$ and suppose that $\gamma \cdot x = x$ for some $x = (x_n)_{n \in \N} \in \Sep(s)$. Pick $n \in \N$ such that $\gamma \in \Phi_n$ and let $i \defeq x_n(\iden)$. Since $\gamma \cdot x_n = x_n$, we have $x_n(\gamma^m) = i$ for all $m \in \Z$. But then the infinite set $\langle \gamma \rangle$ is contained in a single component of the subgraph of $\Cay(\G, \Phi_n)$ induced by $x_n^{-1}(i)$, which contradicts the fact that $x_n \in \Sep(s, \Phi_n)$. %is a contradiction.
    \end{scproof}

    In general, the {free part} $\Free(\Sep(s))$ of $\Sep(s)$, i.e., the set of all points $x \in \Sep(s)$ with trivial stabilizers, is a dense $G_\delta$ subset of $\Sep(s)$, because both $\Sep(s)$ and the free part of $((s+1)^\G)^\N$ are dense $G_\delta$ subsets of $((s+1)^\G)^\N$. Hence, $\Free(\Sep(s))$ is a nonempty zero-dimensional Polish space, and, by definition, the group $\G$ acts on $\Free(\Sep(s))$ freely. We will show that $\Free(\Sep(s))$ is also \sr, and therefore %: %We can now describe our examples of zero-dimensional \sr $\G$-spaces:
    %We claim that $\Free(\Sep(s))$ is \sr:
    %We claim that the space $\Free(\Sep(s))$
    it is an example confirming Theorem~\ref{theo:rich}:

    \begin{tcolorbox}
    \begin{theo}\label{theo:example}
        For any $s \in \N^+$, $\Free(\Sep(s))$ is an \sr $\G$-space.
    \end{theo}
    \end{tcolorbox}

    It is clear from the above discussion that Theorem~\ref{theo:example} implies Theorem~\ref{theo:rich}. In order to explain how Theorem~\ref{theo:example} is proved, we need to say a few words about the relationship between the space $\Sep(s)$ of asymptotic $s$-separators and asymptotic separation index. % of group actions.

    \subsection{Asymptotic separation index and its continuous version}\label{subsubsec:asi}

    For simplicity, we will introduce asymptotic separation index for free group actions, although it is not the most general context in which it can be defined; for example, the original definition in \cite{CJMSTD} is given for so-called \emph{Borel extended metric spaces}.

    \begin{defn}[{Asymptotic separation index~\cite[Defn.~3.2]{CJMSTD}}]\label{defn:asi}
        Let $\G \acts X$ be a free Borel action of $\G$ on a Polish space $X$. The \emphd{asymptotic separation index} of $X$, in symbols $\asi(X)$, is the smallest integer $s \in \N$ such that for every finite set $\Phi \subset \G$, there exists a partition $X = X_0 \sqcup \ldots \sqcup X_s$ into $\Phi$-finite Borel sets. If there is no such $s \in \N$, we let $\asi(X) \defeq \infty$.
    \end{defn}

    We need %Our proof of Theorem~\ref{theo:example} relies on
    a variant of Definition~\ref{defn:asi} in which the sets $X_0$, \ldots, $X_s$ are %required to be
    {clopen} rather than Borel:

    \begin{defn}[Continuous asymptotic separation index]\label{defn:casi}
        Let $X$ be a free zero-dimensional Polish $\G$-space. % $\G \acts X$ be a continuous action of $\G$ on a zero-dimensional Polish space $X$.
        The \emphd{continuous asymptotic separation index} of $X$, in symbols $\casi(X)$, is the smallest integer $s \in \N$ such that for every finite set $\Phi \subset \G$, there exists a partition $X = X_0 \sqcup \ldots \sqcup X_s$ into $\Phi$-finite clopen sets. If there is no such $s \in \N$, we let $\casi(X) \defeq \infty$.
    \end{defn}

    We can equivalently describe the (continuous) asymptotic separation index of $X$ by relating $X$ to the spaces of asymptotic $s$-separators. To this end, note that for any action $\G \acts X$ and $k \in \N^+$, we have a one-to-one correspondence
    \[
	    \big\{\text{functions $X \to k$}\big\} \quad \longleftrightarrow \quad \big\{\text{$\G$-equivariant maps $X \to k^\G$}\big\}.
	\]
	Namely, each function $f \colon X \to k$ gives rise to the $\G$-equivariant \emphd{coding map} $\pi_f \colon X \to k^\G$ via
	\[
	    \big(\pi_f(x)\big)(\gamma) \,\defeq\, f(\gamma \cdot x) \quad \text{for all } x \in X \text{ and } \gamma \in \G.
	\]
	Conversely, if $\pi \colon X \to k^\G$ is $\G$-equivariant, then the function %\pi = \pi_f$ for the function
    $f \colon X \to k$ such that $\pi = \pi_f$ is given by \[f(x) \,\defeq\, (\pi(x))(\iden).\]  %$f \defeq (x \mapsto (\pi(x))(\iden))$. 
    %\[
    %    f \,=\, \big(x \mapsto (\pi(x))(\iden)\big).
    %\]
    %Note that if $X$ is a topological space and an action $\G \acts X$ is continuous, then a function $f \colon X \to k$ is continuous if and only if so is the coding map $\pi_f$. 
    Using this correspondence, we obtain the following:

    \begin{prop}[Dynamical view of continuous asymptotic separation index]\label{prop:map_to_Sep}
        Let $X$ be a free zero-dimensional Polish $\G$-space. %$\G \acts X$ be a free continuous action of $\G$ on a zero-dimensional Polish space $X$.
        The following statements are equivalent for all $s \in \N^+$:
        \begin{enumerate}[label=\ep{\normalfont\roman*}]
            \item\label{item:casi} $\casi(X) \leq s$,
            \item\label{item:term} for every finite set $\Phi \subset \G$, there exists a continuous $\G$-equivariant map $X \to \Sep(s,\Phi)$,
            \item\label{item:product} there exists a continuous $\G$-equivariant map $X \to \Sep(s)$.
        \end{enumerate}
    \end{prop}
    \begin{scproof}
        The equivalence \ref{item:term} $\Longleftrightarrow$ \ref{item:product} is clear from the definition of $\Sep(s)$. Given a finite set $\Phi \subset \G$ and a partition $X = X_0 \sqcup \ldots \sqcup X_s$ of $X$ into $\Phi$-finite clopen sets, we obtain a continuous $\G$-equivariant map $\pi \colon X \to \Sep(s,\Phi)$ by setting $f(x) \defeq i$ for all $x \in X_i$ and taking $\pi \defeq \pi_f$. This proves \ref{item:casi} $\Longrightarrow$ \ref{item:term}. Conversely, given a continuous $\G$-equivariant map $\pi \colon X \to \Sep(s,\Phi)$, we can form a partition of $X$ into $\Phi$-finite clopen sets $X_0$, \ldots, $X_s$ by letting $X_i \defeq \set{x \in X \,:\, (\pi(x))(\iden) = i}$.
    \end{scproof}

    In particular, $\casi(\Free(\Sep(s))) \leq s$ since the inclusion map $\Free(\Sep(s)) \hookrightarrow \Sep(s)$ is continuous and $\G$-equivariant. The same proof gives a version of Proposition~\ref{prop:map_to_Sep} for the ordinary asymptotic separation index, where the word ``continuous'' is replaced by ``Borel.'' %\todo{...} %(though we will not need it here), in which the word ``continuous'' is replaced everywhere by ``Borel.'' % For completeness, we record the version of Proposition~\ref{prop:map_to_Sep} for the ordinary asymptotic separation index as well, although it will not be used in this paper:

    %\begin{prop}[Dynamical view of asymptotic separation index]
    %    Let $\G \acts X$ be a free Borel action of $\G$ on a Polish space $X$. The following statements are equivalent for all $s \in \N^+$:
    %    \begin{enumerate}[label=\ep{\normalfont\roman*}]
    %        \item $\asi(X) \leq s$,
    %        \item for every finite set $F \subset \G$, there exists a Borel $\G$-equivariant map $X \to \Sep(s,F)$,
    %        \item there exists a Borel $\G$-equivariant map $X \to \Sep(s)$.
    %    \end{enumerate}
    %\end{prop}
    %\begin{scproof}
    %    Same as the proof of Proposition~\ref{prop:map_to_Sep}, \emph{mutatis mutandis}.
    %\end{scproof}

    %In \cite{CJMSTD}, Conley \emph{et al.}~showed that for many groups $\G$ (e.g., polycyclic groups, the lamplighter group $\Z_2 \wr \Z$, and the Baumslag--Solitar group $BS(1,2)$, to name a few), all their Borel actions on Polish spaces have asymptotic separation index at most $1$.
    In applications, the important distinction is usually between actions $\G \acts X$ with finite asymptotic separation index and those for which $\asi(X) = \infty$.\footnote{As an aside, we remark that there are currently no known examples with $1 < \asi(X) < \infty$ \cite[3191]{CJMSTD}.} In particular, it turns out that when $\asi(X) < \infty$, many otherwise intractable combinatorial problems can be solved on $X$ in a Borel way \cite{ASIalgorithms,BW,BWKonig,CJMSTD,WeilacherFinDim}. Our insight is that 
    %We observe that
    upgrading the assumption from $\asi(X) < \infty$ to $\casi(X) < \infty$ may yield a {continuous} solution in place of a Borel one. In the next subsection we shall explain how this idea is used to find continuous $\G$-equivariant maps \[\pi \colon \Free(\Sep(s)) \to \Free(\Sep(s))\] witnessing that $\Free(\Sep(s))$ is \sr.

    \subsection{Combinatorial core of the problem and the role of the Lov\'asz Local Lemma}\label{subsubsec:combi+LLL}

    Suppose $X$ is a free zero-dimensional Polish $\G$-space. %space with a free continuous action $\G \acts X$.
    By Proposition~\ref{prop:map_to_Sep}, if $\casi(X) \leq s$, then there exists a continuous $\G$-equivariant map $\pi \colon X \to \Sep(s)$. The technical substance of Theorem~\ref{theo:example} is in the following lemma, which asserts that the map $\pi$ can additionally be chosen to have properties analogous to those in Definition~\ref{defn:amply_syndetic}: % to make \todo{...} %the preimages of finitely many given open sets syndetic:
    
    \begin{lemma}\label{lemma:key}
        For every $s \in \N^+$ and a finite tuple $U_1$, \ldots, $U_k \subseteq ((s+1)^\G)^\N$ of nonempty open sets, there is $n = n(U_1,\ldots, U_k) \in \N$ such that for all finite $F \subset \G$ with $|F| \geq n$, the following holds:
        
        \smallskip
        
        \noindent Every free zero-dimensional Polish $\G$-space $X$ %with a free continuous action $\G \acts X$
        such that $\casi(X) \leq s$ admits a continuous $\G$-equivariant map $\pi \colon X \to \Sep(s)$ such that the sets $\pi^{-1}(U_1)$, \ldots, $\pi^{-1}(U_k)$ are $F$-syndetic.
        
        %Let $X$ be a zero-dimensional Polish space equipped with a free continuous action $\G \acts X$ and suppose that $\casi(X) \leq s \in \N^+$. Then, for every finite tuple
        %\[
        %    U_1, \ \ldots, \ U_k \ \subseteq \  ((s+1)^\G)^\N
        %\]
        %of nonempty open sets, there is $n \in \N$ such that for all finite $F \subset \G$ with $|F| \geq n$,  %and any infinite set $S \subseteq \G$,
        %there exists a continuous $\G$-equivariant map $\pi \colon X \to \Sep(s)$ %\[\pi \colon X \to \Sep(s)\] such that the sets $\pi^{-1}(U_1)$, \ldots, $\pi^{-1}(U_k)$ are $S$-syndetic.
        %making the sets $\pi^{-1}(U_1)$, \ldots, $\pi^{-1}(U_k)$ $F$-syndetic.
        %\[
        %    F^{-1} \cdot \pi^{-1}(U_1) \,=\, \cdots \,=\, F^{-1}\cdot \pi^{-1}(U_k) \,=\, X.
        %\]
    \end{lemma}

    \begin{remk}
        Note that
        in the statement of Lemma~\ref{lemma:key}, $U_1$, \ldots, $U_k$ are taken to be
        nonempty open subsets of $((s+1)^\G)^\N$ rather than of $\Sep(s)$. This is done purely for convenience and makes no difference to the content of the lemma because $\Sep(s)$ is dense in $((s+1)^\G)^\N$.
    \end{remk}

    Lemma~\ref{lemma:key} takes us most of the way toward proving Theorem~\ref{theo:example}. Indeed, suppose we are given %an infinite subset $S \subseteq \G$ and
    nonempty open sets $U_1$, \ldots, $U_k  \subseteq  ((s+1)^\G)^\N$ and let $n= n(U_1,\ldots, U_k) \in \N$ be the integer produced by Lemma~\ref{lemma:key}. Since $\casi(\Free(\Sep(s))) \leq s$, %we may apply Lemma~\ref{lemma:key} with $X = \Free(\Sep(s))$ 
    %to obtain an integer $n = n(U_) \in \N$ such that
    for each finite set $F \subset \G$ of size at least $n$, there exists a continuous $\G$-equivariant map $\pi \colon \Free(\Sep(s)) \to \Sep(s)$ such that the sets $\pi^{-1}(U_1)$, \ldots, $\pi^{-1}(U_k)$ are $F$-syndetic.
    When $\G$ is torsion-free, this shows that the space $\Free(\Sep(s)) = \Sep(s)$ is \sr, as desired. In general, we need to adjust the map $\pi$ to ensure that its range is included in the free part of $\Sep(s)$. This adjustment is performed in \S\ref{sec:final}.

    To prove Lemma~\ref{lemma:key}, given a large finite set $F \subset \G$, we first 
    judiciously choose $F$-syndetic
    subsets $A_1$, \ldots, $A_k \subseteq X$ % such that
        %\[
        %    F^{-1} \cdot A_1 = \cdots = F^{-1} \cdot A_k = X,
        %\]
        and then construct a continuous $\G$-equivariant map $\pi \colon X \to \Sep(s)$ such that $\pi^{-1}(U_i) \supseteq A_i$ for all $1 \leq i \leq k$. To implement this strategy, we need to answer the question:
    \begin{tcolorbox}[width=\linewidth, sharp corners=all, colback=white!95!black, boxrule=0pt,colframe=white] 
        Under what assumptions on subsets $A_1$, \ldots, $A_k \subseteq X$ can we guarantee that there exists a continuous $\G$-equivariant map $\pi \colon X \to \Sep(s)$  with $\pi(x) \in U_i$ for all $x \in A_i$?
    \end{tcolorbox}
    %\begin{quote}
    %    Under what assumptions on sets $A_1$, \ldots, $A_k \subseteq X$ can we guarantee that there exists a continuous $\G$-equivariant map $\pi \colon X \to \Sep(s)$  with $\pi(x) \in U_i$ for all $x \in A_i$? %What properties of subsets $A_1$, \ldots, $A_k \subseteq X$ suffice to ensure that there exists a continuous $\G$-equivariant map $\pi \colon X \to \Sep(s)$  with $\pi(x) \in U_i$ for all $x \in A_i$?
    %\end{quote}
    We show it is enough for $A_1$, \ldots, $A_k$ to be sufficiently ``spaced out,'' in the following sense:

    \begin{defn}[$\Phi$-spaced sets]
        Given an action $\G \acts X$ and a finite set $\Phi \subset \G$, we say that a subset $A \subseteq X$ is \emphd{$\Phi$-spaced} if $(\Phi \cdot x) \cap (\Phi \cdot y) = \0$ for all distinct $x$, $y \in A$. %When we call a subset $A \subseteq \G$ of the group itself $F$-spaced, we mean so with respect to the left multiplication action $\G \acts \G$. %Note that if $F \subset \G$ is a finite set, then every infinite set $S \subseteq \G$ has an infinite $F$-spaced subset $S' \subseteq S$.
    \end{defn}

    \begin{lemma}\label{lemma:from_separation_to_map}
       For every $s \in \N^+$ and a finite tuple $U_1$, \ldots, $U_k \subseteq ((s+1)^\G)^\N$ of nonempty open sets, there is a finite set $\Phi = \Phi(U_1, \ldots, U_k) \subset \G$ with the following property:

       \smallskip
       
       \noindent Suppose $X$ is a free zero-dimensional Polish $\G$-space such that $\casi(X) \leq s$, and let $A_1$, \ldots, $A_k \subseteq X$ be disjoint clopen sets whose union $A \defeq A_1 \sqcup \ldots \sqcup A_k$ is $\Phi$-spaced. Then there exists a continuous $\G$-equivariant map $\pi \colon X \to \Sep(s)$ such that $\pi(x) \in U_i$
        for all $1 \leq i \leq k$ and $x \in A_i$.
    \end{lemma}

    Lemma~\ref{lemma:from_separation_to_map} is proved in \S\ref{sec:spaced}, where the desired map $\pi$ is built explicitly. Assuming Lemma~\ref{lemma:from_separation_to_map}, the following suffices to complete the proof of Lemma~\ref{lemma:key}:

    \begin{lemma}\label{lemma:sparse_syndetic}
        For any $s$, $k \in \N^+$ and a finite set $\Phi \subset \G$, there is $n = n(s,k, \Phi) \in \N$ such that for all finite $F \subset \G$ with $|F| \geq n$, the following holds:

        \smallskip
        
        \noindent For every free zero-dimensional Polish $\G$-space $X$ such that $\casi(X) \leq s$, there exist disjoint $F$-syndetic clopen subsets $A_1$, \ldots, $A_k \subseteq X$ whose union %such that 
        %\begin{enumerate}[label=\ep{\normalfont\roman*}]
          %  \item
         %each set $A_i$ is $F$-syndetic and
          %  \item the set
          $A \defeq A_1 \sqcup \ldots \sqcup A_k$ is $\Phi$-spaced.
        %\end{enumerate}
    \end{lemma}

    \begin{scproof}[ of Lemma~\ref{lemma:key}]
        Let $s$ and $U_1$, \ldots, $U_k$ be as in Lemma~\ref{lemma:key}. Let $\Phi = \Phi(U_1,\ldots, U_k) \subset \G$ be the finite set given by Lemma~\ref{lemma:from_separation_to_map} and let $n = n(s,k, \Phi) \in \N$ be given by  Lemma~\ref{lemma:sparse_syndetic}. Take any finite set $F \subset \G$ with $|F| \geq n$ and an arbitrary free zero-dimensional Polish $\G$-space $X$ such that $\casi(X) \leq s$. By Lemma~\ref{lemma:sparse_syndetic}, there exist disjoint $F$-syndetic clopen subsets $A_1$, \ldots, $A_k \subseteq X$ whose union $A \defeq A_1 \sqcup \ldots \sqcup A_k$ is $\Phi$-spaced. % :
        %\begin{enumerate}[label=\ep{\normalfont\roman*}]
        %    \item each set $A_i$ is $S$-syndetic, and
        %    \item the set $A \defeq A_1 \sqcup \ldots \sqcup A_k$ is $F$-spaced.
        %\end{enumerate}
        By Lemma~\ref{lemma:from_separation_to_map}, we have a continuous $\G$-equivariant map $\pi \colon X \to \Sep(s)$ with $\pi(x) \in U_i$ for all $x \in A_i$, i.e., %. In other words,
        $\pi^{-1}(U_i) \supseteq A_i$, which is $F$-syndetic, as desired.
    \end{scproof}

    \begin{remk}
        Lemma~\ref{lemma:sparse_syndetic} may badly fail without the assumption $\casi(X) < \infty$. For example, let $\mathbb{F}_2$ be the free group with generators $a$, $b$ and let $X$ be the free part of the shift action $\mathbb{F}_2 \acts 2^{\mathbb{F}_2}$. Then, for $\Phi \defeq \set{\iden, a}$ and for any finite set $F \subset \langle b \rangle$ (no matter how large), there is no Borel---let alone clopen---set $A \subseteq X$ that is both $F$-syndetic and $\Phi$-spaced. This is a consequence of a result of Marks \cite[Thm.~1.6]{Marks}, who showed that %there is no Borel (let alone clopen) set $A \subseteq X$ that is both $\langle a \rangle $-syndetic and $\set{b}$-spaced. In fact,
        for any Borel set $A \subseteq X$, there exists an element $x \in X$ such that either $\langle a \rangle \cdot x \subseteq A$ or $(\langle b \rangle \cdot x) \cap A = \0$.
        %, and hence $A$ is not $F$-syndetic with respect to any $F \subset \langle b \rangle$ 
        %
        %
        %$(\langle a \rangle \cdot x) \cap A = \0$, which means that $A$ is not $F$-syndetic with respect to any subset $F$ of the infinite set $\langle a \rangle$, or $\langle b \rangle \cdot x \subseteq A$---and hence $A$ is far from being $\set{\iden, b}$-spaced
        %\cite[Thm.~1.6]{Marks}.
    \end{remk}

    The proof of Lemma~\ref{lemma:sparse_syndetic} is where the Lov\'asz Local Lemma comes into play. %The precise statement of the LLL is somewhat technical, so we postpone it until \S\todo{ref}.  
    Consider the following purely combinatorial fact (which is implied by Lemma~\ref{lemma:sparse_syndetic}):

    \begin{quote}
        \textsl{For all finite $\Phi \subset \G$ and sufficiently large finite $F \subset \G$, %an infinite set $S \subseteq \G$,
        there is a set $A \subseteq \G$ that is both $F$-syndetic and $\Phi$-spaced \ep{with respect to the left multiplication action $\G \acts \G$}.}
    \end{quote}

    \noindent This statement has a routine proof using the LLL, and as far as we are aware, there is no other known way to establish it in full generality. This means that, to prove Lemma~\ref{lemma:sparse_syndetic}, we must rely on some version of the LLL; furthermore, this version must be able to produce \emph{clopen} sets with the desired combinatorial properties. As mentioned previously, Weilacher and the first named author \cite{BW} gave a {Borel} version of the LLL that can be used in the context of an action $\G \acts X$ with $\asi(X) < \infty$. In \S\ref{sec:CLLL}, we observe that essentially the same argument yields a {continuous} version of the LLL under the condition $\casi(X) < \infty$, which can then be applied to deduce Lemma~\ref{lemma:sparse_syndetic}.

    \begin{figure}[t]
			\centering
			\begin{tikzpicture}[yscale=0.9]
                    \node[draw, rectangle, align=center] (MainThm) at (0,0) {\small Theorem~\ref{theo:main_prime}: \\ \small An $\mathcal{F}_n$-minimal $\G$-flow for all $n \in \N$};

                    %\node[draw, rectangle,align=center] (Corl) at (0,-2) {\small Corollary~\ref{corl:subshift}: \\ \small An $\mathcal{F}_n$-minimal subshift for all $n \in \N$};

                    %\draw [-{Stealth},very thick] (MainThm) to node[midway,anchor=east] {\small\S\ref{sec:subshift}}  (Corl);

                    \node[draw, rectangle, align=center] (Generic) at (-3,2.5) {\small Theorem~\ref{theo:generic}: \\ \small $\mathcal{F}$-minimal subflows are generic};

                    \node[draw, rectangle, align=center] (Ample) at (4,2.5) {\small Theorem~\ref{theo:rich}/Theorem~\ref{theo:example}: \\ \small An \sr $\G$-space, namely $\Free(\Sep(s))$};

                    \draw [very thick] (Generic) to (0,1.5);
                    \draw [very thick] (Ample) to (0,1.5);
                    \draw [-{Stealth},very thick] (0,1.5) to node[midway,anchor=east] {\small\S\ref{subsubsec:ample}}  (MainThm);

                    \node[left=0.1cm of Generic] {\small\S\ref{sec:generic}};

                    %\node[draw, rectangle, align=center] (Example) at (3,4.5) {\small Theorem~\ref{theo:example}: \\ \small $\Free(\Sep(s))$ is \sr};

                    %\draw [-{Stealth},very thick] (Example) to node[midway,anchor=east] {\small\S\ref{subsubsec:example}}  (Ample);

                    \node[draw, rectangle, align=center] (Key) at (4,4.5) {\small Lemma~\ref{lemma:key}: \\ \small Maps to $\Sep(s)$ making open sets $F$-syndetic};

                    \draw [-{Stealth},very thick] (Key) to node[midway,anchor=east] {\small\S\ref{sec:final}}  (Ample);

                    \node[draw, rectangle, align=center] (Spaced) at (6.5,7) {\small Lemma~\ref{lemma:from_separation_to_map}: \\ \small Maps to $\Sep(s)$ with $\Phi$-spaced constraints};

                    \node[draw, rectangle, align=center] (Combi) at (-1,7) {\small Lemma~\ref{lemma:sparse_syndetic}: \\ \small $F$-syndetic $\Phi$-spaced sets};

                    \draw [very thick] (Combi) to (4,6);
                    \draw [very thick] (Spaced) to (4,6);
                    \draw [-{Stealth},very thick] (4,6) to node[midway,anchor=east] {\small\S\ref{subsubsec:combi+LLL}}  (Key);

                    \node[left=0.1cm of Spaced] {\small\S\ref{sec:spaced}};

                    \node[draw, rectangle, align=center] (LLL) at (-1,9) {\small Theorem~\ref{theo:cont_LLL}: \\ \small Continuous LLL under finite $\casi$ assumption};

                    \draw [-{Stealth},very thick] (LLL) to node[midway,anchor=east] {\small\S\ref{subsec:synd_and_spaced}}  (Combi);

                    \node[left=0.1cm of LLL] {\small\S\ref{sec:CLLL}};
			\end{tikzpicture}
        \caption{A flowchart for the proof of Theorem~\ref{theo:main_prime}.}\label{fig:flowchart}
	\end{figure}

    \subsection{Summary and a road map for the remainder of the paper}

    Figure~\ref{fig:flowchart} shows the structure of the proof of Theorem~\ref{theo:main_prime} described above. The remainder of this paper is organized as follows. After some preliminary remarks in \S\ref{sec:prelims}, we prove the remaining facts needed to deduce Theorem~\ref{theo:main_prime}, namely:
    \begin{itemize}
        \item
        Theorem~\ref{theo:generic} in \S\ref{sec:generic},
        \item
        Lemma~\ref{lemma:from_separation_to_map} in \S\ref{sec:spaced}, and
        \item
        a continuous variant of the LLL together with its consequence, Lemma~\ref{lemma:sparse_syndetic}, in \S\ref{sec:CLLL}. 
    \end{itemize}
    %proceed to prove Theorem~\ref{theo:generic} in \S\ref{sec:generic}. Then, in \S\ref{sec:spaced}, we verify Lemma~\ref{lemma:from_separation_to_map}. Next, in \S\ref{sec:CLLL}, we introduce our continuous version of the LLL and use it to establish Lemma~\ref{lemma:sparse_syndetic}.
    As explained in \S\ref{subsubsec:combi+LLL}, Lemma~\ref{lemma:from_separation_to_map} and Lemma~\ref{lemma:sparse_syndetic} combined imply Lemma~\ref{lemma:key}, which we use in \S\ref{sec:final} to complete the proof of Theorem~\ref{theo:example} and hence also of Theorem~\ref{theo:rich}. As explained in \S\ref{subsubsec:ample}, Theorems~\ref{theo:rich} and \ref{theo:generic} yield our main result, Theorem~\ref{theo:main_prime}.
    
    In the remaining sections, \S\S\ref{sec:subshift}--\ref{subsec:subshift_appl}, we establish the corollaries of our main result. Namely, we prove Corollary~\ref{corl:subshift} (the existence of subshifts with minimal subdynamics) in \S\ref{sec:subshift}, apply our results to construct disjoint $\G$-flows in \S\ref{sec:disjoint}, and study Borel complete sections in \S\ref{subsec:subshift_appl}. % use \todo{edit}
    %Finally, in \S\ref{sec:subshift} we show how  to deduce Corollary~\ref{corl:subshift} from Theorem~\ref{theo:main_prime} and then use it to prove Corollaries~\ref{corl:DST1} and \ref{corl:DST}.
    %
    %
    %We then turn our attention to Theorem~\ref{theo:rich}. %, which is proved in \S\ref{sec:construction}.
    %We review the relevant background concerning asymptotic separation index in \S\ref{subsec:asi+example} and then give an explicit example of an \sr $\G$-space in Theorem~\ref{theo:example}. Next, we state and prove the continuous variant of the Lov\'asz Local Lemma needed for the proof of Theorem~\ref{theo:example} in \S\ref{subsec:CLLL}. Finally, we complete the proof of Theorem~\ref{theo:example} in \S\ref{subsec:finish}.

    \section{Remarks on zero-dimensional spaces}\label{sec:prelims}

    %This section collects some basic definitions and facts. This material is very standard, and the reader is encouraged to skip this section and only consult it if and when needed.

    %\subsection{Natural numbers}

    %We write $\N = \set{0,1,2\ldots}$ and $\N^+ = \set{1,2,3,\ldots}$. Each natural number $k \in \N$ with is identified with the $k$-element set $\set{i \in \N \,:\, i < k}$ and given the discrete topology. 

    %\subsection{The free part of an action}

    %Given an action $\G \acts X$, we let $\Free(X)$ be the \emphd{free part} of $X$, i.e., the set of all points $x \in X$ with trivial stabilizers. In other words, $\Free(X)$ is the largest $\G$-invariant subset of $X$ on which the action of $\G$ is free. Note that when $X$ is a Polish space and the action $\G \acts X$ is continuous, $\Free(X)$ is a $G_\delta$ subset of $X$, and hence it is also a Polish space. %\cite[Thm.~3.11]{KechrisDST}.

    %\subsection{Shifts and coding maps}\label{subsec:shift}

    %We identify each natural number $k \in \N$ with the $k$-element set $\set{i \in \N \,:\, i < k}$ and equip it with the discrete topology.
    
    %For $k \in \N^+$, the \emphd{shift} $\G \acts k^\G$ is the action of $\G$ on the product space $k^\G$ given by
    %\[
    %    (\gamma \cdot x)(\delta) \,\defeq\, x(\delta \gamma) \quad \text{for all $x \colon \G \to k$ and $\gamma$, $\delta \in \G$}.
    %\]
    %This makes $k^\G$ a zero-dimensional Polish $\G$-flow. For an action $\G \acts X$, %there is
    %we have a %natural one-to-one
    %correspondence
	%\[
	%    \big\{\text{functions $X \to k$}\big\} \quad \longleftrightarrow \quad \big\{\text{$\G$-equivariant maps $X \to k^\G$}\big\}.
	%\]
	%Namely, each function $f \colon X \to k$ gives rise to the $\G$-equivariant \emphd{coding map} $\pi_f \colon X \to k^\G$ via
	%\[
	%    \big(\pi_f(x)\big)(\gamma) \,\defeq\, f(\gamma \cdot x) \quad \text{for all } x \in X \text{ and } \gamma \in \G.
	%\]
	%Conversely, if $\pi \colon X \to k^\G$ is $\G$-equivariant, then $\pi = \pi_f$ for the function
    %\[
    %    f \,=\, \big(x \mapsto (\pi(x))(\iden)\big).
    %\]
    %Note that if $X$ is a topological space and an action $\G \acts X$ is continuous, then a function $f \colon X \to k$ is continuous if and only if so is the coding map $\pi_f$. 
    %Shift actions and coding maps are central to the field of \emph{symbolic dynamics}; see \cite{subshifts1,subshifts2} for more information. % Given any continuous action $\G \acts X$ on a topological space $X$, there is a natural one-to-one correspondence
	%\[
	%    \big\{\text{continuous functions $X \to k$}\big\} \quad \longleftrightarrow \quad \big\{\text{continuous $\G$-equivariant maps $X \to k^\G$}\big\}.
	%\]
	%Namely, each continuous function $f \colon X \to k$ gives rise to the \emphd{coding map} $\pi_f \colon X \to k^\G$ given by
	%\[
	%    \big(\pi_f(x)\big)(\gamma) \,\defeq\, f(\gamma \cdot x) \quad \text{for all } x \in X \text{ and } \gamma \in \G.
	%\]
	%Conversely, if $\pi \colon X \to k^\G$ is continuous and $\G$-equivariant, then $\pi = \pi_f$ for the function
    %\[
    %    f \,=\, \big(x \mapsto (\pi(x))(\iden)\big).
    %\]
    %Shift actions and coding maps are central to the field of \emph{symbolic dynamics}; see \todo{cite}. 

    %\subsection{Zero-dimensional spaces}

    Here we record two simple facts about zero-dimensional spaces. The first of these facts was used to derive Theorem~\ref{theo:main_prime} from Theorems~\ref{theo:rich} and \ref{theo:generic} in \S\ref{subsubsec:ample}.

    \begin{prop}\label{prop:map_to_free_flow}
        If $X$ is a free zero-dimensional Polish $\G$-space,  %$\G \acts X$ be a free continuous action of $\G$ on a zero-dimensional Polish space $X$,
        then there exists a Polish $\G$-flow $Y$ with a continuous $\G$-equivariant map $\rho \colon X \to Y$.
    \end{prop}
    \begin{scproof}
        Let $(\gamma_n)_{n \in \N}$ be an enumeration of the non-identity elements of $\G$. By \cite[Lem.~2.3]{Ber_cont}, for each $n \in \N$, there is a continuous function $f_n \colon X \to 3$ such $f_n(x) \neq f_n(\gamma_n \cdot x)$ for all $x \in X$. (This is a continuous refinement of a result of Kechris, Solecki, and Todorcevic \cite[Prop.~4.6]{KST}.) Let
        \[
            Y_n \,\defeq\, \big\{y \in 3^\G \,:\, y(\delta) \neq y(\gamma_n \delta) \text{ for all } \delta \in \G\big\}.
        \]
        Then $Y_n \subset 3^\G$ is a subshift, % closed $\G$-invariant subset (i.e., a subflow) of $3^\G$,
        and we have the coding map $\pi_{f_n} \colon X \to Y_n$. Let $Y \defeq \prod_{i \in \N} Y_i$ (equipped with the diagonal action of $\G$) and define a map $\rho \colon X \to Y$ by
        \[
            \rho(x) \,\defeq\, (\pi_{f_i}(x))_{i \in \N} \,\in\, Y.
        \]
        It remains to note that the action $\G \acts Y$ is free, since for each $n \in \N$ and $y = (y_i)_{i \in \N} \in Y$, %we have
        \[
            (\gamma_n \cdot y_n)(\iden) \,=\, y_n(\gamma_n) \,\neq\, y_n(\iden),
        \]
        and hence $\gamma_n \cdot y \neq y$. 
    \end{scproof}

    The next observation will be used in the proof of the continuous version of the LLL in \S\ref{sec:CLLL}:

    \begin{prop}\label{prop:order}
        If $X$ is a zero-dimensional Polish space, then there exists a linear order $\leq$ on $X$ such that the set $\set{(x,y) \in X^2 \,:\, x \leq y}$ is closed.
    \end{prop}
    \begin{scproof}
        Since every zero-dimensional Polish space can be embedded into the Cantor space $2^\N$ \cite[Thm.~7.8]{KechrisDST}, it is enough to consider $X = 2^\N$. On $2^\N$, the lexicographic order works. 
    \end{scproof}

    \section{Proof of Theorem~\ref{theo:generic}: $\mathcal{F}$-minimality is generic}\label{sec:generic}

    %Here we give the proof of Theorem~\ref{theo:generic}.
    Throughout this section, we work in a fixed Polish $\G$-flow $Y$. Recall that $\Sub(Y)$ is the space of all subflows of $Y$, equipped with the Vietoris topology (defined on p.~\pageref{page:Vietoris}).  %The proof is split into two parts: showing that the set of $S$-minimal subflows is $G_\delta$
    

    %\subsection{$S$-minimality is a $G_\delta$ property}

    %Recall that $\Sub(Y)$ is the set of all subflows of $Y$. The \emphd{Vietoris topology} on $\Sub(Y)$ has a basis consisting of all sets of the form
    %\[
    %    \llbracket U_0; U_1, \ldots, U_k\rrbracket \,\defeq\, \set{Z \in \Sub(Y) \,:\, Z \subseteq U_0, \, Z \cap U_1 \neq \0, \, \ldots, \, Z \cap U_k \neq \0},
    %\]
    %where $k \in \N$ and $U_0$, \ldots, $U_k$ are open subsets of $Y$. It is a standard fact that this topology makes $\Sub(Y)$ a compact Polish space \todo{cite}.
    
    Given $U \subseteq Y$, we define the following two subsets of $\Sub(Y)$:
    \begin{align*}
        \In(U) \,&\defeq\, \set{Z \in \Sub(Y) \,:\, Z \subseteq U},\\
        \Out(U) \,&\defeq\, \set{Z \in \Sub(Y) \,:\, U \cap Z = \0} \,=\, \In(Y \setminus U).
    \end{align*}

    \begin{lemma}\label{lemma:in-n-out}
        If $U \subseteq Y$ is an open set, then $\In(U)$ is open and $\Out(U)$ is closed in $\Sub(Y)$.
    \end{lemma}
    \begin{scproof}
        Indeed, $\In(U) = \llbracket U;\,\rrbracket$ and $\Out(U)$ is the complement of $\llbracket \Sub(Y); U\rrbracket$.
    \end{scproof}
    
    Now we fix an unbounded family $\mathcal{F}$ of finite subsets of $\G$. Given a set $U \subseteq Y$, % and a family $\mathcal{F}$ of finite subsets of $\G$,
    we let
    \[
        \Synd(\mathcal{F},U) \,\defeq\, \set{Z \in \Sub(Y) \,:\, \text{$U \cap Z$ is $F$-syndetic in $Z$ for some $F \in \mathcal{F}$}}.
    \]
    %\begin{align*}
    %    \Synd(\mathcal{F},U) \,&\defeq\, \set{Z \in \Sub(Y) \,:\, \text{$U \cap Z$ is $F$-syndetic in $Z$ for some $F \in \mathcal{F}$}},\\
    %    \Min(\mathcal{F}) \,&\defeq\, \set{Z \in \Sub(Y) \,:\, \text{$Z$ is $\mathcal{F}$-minimal}}.
    %\end{align*}
    %\begin{align*}
    %    \Sub^S(Y) \,&\defeq\, \set{Z \in \Sub(Y) \,:\, \text{$Z$ is $S$-minimal}},\\
    %    \Sub^{S,U}(Y) \,&\defeq\, \set{Z \in \Sub(Y) \,:\, \text{$U \cap Z$ is $S$-syndetic in $Z$}}.
    %\end{align*}

    \begin{lemma}\label{lemma:SU}
        For each open subset $U \subseteq Y$, % and every family $\mathcal{F}$ of finite subsets of $\G$, % $S \subseteq \G$ and every open $U \subseteq Y$,
        the set $\Synd(\mathcal{F},U)$ is open in $\Sub(Y)$.
    \end{lemma}
    \begin{scproof}
        Indeed, ${\Synd}(\mathcal{F},U)$ $= \bigcup_{F \in \mathcal{F}} \In(F^{-1} \cdot U)$, which is open by Lemma~\ref{lemma:in-n-out}. % is a union of open sets.
    \end{scproof}

    %Define $\Min(\mathcal{F}) \defeq \set{Z \in \Sub(Y) \,:\, \text{$Z$ is $\mathcal{F}$-minimal}}$.

    \begin{prop}\label{prop:Gdelta}
        %For each family $\mathcal{F}$ of finite subsets of $\G$,
        The set $\Min(\mathcal{F}) \defeq \set{Z \in \Sub(Y) \,:\, \text{$Z$ is $\mathcal{F}$-minimal}}$ is $G_\delta$ in $\Sub(Y)$.
    \end{prop}
    \begin{scproof}
        Let $(U_n)_{n \in \N}$ be a countable basis for the topology on $Y$. Then $Z \in \Sub(Y)$ is $\mathcal{F}$-minimal if and only if for all $n \in \N$, either $U_n \cap Z = \0$ or $U_n \cap Z$ is $F$-syndetic in $Z$ for some $F \in \mathcal{F}$. Hence,
        \begin{equation}\label{eq:min}
            \Min(\mathcal{F}) \,=\, \bigcap_{n \in \N} \big(\Out(U_n) \,\cup\, \Synd(\mathcal{F}, U_n)\big),
        \end{equation}
        which is $G_\delta$ by Lemmas~\ref{lemma:in-n-out} and \ref{lemma:SU}.
    \end{scproof}

    We also need the following observation:

    \begin{lemma}\label{lemma:closure}
        Let $V \subseteq Y$ and suppose that $A \subseteq Y$ is a $\G$-invariant set such that $V \cap A$ is $F$-syndetic in $A$ for some finite set $F \subset \G$. Then $\overline{V} \cap \overline{A}$ is $F$-syndetic in $\overline{A}$.
    \end{lemma}
    \begin{scproof}
        By assumption, %Let $F \subseteq S$ be a finite set such that
        $F^{-1} \cdot V \supseteq A$. Note that for each $\sigma \in F$,
        \[\sigma^{-1} \cdot \overline{V} \,=\, \overline {\sigma^{-1} \cdot V}.\]
        Since $F$ is finite, it follows that $F^{-1} \cdot \overline{V} = \overline{F^{-1} \cdot V} \supseteq \overline{A}$, as desired.
    \end{scproof}

    %\subsection{\Sr $\G$-spaces and density of $S$-minimal subflows}

    At this point, we fix an \sr $\G$-space $X$ and recall that
    \[
        \Sub_X(Y) \,\defeq\, \overline{\big\{\overline{\rho(X)} \,:\, \text{$\rho \colon X \to Y$ is a continuous $\G$-equivariant map}\big\}} \,\subseteq\, \Sub(Y).
    \]
    We are interested in the following subset of $\Sub_X(Y)$:
    \[
        \Min_X(\mathcal{F}) \,\defeq\, \Min(\mathcal{F}) \cap \Sub_X(Y) \,=\, \set{Z \in \Sub_X(Y) \,:\, \text{$Z$ is $\mathcal{F}$-minimal}}.
    \]
    Our goal is to show that $\Min_X(\mathcal{F})$ is a dense $G_\delta$ subset of $\Sub_X(Y)$. Since $\Min_X(\mathcal{F})$ is a $G_\delta$ set by Proposition~\ref{prop:Gdelta}, the only thing left to check is density.

    \begin{prop}\label{prop:dense}
        The set $\Min_X(\mathcal{F})$ is dense in $\Sub_X(Y)$.
    \end{prop}
    \begin{scproof}
        Thanks to formula \eqref{eq:min} and the Baire category theorem, it suffices to argue that the set
        \[
            \big(\Out(U) \,\cup\, \Synd(\mathcal{F}, U)\big) \,\cap\, \Sub_X(Y)
        \]
        is dense in $\Sub_X(Y)$ for every open $U \subseteq Y$. %, since countable intersections of dense $G_\delta$ subsets of $\Sub_X(Y)$ are themselves dense by the Baire category theorem.
        To this end, suppose %let $\llbracket U_0; U_1, \ldots, U_k \rrbracket$ be a basic open subset of $\Sub(Y)$ % $U_0$, \ldots, $U_k$ be open subsets of $Y$
        %satisfying
        \begin{equation}\label{eq:nonempty_basic}
            \llbracket U_0; U_1, \ldots, U_k \rrbracket \,\cap\, \Sub_X(Y) \,\neq \,\0
        \end{equation}
        for some basic open set $\llbracket U_0; U_1, \ldots, U_k \rrbracket$ in $\Sub(Y)$. We need to find a subflow
        \[%\begin{equation}\label{eq:witness}
            Z \,\in\, \llbracket U_0; U_1, \ldots, U_k \rrbracket \,\cap\, \big(\Out(U) \,\cup\, \Synd(\mathcal{F}, U)\big) \,\cap\, \Sub_X(Y).
        \]%\end{equation}
        By \eqref{eq:nonempty_basic} and the definition of $\Sub_X(Y)$, there is a continuous $\G$-equivariant map $\rho \colon X \to Y$ with
        \[
            \overline{\rho(X)} \,\in\, \llbracket U_0; U_1, \ldots, U_k \rrbracket.
        \]
        If $\overline{\rho(X)} \in \Out(U)$, then taking $Z \defeq \overline{\rho(X)}$ completes the proof. Otherwise, we have
        \[
            \overline{\rho(X)} \cap U_1 \neq \0, \quad \ldots, \quad \overline{\rho(X)} \cap U_k \neq \0, \quad \overline{\rho(X)} \cap U \neq \0.
        \]
        Since the sets $U_1$, \ldots, $U_k$, $U$ are open, this is equivalent to
        \[
            {\rho(X)} \cap U_1 \neq \0, \quad \ldots, \quad {\rho(X)} \cap U_k \neq \0, \quad {\rho(X)} \cap U \neq \0.
        \]
        We may pick open sets $V_1$, \ldots, $V_k$, $V$ such that $\overline{V_1} \subseteq U_1$, \ldots, $\overline{V_k} \subseteq U_k$, $\overline{V} \subseteq U$ and
        \[
            {\rho(X)} \cap V_1 \neq \0, \quad \ldots, \quad {\rho(X)} \cap V_k \neq \0, \quad {\rho(X)} \cap V \neq \0.
        \]
        Then $\rho^{-1}(V_1)$, \ldots, $\rho^{-1}(V_k)$, $\rho^{-1}(V)$ are nonempty open subsets of $X$. Since $X$ is \sr, for any sufficiently large set $F \in \mathcal{F}$ (which exists because $\mathcal{F}$ is unbounded), we have a
        %
        %Now we use the fact that $X$ is \sr and $$ to obtain a
        continuous $\G$-equivariant map $\pi \colon X \to X$ such that the sets
        \[
            (\rho \circ \pi)^{-1}(V_1), \quad \ldots, \quad (\rho \circ \pi)^{-1}(V_k), \quad (\rho \circ \pi)^{-1}(V)
        \]
        are $F$-syndetic in $X$. Equivalently, the sets
        \[
            V_1 \cap (\rho \circ \pi)(X), \quad \ldots, \quad V_k \cap (\rho \circ \pi)(X), \quad V \cap (\rho \circ \pi)(X)
        \]
        are $F$-syndetic in $(\rho \circ \pi)(X)$. %        . Let $\sigma \defeq \rho \circ \pi \colon X \to Y$.
        Now we claim that we can take %Now we claim that the subflow
        \[
            Z \,\defeq\, \overline{(\rho \circ \pi)(X)}. %\overline{\sigma(X)}.
        \]
        %satisfies \eqref{eq:witness}.
        Indeed, $\rho \circ \pi \colon X \to Y$ is a continuous $\G$-equivariant map, and hence $Z \in \Sub_X(Y)$. %by definition. %The set %$\sigma^{-1}(V)$ is $S$-syndetic in $X$, or, equivalently, %which is equivalent to the set
        %$V \cap (\rho \circ \pi)(X)$ is $S$-syndetic in $\sigma(X)$.
        Since $U \supseteq \overline{V}$, the set $U \cap Z$ is $F$-syndetic in $Z$ by Lemma~\ref{lemma:closure}, and thus $Z \in \Synd(\mathcal{F},U)$. The same reasoning shows that the sets $U_1 \cap Z$, \ldots, $U_k \cap Z$ are $F$-syndetic in $Z$ as well; in particular, they are nonempty.  Finally, $Z \subseteq \overline{\rho(X)} \subseteq U_0$ by construction, and therefore $Z \in \llbracket U_0; U_1, \ldots, U_k\rrbracket$, as desired.
    \end{scproof}

    \begin{scproof}[ of Theorem~\ref{theo:generic}]
        Follows by combining Propositions~\ref{prop:Gdelta} and \ref{prop:dense}.
    \end{scproof}

    \section{Proof of Lemma~\ref{lemma:from_separation_to_map}: Maps to $\Sep(s)$ with $\Phi$-spaced constraints}\label{sec:spaced}

    In this section we prove Lemma~\ref{lemma:from_separation_to_map}, restated here for ease of reference:

    \begin{lemmacopy}{lemma:from_separation_to_map}
       For every $s \in \N^+$ and a finite tuple $U_1$, \ldots, $U_k \subseteq ((s+1)^\G)^\N$ of nonempty open sets, there is a finite set $\Phi = \Phi(U_1, \ldots, U_k) \subset \G$ with the following property:

       \smallskip
       
       \noindent Suppose $X$ is a free zero-dimensional Polish $\G$-space such that $\casi(X) \leq s$, and let $A_1$, \ldots, $A_k \subseteq X$ be disjoint clopen sets whose union $A \defeq A_1 \sqcup \ldots \sqcup A_k$ is $\Phi$-spaced. Then there exists a continuous $\G$-equivariant map $\pi \colon X \to \Sep(s)$ such that $\pi(x) \in U_i$
        for all $1 \leq i \leq k$ and $x \in A_i$.
    \end{lemmacopy}
    \begin{scproof}
        We may assume that there exist a finite set $D \subset \G$ containing $\iden$, a natural number $N \in \N$, and mappings $\phi_{i,n} \colon D \to (s+1)$ for $1 \leq i \leq k$, $0 \leq n < N$ such that %\[\phi_1,\ \ldots,\ \phi_k \colon (Q \times N) \to (s+1),\] such that for all $1 \leq i \leq k$,
        $
            U_i = \prod_{n \in \N} U_{i,n} %\prod_{n = 0}^{N-1} U_{i,n} \times \prod_{n \geq N} (s+1)^\G,
        $,
        where %for all $1 \leq i \leq k$ and $n\in \N$,
        \[
            U_{i,n} \,\defeq\, \begin{cases}
                \big\{x \in (s+1)^{\G} \,:\, \rest{x}{D} = \phi_{i,n} \big\} &\text{if } n < N,\\
                (s+1)^\G &\text{if } n \geq N.
            \end{cases}
        \]
        No generality is lost because sets of this type form a basis for the topology on $((s+1)^{\G})^{\N}$.
        
        Recall that $\Sep(s) = \prod_{n \in \N} \Sep(s, \Phi_n)$, where $(\Phi_n)_{n \in \N}$ is a list of all finite subsets of $\G$ that includes each set infinitely often. We let $\Phi_n^* \defeq \Phi_n \cup \Phi_n^{-1} \cup \set{\iden}$ and define
        \[
             \Phi \,\defeq\, \bigcup_{n < N} \Phi_n^* DD^{-1}\Phi_n^*.
        \]
        We claim that this finite set $\Phi$ works.
        
        Let $X$, $A_1$, \ldots, $A_k$, $A$ be as in the statement of the lemma. We seek a continuous $\G$-equivariant map $\pi \colon X \to \Sep(s)$ such that $\pi(x) \in U_i$ for all $x \in A_i$. Since $U_i = \prod_{n \in \N} U_{i,n}$, it is enough to argue that for each $n \in \N$, there exists a continuous $\G$-equivariant map
            \[
                \pi_n \colon X \to \Sep(s, \Phi_n)
        \]
        such that $\pi_n(x) \in U_{i,n}$ for all $x \in A_i$, as then we can define the desired map $\pi \colon X \to \Sep(s)$ via
        \[  
            \pi(x) \,\defeq\, (\pi_n(x))_{n \in \N}.
        \]
        Thus, for the remainder of the proof, we fix some $n \in \N$.
        
        If $n \geq N$, then %$U_{i,n} = (s+1)^\G$, so
        we can let $\pi_n \colon X \to \Sep(s,\Phi_n)$ be an arbitrary continuous $\G$-equivariant map, which exists by Proposition~\ref{prop:map_to_Sep} as $\casi(X) \leq s$.
            
            Now suppose that $n < N$. %Note that since $D \supseteq Q$, for each element $x \in Q\cdot C$, there is a unique pair $(i,\delta) \in \set{1,\ldots, k} \times Q$, which we call the \emphd{type} of $x$, such that $\delta^{-1} \cdot x \in C_i$.
            Since $\casi(X) \leq s$, we may fix a continuous map $h \colon X \to (s+1)$ such that the sets $h^{-1}(0)$, \ldots, $h^{-1}(s)$ are $\Phi$-finite. % $\G$-equivariant map \[\rho \colon X \to \Sep(s,F)\] and let $h \colon X \to (s+1) \colon x \mapsto (\rho(x))(\iden)$ be the function such that $\rho = \pi_h$ (see \S\ref{subsec:shift}). %We modify $h$ by copying the mapping $\phi_{i,n} \colon Q \to (s+1)$ onto the set $Q \cdot y$ for each $y \in C_i$ and denote the resulting function by $f \colon X \to (s+1)$. That is, we let 
            %\[
            %    (\pi_n(x))(\gamma) \,\defeq\, \begin{cases}
            %        \phi_{i,n}(\delta) &\text{if $\gamma \cdot x \in Q\cdot C$ and $\mathsf{type}(\gamma \cdot x) = (i,\delta)$},\\
            %        (\rho(x))(\gamma) &\text{if $\gamma \cdot x \notin Q\cdot C$}.
            %    \end{cases}
            %\]
            %
            Define
            \begin{equation}\label{eq:modify}
                f(x) \,\defeq\, \begin{cases}
                    \phi_{i,n}(\delta) &\text{if $\delta \in D$ and $\delta^{-1}\cdot x \in A_i$},\\
                    h(x) &\text{if $x \notin D \cdot A$}.
                \end{cases}
            \end{equation}
            Note that since $\Phi \supseteq D$ and $A$ is $\Phi$-spaced, for each point $x \in D \cdot A$, there exists a unique pair $(\delta,i) \in D \times \set{1, \ldots, k}$ such that $\delta^{-1}\cdot x \in A_i$, and hence formula \eqref{eq:modify} describes a well-defined continuous function $f \colon X \to (s+1)$. Informally, to obtain the function $f$, we modify $h$ by copying the mapping $\phi_{i,n} \colon D \to (s+1)$ onto each set of the form $D \cdot x$ with $x \in A_i$.
            
            We claim that the coding map $\pi_n \defeq \pi_f \colon X \to (s+1)^\G$ corresponding to $f$ has the desired properties. The definition of $f$ ensures that for all $x \in A_i$, we have $\rest{\pi_n(x)}{D} = \phi_{i,n}$, i.e., $\pi_n(x) \in U_{i,n}$. It remains to verify that $\pi_n(x) \in \Sep(s, \Phi_n)$ for all $x \in X$. %That is, we need to show that for all $0 \leq t \leq s$, the set $f^{-1}(t)$ is $F_n$-finite.

            To this end, fix $0 \leq i \leq s$ and let $G_i$ be the subgraph of $\Sch(X, \Phi_n)$ induced by the set $f^{-1}(i)$. Similarly, let $H_i$ be the subgraph of $\Sch(X,\Phi)$ induced by $h^{-1}(i)$. %respectively $h^{-1}(t)$. %We know that every connected component of the graph $H_t$ is finite.
            We need to argue that all connected components of $G_i$ are finite. 
            Suppose, toward a contradiction, that $C$ is an infinite component of $G_i$. Since the set $A$ is $\Phi$-spaced and $\Phi \supseteq \Phi_n^* D$, there is no infinite path in $\Sch(X,\Phi_n)$ consisting solely of vertices in $D \cdot A$. Hence, the set $C \setminus (D \cdot A)$ must be infinite. This, however, is impossible, as any two vertices in $C \setminus (D \cdot A)$ belong to the same \ep{finite} component of $H_i$. Indeed, suppose %for contradiction
            that $(x_0,x_1,\ldots,x_\ell)$ is a shortest path in $C$ such that $x_0$, $x_\ell \in C \setminus (D \cdot A)$ belong to distinct components of $H_i$ (note that they must belong to $H_i$ because $h$ and $f$ may only differ on $D \cdot A$). As $\Phi \supseteq \Phi_n$, we have $\ell \geq 2$. By the minimality of the path, $x_1$, \ldots, $x_{\ell - 1} \in D \cdot A$, and since $A$ is $\Phi$-spaced, we conclude that $x_1$, \ldots, $x_{\ell - 1} \in D \cdot y$ for some $y \in A$. %Since $D \supseteq (F_n \cup F_n^{-1} \cup \set{\iden})$ and the set $A$ is $D$-separated, it follows that in fact $y_1 = \cdots = y_{\ell-1}$.
            But then
            \begin{align*}
                x_{\ell} \,\in\, \Phi_n^* \cdot x_{\ell - 1} \,\subseteq\, \Phi_n^*D \cdot y \,\subseteq\, \Phi_n^*DD^{-1} \cdot x_1 \,\subseteq\,  \Phi_n^*DD^{-1}\Phi_n^* \cdot x_0 \,\subseteq\, \Phi \cdot x_0,
            \end{align*}
            which implies that $x_\ell$ and $x_0$ are either equal or adjacent in $H_i$; a contradiction.
    \end{scproof}
    

    \section{Continuous Lov\'asz Local Lemma via continuous asymptotic separation index}\label{sec:CLLL}

    \subsection{Classical Lov\'asz Local Lemma}\label{subsec:LLL}

    Here we state the LLL in its classical combinatorial form, before moving on to the continuous setting in \S\ref{subsec:CLLL}. Given a family $\mathcal{B}$ of ``bad'' random events in a probability space, the LLL provides a sufficient condition that guarantees that all these ``bad'' events can be avoided. To formulate it, we need the following notion:

    \begin{defn}[Dependency relations]
        Let $\mathcal{B}$ be a family of random events in a probability space $(\Omega, \P)$. A \emphd{dependency relation} on $\mathcal{B}$ is a reflexive binary relation $\sim$ %on $\mathcal{B}$
        such that every event $B \in \mathcal{B}$ is mutually independent\footnote{Recall that random events $A$, $A'$ are \emphd{independent} if $\P[A \cap A'] = \P[A] \P[A']$, and an event $A$ is \emphd{mutually independent} from a family of events $\mathcal{A}$ if it is independent from every Boolean combination of the events in $\mathcal{A}$.} from the events $\set{B' \in \mathcal{B} \,:\, B' \not\sim B}$.
    \end{defn}

    \begin{exmp}\label{exmp:vbl}
        Let $(\Omega_i, \P_i)_{i \in I}$ be a family of probability spaces and let $(\Omega, \P) \defeq \prod_{i \in I} (\Omega_i, \P_i)$ be their product space. Suppose $\mathcal{B}$ is a family of random events in $\Omega$ such that the outcome of each event $B \in \mathcal{B}$ is determined by some nonempty subset $\dom(B) \subseteq I$ of coordinates. Then the relation $\sim$ such that $B \sim B'$ if and only if $\dom(B) \cap \dom(B' )\neq \0$ is a dependency relation on $\mathcal{B}$ \cite[41]{MolloyReed}. This is by far the most common setting in which the LLL is applied.
    \end{exmp}

    \begin{theo}[{Lov\'asz Local Lemma \cites{EL}{SpencerRamsey}[Corl.~5.1.2]{AS}}]\label{theo:LLL}
        Let $\mathcal{B}$ be a finite family of random events in a probability space $(\Omega, \P)$ and let $\sim$ be a dependency relation on $\mathcal{B}$. Suppose that there exist $p \in [0,1)$ and $d \in \N^+$ such that:
        \begin{itemize}
            \item $\P[B] \leq p$ and $|\set{B' \in \mathcal{B} \,:\, B' \sim B}| \leq d$ for all $B \in \mathcal{B}$,
            \item $\mathsf{e} p d \leq 1$, where $\mathsf{e} = 2.718\ldots$ is the base of the natural logarithm. 
        \end{itemize}
        Then $\P\left[\bigcup \mathcal{B}\right] < 1$ and, in particular, $\Omega \setminus \bigcup \mathcal{B} \neq \0$.
    \end{theo}

    In practice, the ``in particular'' part of Theorem~\ref{theo:LLL} is often valid for infinite families $\mathcal{B}$ as well. For example, we have the following:

    \begin{corl}[{Infinite LLL}]\label{corl:inf_LLL}
        Suppose $\Omega$ is a compact space endowed with a Borel probability measure $\P$. Let $\mathcal{B}$ be a family of open subsets of $\Omega$ and let $\sim$ be a dependency relation on $\mathcal{B}$. If there exist $p \in [0,1)$ and $d \in \N^+$ as in Theorem~\ref{theo:LLL},  %such that:
        %\begin{itemize}
        %    \item $\P[B] \leq p$ and $|\set{B' \in \mathcal{B} \,:\, B' \sim B}| \leq d$ for all $B \in \mathcal{B}$,
        %    \item $\mathsf{e} p d \leq 1$, where $\mathsf{e} = 2.718\ldots$ is the base of the natural logarithm. 
        %\end{itemize}
        then $\Omega \setminus \bigcup \mathcal{B} \neq \0$.
    \end{corl}
    \begin{scproof}
        By compactness, we just need to argue that $\Omega \setminus \bigcup \mathcal{B}' \neq \0$ for every finite subfamily $\mathcal{B}' \subseteq \mathcal{B}$, which is true by Theorem~\ref{theo:LLL}.
    \end{scproof}
    

    Let us now describe the framework in which the LLL shall be employed in this paper. Fix $k \in \N^+$ and let $W \subset \G$ be a finite set, called a \emphd{window}. Given a set $\mathcal{P} \subseteq k^W$ of \emphd{patterns}, we define
    \[
        \Sigma(k, W, \mathcal{P}) \,\defeq\, \big\{x \in k^\G \,:\, \rest{(\gamma \cdot x)}{W} \in \mathcal{P} \text{ for all } \gamma \in \G\big\},
    \]
    where, as usual, the vertical line indicates the restriction of a function to a subset of its domain. The set $\Sigma(k, W, \mathcal{P})$ is closed and shift-invariant, so when $\Sigma(k,W,\mathcal{P}) \neq \0$, it is a subshift. Subshifts of the form $\Sigma(k, W, \mathcal{P})$ are called \emphd{subshifts of finite type}; see \cite{subshifts1,subshifts2} for more background on their role in symbolic dynamics.
    
    A simple application of the LLL provides a condition that implies $\Sigma(k, W, \mathcal{P})$ is nonempty:

    \begin{prop}\label{prop:SFT}
       Let $k \in \N^+$, let $\0\neq W \subset \G$ be a finite set, and let $\mathcal{P} \subseteq k^W$. Suppose that
       \[
            \mathsf{e} \left(1 - \frac{|\mathcal{P}|}{k^{|W|}}\right) |W|^2 \,\leq\, 1.
       \]
       Then $\Sigma(k, W, \mathcal{P}) \neq \0$.
    \end{prop}
    \begin{scproof}
        %We may assume $W \neq \0$.
        We equip the set $k = \set{0,\ldots, k-1}$ with the uniform probability measure and work in the space $k^\G$ with the product measure $\P$. For $\gamma \in \G$, define $B_\gamma \subseteq k^\G$ by
        \[
            B_\gamma \,\defeq\, \big\{x \in k^\G \,:\, \rest{(\gamma \cdot x)}{W} \notin \mathcal{P}\big\},
        \]
        and let $\mathcal{B} \defeq \set{B_\gamma \,:\, \gamma \in \G}$. Note that each set $B_\gamma$ is clopen.  % and observe that, by definition,
        %\[
        %    \Sigma(k, W, \mathcal{P}) \,=\, k^\G \setminus \bigcup \mathcal{B}.
        %\]
        As in Example~\ref{exmp:vbl}, we observe that the membership of a point $x \in k^\G$ in $B_\gamma$ is determined by the restriction of $x$ to $W\gamma$, and hence we can define a dependency relation $\sim$ on $\mathcal{B}$ via %by making
        $B_\gamma \sim B_{\gamma'}$ $\Longleftrightarrow$ %if and only if
        $(W\gamma) \cap (W\gamma') \neq \0$. Each event $B_\gamma$ satisfies
        \[
            \P[B_\gamma] \,=\, 1 - \frac{|\mathcal{P}|}{k^{|W|}} \,\eqqcolon\, p \qquad \text{and} \qquad |\set{B_{\gamma'} \in \mathcal{B} \,:\, B_\gamma \sim B_{\gamma'}}| \,=\, |W^{-1}W\gamma| \,\leq\, |W|^2 \,\eqqcolon\, d,
        \]
        so we may apply Corollary~\ref{corl:inf_LLL} to conclude that $\Sigma(k, W, \mathcal{P}) = k^\G \setminus \bigcup \mathcal{B} \neq \0$, as desired.
    \end{scproof}

    \subsection{A continuous version}\label{subsec:CLLL}

    We can now formulate the continuous variant of the LLL that we use to prove Lemma~\ref{lemma:sparse_syndetic}:

    \begin{tcolorbox}
    \begin{theo}\label{theo:cont_LLL}
        Let $s$, $k \in \N^+$, let $\0\neq W \subset \G$ be a finite set, and let $\mathcal{P} \subseteq k^W$. Suppose that
       \begin{equation}\label{eq:cont_LLL}
            \mathsf{e}^{s+1} \left(1 - \frac{|\mathcal{P}|}{k^{|W|}}\right) |W|^{2(s+1)} \,\leq\, 1.
       \end{equation}
       Then every free zero-dimensional Polish $\G$-space $X$ % continuous action $\G \acts X$ on a zero-dimensional Polish space
       such that $\casi(X) \leq s$ admits a continuous $\G$-equivariant map $\pi \colon X \to \Sigma(k, W, \mathcal{P})$.
    \end{theo}
    \end{tcolorbox}

    A straightforward application of a result of Weilacher and the first named author \cite[Thm.~1.29]{BW} %can be used to obtain
    yields a {Borel} $\G$-equivariant map $\pi \colon X \to \Sigma(k,W,\mathcal{P})$ assuming $\asi(X) \leq s$. To prove Theorem~\ref{theo:cont_LLL}, we observe that when $\casi(X) \leq s$, a careful implementation of the construction from \cite{BW} %produces a continuous map $\pi$.
        makes the map $\pi$ continuous.
        %For completeness
        To make the presentation self-contained, we describe the construction here. The main tool we use is the \emph{method of conditional probabilities}, which is a common derandomization technique in computer science % from computer science, where it is commonly used for derandomization
        \cites[\S16]{AS}[\S5.6]{RandAlg}. An analogous approach was employed in \cite{Ber_cont} to establish a different continuous version of the LLL, and similar arguments have appeared in other related contexts as well, for example in \cite[Thm.~3.6]{FG} by Fischer and Ghaffari.
    
    \begin{scproof}
        

        %Let $s$, $k$, $W$, $\mathcal{P}$, and $X$ be as in the statement of Theorem~\ref{theo:cont_LLL}.
        
        %We may assume $W \neq \0$.
        Thanks to the coding map correspondence $f \longleftrightarrow \pi_f$, our goal is to find a continuous function $f \colon X \to k$ such that for all $x \in X$, the map $W \to k$ sending each $\sigma \in W$ to $f(\sigma \cdot x)$ belongs to $\mathcal{P}$. To this end, we introduce a family of ``bad'' random events analogous to the one in the proof of Proposition~\ref{prop:SFT}. Let $k = \set{0,\ldots, k-1}$ carry the uniform probability measure and consider the %(in general, non-separable)
        space $k^X$ with the product measure $\P$. For each $x \in X$, let %define $B_x \subseteq k^X$ by
        \[
            B_x \,\defeq\, \big\{f \in k^X \,:\, \text{the map $W \to k \colon \sigma \mapsto f(\sigma \cdot x)$ is not in $\mathcal{P}$}\big\},
        \]
        and set $\mathcal{B} \defeq \set{B_x \,:\, x \in X}$. Each set $B_x$ is clopen in $k^X$. The relation $\sim$ on $\mathcal{B}$ given by $B_x \sim B_y$ $\Longleftrightarrow$ $(W \cdot x) \cap (W \cdot y) \neq \0$ is a dependency relation, and every event $B_x$ satisfies
        \[
            \P[B_x] \,=\, 1 - \frac{|\mathcal{P}|}{k^{|W|}} \,\eqqcolon\, p \qquad \text{and} \qquad |\set{B_{y} \in \mathcal{B} \,:\, B_x \sim B_{y}}| \,=\, |W^{-1}W \cdot x| \,\leq\, |W|^2 \,\eqqcolon\, d.
        \]
        At this point we may apply Corollary~\ref{corl:inf_LLL} to obtain a function $f \in k^X \setminus \bigcup \mathcal{B}$. However, this is not enough for our purposes, as we also need $f$ to be {continuous}.

        To achieve this aim, we will use the assumption that $\casi(X) \leq s$. Define $\Phi \defeq WW^{-1}$, and fix a partition $X = X_0 \sqcup \ldots \sqcup X_s$ into $\Phi$-finite clopen sets. Let $G_i$ be the subgraph of the Schreier graph $\Sch(X,\Phi)$ induced by $X_i$ and let $\mathcal{C}_i$ be the set of all connected components of $G_i$. Note that all components of $G_i$ are finite. The set $\Phi$ is chosen to guarantee that each set of the form $W \cdot x$ meets at most one component of $G_i$.
        
        %We also fix a linear order $\leq$ on $X$ such that the set $\set{(x,y) \in X^2 \,:\, x \leq y}$ is closed, given by Proposition~\ref{prop:order}. %Thanks to the coding map correspondence $f \longleftrightarrow \pi_f$ (see \S\todo{ref}), our goal is to find a continuous function $f \colon X \to k$ such that for all $x \in X$, the mapping $W \to k \colon \gamma \mapsto f(\gamma \cdot x)$ belongs to $\mathcal{P}$.  
        %For $0 \leq i \leq s$, let $\mathcal{C}_i$ be the set of all connected components of the subgraph of $\Sch(X, F)$ induced by $X_i$, and let $\mathcal{C} \defeq \bigcup_{i=0}^s \mathcal{C}_i$. We call the members of $\mathcal{C}$ \emphd{clusters}. By the choice of $X_0$, \ldots, $X_s$, all clusters are finite. %The \emphd{cluster graph} $G$ has vertex set $\mathcal{C}$ and includes an edge between two distinct clusters $C$, $C'$ if and only if $\Sch(X, F)$ has an edge joining a vertex in $C$ to a vertex in $C'$.
        %For $x \in X$, we let $[x] \in \mathcal{C}$ be the unique cluster containing $x$.
        %
        Given a partial map $h \colon X \pto k$ and a point $x \in X$, we write $\P[B_x \,\vert\, h]$ for the conditional probability that a random function $f \in k^X$ is in $B_x$
        given that $f$ agrees with $h$ on $(W \cdot x) \cap \dom(h)$. Note that when $\dom(h) = X$, we have $\P[B_x \,\vert\, h] = 1$ if $h \in B_x$ and $\P[B_x \,\vert\, h] = 0$ otherwise. 
        %Note that the value $\P[B_x \,\vert\, h]$ is determined by the restriction of $h$ to the finite set $(W \cdot x) \cap \dom(h)$.
        Inducting on $0 \leq i \leq s$, we shall define continuous functions $f_i \colon X_i \to k$ such that for all $x \in X$,
        \[
            \P[B_x \,\vert\, f_0 \cup \ldots \cup f_i]  \,<\, (\mathsf{e} d)^{-(s - i)}.
        \]
        Then the map $f \defeq f_0 \cup \ldots \cup f_s \colon X \to k$ satisfies $\P[B_x \,\vert\, f] < 1$ for all $x \in X$. As $\dom(f) = X$, we conclude that $\P[B_x \,\vert\, f] = 0$ for all $x \in X$, %so $f$ is as desired.
        i.e., $f \notin \bigcup \mathcal{B}$, as desired.

        To construct the functions $f_0$, \ldots, $f_s$, we first observe that, by \eqref{eq:cont_LLL}, 
        \[
            \P[B_x] \,=\, p \,<\, (\mathsf{e} d)^{-(s+1)} \quad \text{for all } x \in B_x.
        \]
        (The inequality is strict because $p$ is rational, while $(\mathsf{e}d)^{s+1}$ is not.) Now suppose we already have continuous functions $f_0 \colon X_0 \to k$, \ldots, $f_{i-1} \colon X_{i-1} \to k$ for some $0 \leq i \leq s$ such that for all $x \in X$,
        \begin{equation}\label{eq:ind}
            \P[B_x \,\vert\, f_0 \cup \ldots \cup f_{i-1}] \,<\, (\mathsf{e}d)^{-(s-i+1)}.
        \end{equation}
        %Working in the space $k^{X_i}$ with the product measure $\P_i$, we define random events
        %\[
        %    B_{x,i} \,\defeq\, \big\{h \in k^{X_i} \,:\, \P[B_x \,\vert\, f_0 \cup \ldots \cup f_{i-1} \cup h] \geq (\mathsf{e}d)^{-(s-i)}\big\},
        %\]
        %and set $\mathcal{B}_i \defeq \set{B_{x,i} \,:\, x \in X}$. To complete the proof of Theorem~\ref{theo:cont_LLL}, we need to argue that there exists a continuous function $f_i \in k^{X_i} \setminus \bigcup \mathcal{B}_i$.
        %Let $G_i$ be the subgraph of the Schreier graph $\Sch(X,F)$ induced by $X_i$ and let $\mathcal{C}_i$ be the set of all connected components of $G_i$. Recall that, by the choice of $X_i$, all components of $G_i$ are finite. Note that each set of the form $W \cdot x$ meets at most one component of $G_i$. 
        To define $f_i$, we consider the components of the graph $G_i$ individually. For $C \in \mathcal{C}_i$, we say that a function $h \colon C \to k$ is \emphd{good} if% for $C$ if  %Note that each set of the form $W \cdot x$ meets at most one component of $G_i$.
        \[
                \P[B_x \,\vert\, f_0 \cup \ldots \cup f_{i-1} \cup h] \,<\, (\mathsf{e}d)^{-(s-i)} \quad \text{ for all } x \in X.
            \]
        
        \begin{claim}
            For each component $C \in \mathcal{C}_i$ of $G_i$, there exists a good function $h \colon C \to k$. % that is good for $C$. %such that
            %\[
            %    \P[B_x \,\vert\, f_0 \cup \ldots \cup f_{i-1} \cup h] < (\mathsf{e}d)^{-(s-i)} \quad \text{ for all } x \in X.
            %\]
        \end{claim}
        \begin{claimproof}
            Working in the space $k^C$ with the product measure $\P_C$, we consider the random events
            \[
                B_{x,C} \,\defeq\, \big\{h \in k^C \,:\, \P[B_x \,\vert\, f_0 \cup \ldots \cup f_{i-1} \cup h] \geq (\mathsf{e}d)^{-(s-i)} \big\},
            \]
            and set $\mathcal{B}_{C} \defeq \set{B_{x,C} \,:\, x \in X}$. Since each event $B_{x,C}$ is determined by the restriction of $h$ to the set $(W \cdot x) \cap C \subseteq W \cdot x$, %we can define a %slightly abuse notation and employ the symbol $\sim$ to denote the
            %dependency relation on $\mathcal{B}_{x,C}$ by
            relating $B_{x,C}$ and $B_{y,C}$ if and only if $(W \cdot x) \cap (W \cdot y)$ gives a dependency relation on $\mathcal{B}_{x,C}$. 
            Using Markov's inequality and \eqref{eq:ind}, we see that for all $x \in X$,
            \[
                \P_C[B_{x,C}] \,\leq\, \frac{\P[B_x \,\vert\, f_0 \cup \ldots \cup f_{i-1}]}{(\mathsf{e}d)^{-(s-i)}} \,<\, \frac{(\mathsf{e}d)^{-(s-i+1)}}{(\mathsf{e}d)^{-(s-i)}} \,=\, \frac{1}{\mathsf{e}d}.
            \]
            It follows that $\mathcal{B}_{C}$ satisfies the assumptions of the LLL, and hence a desired function $h \in k^C \setminus \bigcup \mathcal{B}_C$ exists by Corollary~\ref{corl:inf_LLL}.
        \end{claimproof}

        Fix a closed linear order $\leq$ on $X$ given by Proposition~\ref{prop:order}. Take any $C \in \mathcal{C}_i$ and list the elements of $C$ in the increasing order as $x_1 < \cdots < x_{|C|}$. A function $h \colon C \to k$ can then be represented by the tuple $(h(x_1), \ldots, h(x_{|C|}))$, and we let $h_C \colon C \to k$ be the good function %whose corresponding tuple  
        for which this tuple is lexicographically minimal. % represented by the lexicographically least such tuple. 
        Define
        $
            f_i \defeq \bigcup_{C \in \mathcal{C}_i} h_C
        $.
        We claim that $f_i \colon X_i \to k$ is as desired.
        
        First, we observe that for any $x \in X$,
        \[
            \P[B_x \,\vert\, f_0 \cup \ldots \cup f_{i-1} \cup f_i] \,<\, (\mathsf{e}d)^{-(s-i)}.
        \]
        Indeed, the value $\P[B_x \,\vert\, f_0 \cup \ldots \cup f_{i-1} \cup f_i]$ is determined by the restriction of $f_i$ to $(W \cdot x) \cap X_i$. If $(W \cdot x) \cap X_i = \0$, then the desired bound holds by \eqref{eq:ind}. 
        %\[
        %    \P[B_x \,\vert\, f_0 \cup \ldots \cup f_{i-1} \cup f_i] \,=\, \P[B_x \,\vert\, f_0 \cup \ldots \cup f_{i-1}] \,<\, (\mathsf{e}d)^{-(s-i+1)} \,<\, (\mathsf{e}d)^{-(s-i)}
        %\]
        %by \eqref{eq:ind}.
        On the other hand, if $(W \cdot x) \cap X_i \neq \0$, then there is a unique component $C \in \mathcal{C}_i$ such that $(W \cdot x) \cap X_i = (W \cdot x) \cap C$, in which case
        \[
            \P[B_x \,\vert\, f_0 \cup \ldots \cup f_{i-1} \cup f_i] \,=\, \P[B_x \,\vert\, f_0 \cup \ldots \cup f_{i-1} \cup h_C] \,<\, (\mathsf{e}d)^{-(s-i)},
        \]
        because $h_C$ is good.

        It remains to verify that $f_i$ is continuous. To this end, consider any point $x \in X_i$ and let $R_x$ be the (finite) set of all group elements $\gamma \in \G$ such that $\gamma \cdot x$ is in the same component of $G_i$ as $x$. It is evident from the definition of $f_i$ that the value $f_i(x)$ is determined by the following information:
        \begin{itemize}
            \item the set $R_x$,
            \item the restrictions of the functions $f_0$, \ldots, $f_{i-1}$ to $\Phi R_x \cdot x$,
            \item the restriction of the order $\leq$ to the set $R_x \cdot x$.
        \end{itemize}
        Since the sets $X_0$, \ldots, $X_s$ are clopen, the functions $f_0$, \ldots, $f_{i-1}$ and the action %are continuous, the action
        $\G \acts X$ are continuous, %is continuous and free,
        and the order $\leq$ is closed, this shows that $f_i(x) = f_i(y)$ for all $y$ in some open neighborhood of $x$, and hence $f_i$ is continuous, as desired.
     \end{scproof}


    \subsection{Application: Proof of Lemma~\ref{lemma:sparse_syndetic}}\label{subsec:synd_and_spaced}

    We are now ready to prove Lemma~\ref{lemma:sparse_syndetic}, restated here for ease of reference:

    \begin{lemmacopy}{lemma:sparse_syndetic}
        For any $s$, $k \in \N^+$ and a finite set $\Phi \subset \G$, there is $n = n(s,k, \Phi) \in \N$ such that for all finite $F \subset \G$ with $|F| \geq n$, the following holds:

        \smallskip
        
        \noindent For every free zero-dimensional Polish $\G$-space $X$ such that $\casi(X) \leq s$, there exist disjoint $F$-syndetic clopen subsets $A_1$, \ldots, $A_k \subseteq X$ whose union %such that 
        %\begin{enumerate}[label=\ep{\normalfont\roman*}]
          %  \item
         %each set $A_i$ is $F$-syndetic and
          %  \item the set
          $A \defeq A_1 \sqcup \ldots \sqcup A_k$ is $\Phi$-spaced.
        %\end{enumerate}
    \end{lemmacopy}
    \begin{scproof}
        Let $D \defeq \Phi^{-1}\Phi$ and $d \defeq |D|$. Fix an integer $r \in \N^+$ so large that
        \begin{equation}\label{eq:large_m}
            \mathsf{e}^{s+1} \, k \left(1 - \frac{1}{(k+1)^d}\right)^r \, r^{2(s+1)} \,\leq\, 1.
        \end{equation}
        This is possible since the left-hand side of \eqref{eq:large_m} tends to $0$ as $r \to \infty$. Now set $n \defeq d^2 r$.

        Suppose $F \subset \G$ is a finite set with $|F| \geq n$. Then $F$ has a $D$-spaced subset $R \subseteq F$ of size $|R| = r$ \cite[Lem.~4.1]{LLLDyn3}. 
        %
        %It is easy to see that any infinite set $S \subseteq \G$ has an infinite $D$-spaced subset $S' \subseteq S$. Hence, % Since $S$ is infinite,
        %we may pick a $D$-spaced subset $R \subset S$ of size exactly $|R| = r$.
        Let $W \defeq DR$. Call a mapping $\phi \colon W \to (k+1)$ \emphd{acceptable} if for each $1 \leq i \leq k$, there is $\sigma \in R$ such that
        \begin{equation}\label{eq:acceptable}
            \phi(\delta \sigma) \,=\, \begin{cases}
                i &\text{if } \delta = \iden,\\
                0 &\text{if } \delta \neq \iden,
            \end{cases} \qquad \text{for all }\delta \in D.
        \end{equation}
        %\begin{itemize}
        %    \item $p(\sigma) = i$, and
        %    \item $p(\delta\sigma) = 0$ for all $\delta \in D \setminus \set{\iden}$.
        %\end{itemize}
        Let $\mathcal{P} \subseteq k^W$ be the set of all acceptable maps $\phi \colon W \to k$.

        \begin{claim}\label{claim:bound_on_prob}
            $\displaystyle 1 - \frac{|\mathcal{P}|}{(k+1)^{|W|}} \,\leq\, k \left(1 - \frac{1}{(k+1)^d}\right)^r$.
        \end{claim}
        \begin{claimproof}
            We seek an upper bound on the probability that a uniformly random map $\phi \colon W \to (k+1)$ is not acceptable. %To this end, say that $p \colon W \to (k+1)$ is \emphd{$i$-acceptable} if there is $\sigma \in R$ satisfying \eqref{eq:acceptable}.
            The probability \eqref{eq:acceptable} holds for fixed $1 \leq i \leq k$ and a specific element $\sigma \in R$ is precisely $1/(k+1)^d$. Since the set $R$ is $D$-spaced, the conditions for different choices of $\sigma \in R$ are independent, so the probability \eqref{eq:acceptable} fails for some $1 \leq i \leq k$ and all $\sigma \in R$ is
            \[
                \left(1 - \frac{1}{(k+1)^d}\right)^r.
            \]
            Summing over all $1 \leq i \leq k$ yields the desired bound.
        \end{claimproof}

        Thanks to Claim~\ref{claim:bound_on_prob} and \eqref{eq:large_m}, we may apply Theorem~\ref{theo:cont_LLL} to obtain a continuous $\G$-equivariant map $\pi \colon X \to \Sigma(k+1, W, \mathcal{P})$. To complete the proof, we let $A_i$ be the set of all $x \in X$ such that
        \[
            (\pi(x))(\delta) \,=\, \begin{cases}
                i &\text{if } \delta = \iden,\\
                0 &\text{if } \delta \neq \iden,
            \end{cases} \qquad \text{for all }\delta \in D.
        \]
        The sets $A_i$ for $1 \leq i \leq k$ are clearly clopen and pairwise disjoint. By the definition of $\mathcal{P}$, we have $R^{-1} \cdot A_i = X$ for all $1 \leq i \leq k$, so $A_i$ is $F$-syndetic as $F \supseteq R$. Finally, if $x$, $y \in A \defeq A_1 \sqcup \ldots \sqcup A_k$ are distinct points, then $(\pi(y))(\iden) \neq 0$, and thus $y \notin (D \cdot x)$, i.e., $(\Phi \cdot x) \cap (\Phi \cdot y) = \0$, as desired.
    \end{scproof}

    \section{Finishing the proof of Theorem~\ref{theo:example}}\label{sec:final}

    Here we complete the proof that for $s \in \N^+$, $\Free(\Sep(s))$ is an \sr $\G$-space. Fix nonempty open sets $U_1$, \ldots, $U_k \subseteq ((s+1)^\G)^\N$. % and let $n = n(U_1, \ldots, U_k) \in \N$ be  given by Lemma~\ref{lemma:key}. 
    %
    %
    %Take an infinite subset $S \subseteq \G$ and nonempty open sets
    %        \[
    %        U_1, \ \ldots, \ U_k \ \subseteq \  ((s+1)^\G)^\N.
    %    \]
        As explained in \S\ref{subsubsec:combi+LLL}, Lemma~\ref{lemma:key} implies that for every large enough finite set $F \subset \G$, % of size at least $n$,
        we have %Lemma~\ref{lemma:key} %applied with $X = \Free(\Sep(s))$
        %yields
        a continuous $\G$-equivariant map 
        \[
            \pi \colon \Free(\Sep(s)) \to \Sep(s)
        \]
        such that the sets $\pi^{-1}(U_1)$, \ldots, $\pi^{-1}(U_k)$ are $F$-syndetic. It remains to modify the map $\pi$ to make its range included in $\Free(\Sep(s))$. (This step is only needed if $\G$ has nontrivial torsion elements.) 
        
        Recall that $(\Phi_n)_{n \in \N}$ is a list of all finite subsets of $\G$ that includes each set infinitely often, and
        \[
            \Sep(s) \,=\, \prod_{n \in \N} \Sep(s, \Phi_n).
        \]
        We may assume without loss of generality that there is $N \in \N$ such that each set $U_i$ is of the form
        \[
            U_i \,=\, V_i \times \prod_{n \geq N} (s+1)^\G,
        \]
        for some nonempty open $V_i \subseteq ((s+1)^\G)^N$. By composing $\pi$ with the coordinate projection
        \[
            \Sep(s) \,=\, \prod_{n \in \N} \Sep(s,\Phi_n) \,\twoheadrightarrow\, \prod_{n < N} \Sep(s,\Phi_n),
        \]
        %on the first $N$ coordinates produces
        we obtain a continuous $\G$-equivariant map
        \[
            \rho \colon \Free(\Sep(s)) \to \prod_{n < N}  \Sep(s,\Phi_n)
        \]
        such that the sets $\rho^{-1}(V_1)$, \ldots, $\rho^{-1}(V_k)$ are $F$-syndetic. Next we observe that, since every finite subset of $\G$ appears in the sequence $(\Phi_n)_{n \in \N}$ infinitely often, there is a $\G$-equivariant homeomorphism %we have %we can reorder the terms in the product to see that
        \[
            \sigma \colon \Sep(s) \to \prod_{n \geq N} \Sep(s,\Phi_n).
        \]
        The restriction of $\sigma$ to $\Free(\Sep(s))$ establishes a $\G$-equivariant homeomorphism
        \[
            \Free(\Sep(s)) \,\cong\, \Free\left(\prod_{n \geq N} \Sep(s,\Phi_n)\right).
        \]
        %In particular, there is a continuous $\G$-equivariant map
        %\[
        %    \rho \colon \Free(\Sep(s)) \to \prod_{n \geq N} \Sep(s,F_n)
        %\]
        %that maps $X$ to the free part of $\prod_{n \geq N} \Sep(d,F_n)$.
        Now we put the maps $\rho$ and $\sigma$ together and define
        \[
            \pi' \colon \Free(\Sep(s)) \to \prod_{n < N} \Sep(s,\Phi_n) \times \Free\left(\prod_{n \geq N} \Sep(s,\Phi_n)\right)
        \]
        via $\pi'(x) \defeq (\rho(x), \sigma(x))$. It remains to note that
        \[
            \prod_{n < N} \Sep(s,\Phi_n) \times \Free\left(\prod_{n \geq N} \Sep(s,\Phi_n)\right) \,\subseteq\, \Free\left(\prod_{n \in \N} \Sep(s, \Phi_n)\right) \,=\, \Free(\Sep(s)),
        \]
        so $\pi'$ has all the desired properties.

    \section{Subshifts with minimal subdynamics}%Proof of Corollary~\ref{corl:subshift}}
    \label{sec:subshift}

    %\subsection{Proof of Corollary~\ref{corl:subshift_prime}}

    In this section we prove Corollary~\ref{corl:subshift}, i.e., we show that the $\G$-flow in Theorem~\ref{theo:main_prime} can be taken to be a subshift. %We begin with a lemma:
    We shall use the following result of Seward and Tucker-Drob:

    \begin{theo}[{Seward--Tucker-Drob \cite{ST-D}}]\label{theo:ST-D}
        There exists a free subshift $Y \subset 2^\G$ such that every free Borel action $\G \acts X$ on a Polish space $X$ admits a Borel $\G$-equivariant map $\pi \colon X \to Y$.
    \end{theo}

    See also \cite{Ber_cont} for a shorter probabilistic proof of Theorem~\ref{theo:ST-D}. In addition to Theorem~\ref{theo:ST-D}, we need the following lemma:

    \begin{lemma}\label{lemma:transfer}
        Let $(\mathcal{F}_i)_{i \in I}$ be a collection of families of finite subsets of $\G$ and let % be a family of subsets of $\G$ and let
        $X$ be a %Polish
        $\G$-flow that is $\mathcal{F}_i$-minimal for all $i \in I$. If $Y$ is a $\G$-flow such that there exists a Baire-measurable $\G$-equivariant map $\pi \colon X \to Y$, then $Y$ has a subflow $Z \subseteq Y$ that is $\mathcal{F}_i$-minimal for all $i \in I$.
    \end{lemma}
    \begin{scproof}
        Call a closed set $C \subseteq Y$ \emphd{huge} if $\pi^{-1}(C)$ is comeager in $X$. Let the set of all huge closed subsets of $Y$ be denoted by $\mathcal{H}$. Note that if $C \in \mathcal{H}$, then $\gamma \cdot C \in \mathcal{H}$ as well for all $\gamma \in \G$, because
        $
            \pi^{-1}(\gamma \cdot C) = \gamma \cdot \pi^{-1}(C)
        $. Hence, the following is a closed $\G$-invariant subset of $Y$: % Let
        \[
            Z \,\defeq\, \bigcap_{C \in \mathcal{H}} C.
        \]
        %By construction, the set $Y'$ is closed. Furthermore, it is $\G$-invariant, because Note that if $C \in \mathcal{H}$, then $\gamma \cdot C \in \mathcal{H}$ as well for all $\gamma \in \G$. %because
        %\[
        %    \pi^{-1}(\gamma \cdot C) \,=\, \gamma \cdot \pi^{-1}(C)
        %\]
        %and the map $X \mapsto X \colon x \mapsto \gamma \cdot x$ is a homeomorphism.
        %Moreover, $Z$ is nonempty, i.e., it is a subflow of $Y$. Indeed,
        The intersection of finitely (or even countably) many huge sets is still huge, and hence it is nonempty by the Baire category theorem applied in $X$. Since $Y$ is compact, this implies that $Z \neq \0$, so $Z$ is a subflow of $Y$. (Moreover, if $Y$ is Polish, then $Z$ itself is huge.) 
        %Note that, since $Y$ is second-countable, there is a countable family $(C_n)_{n \in \N}$ of huge closed sets such that $Z = \bigcap_{n \in \N} C_n$. It follows that $\pi^{-1}(Z) = \bigcap_{n \in \N} \pi^{-1}(C_n)$ is comeager in $X$, so $Z$ itself is huge. In particular, $Z$ is nonempty, i.e., it is a subflow of $Y$. %, so $Z$ cannot be empty.
        
        %compact, it is enough to argue that $C_1 \cap \ldots \cap C_k \neq \0$ for every finite tuple of sets $C_1$, \ldots, $C_k \in \mathcal{H}$, which holds since the set
        %\[
        %    \pi^{-1}(C_1 \cap \ldots \cap C_k) \,=\, \pi^{-1}(C_1) \cap \ldots \cap \pi^{-1}(C_k)
        %\]
        %is comeager in $X$, hence nonempty.
        
        We claim that $Z$ is $\mathcal{F}_i$-minimal for all $i \in I$. Indeed, take any $i \in I$ and suppose $U \subseteq Y$ is an open set such that $U \cap Z \neq \0$. We need to find a set $F \in \mathcal{F}_i$ with $F^{-1} \cdot U \supseteq Z$. %show that $U \cap Z$ is $\mathcal{F}_i$-syndetic in $Z$.
        Since $Y$ is compact and Hausdorff, we may pick an open set $W \subseteq U$ such that $W \cap Z \neq \0$ and $C \defeq \overline{W} \subseteq U$. 
        The condition $W \cap Z \neq \0$ means that
        the closed set $Y \setminus W$ is not huge, i.e., $\pi^{-1}(W)$ is non-meager in $X$. As $C \supseteq W$, $\pi^{-1}(C)$ is non-meager as well. %Since $Y$ is Polish, every open set in $Y$ is $F_\sigma$, from which it follows that there is a closed subset $C \subseteq U$ such that $\pi^{-1}(C)$ is also non-meager in $X$.
        The set $\pi^{-1}(C)$ is Baire-measurable, so there is a nonempty open set $V \subseteq X$ such that the symmetric difference $V \symdif \pi^{-1}(C)$ is meager.  As $X$ is $\mathcal{F}_i$-minimal, there is $F \in \mathcal{F}_i$ such that $F^{-1}\cdot V = X$, and hence $F^{-1} \cdot \pi^{-1}(C) = \pi^{-1}(F^{-1} \cdot C)$ is comeager in $X$. In other words, the closed set $F^{-1} \cdot C$ is huge, and therefore $Z \subseteq F^{-1} \cdot C \subseteq F^{-1} \cdot U$, as desired.
    \end{scproof}

    \begin{scproof}[ of Corollary~\ref{corl:subshift}]
        Let $(\mathcal{F}_n)_{n \in \N}$ be a sequence of unbounded families of finite subsets of $\G$. By Theorem~\ref{theo:main_prime}, there exists a free $\G$-flow $X$ that is $\mathcal{F}_n$-minimal for all $n \in \N$. Furthermore, the proof of Theorem~\ref{theo:main_prime} presented in \S\ref{subsubsec:ample} yields such a $\G$-flow $X$ whose underlying space is Polish. By Theorem~\ref{theo:ST-D}, there exist a free subshift $Y \subset 2^\G$ and a Borel $\G$-equivariant map $\pi \colon X \to Y$. Since $\pi$ is Borel, it is in particular Baire-measurable, so, by Lemma~\ref{lemma:transfer}, $Y$ has a subflow $Z \subseteq Y$
        that is $\mathcal{F}_n$-minimal for all $n \in \N$, as desired.
    \end{scproof}

    \section{Disjoint flows}\label{sec:disjoint}

    %Let $X$ be a $\G$-flow. For a subset $U \subseteq X$, its \emphd{return set} is
    %\[
    %    \Ret(U) \,\defeq\, \set{\gamma \in \G \,:\, U \cap (\gamma \cdot U) \neq \0}.
    %\]
    %In particular, $U$ is wandering if and only if the return set $\Ret(U)$ is finite. Let us record the following two well-known simple observations:

    %\begin{lemma}\mbox{}
    %
    %    \begin{enumerate}[label=\ep{\normalfont\roman*}]
    %    \item\label{item:min} If $X$ and $Y$ are disjoint $\G$-flows, then at least one of them is minimal.
    %        \item\label{item:no_wandering} A minimal $\G$-flow $X$ has no wandering points.
    %    \end{enumerate}
    %\end{lemma}
    %\begin{scproof}
    %    \ref{item:min} If $X' \subset X$, $Y' \subset Y$ are proper subflows, then $Z \defeq (X \times Y') \cup (X' \times Y)$ is a joining distinct from $X \times Y$.

    %    \ref{item:no_wandering} Suppose $U \subseteq X$ is a wandering nonempty open set. Since $X$ is minimal, we have $X = \G \cdot U$, and, by compactness, there is a finite set $T \subset \G$ such that $X = T \cdot U$. But this is impossible as $(T \cdot U) \cap (\gamma \cdot U) = \0$ for all group elements $\gamma$ not in the finite set $T \, \Ret(U)$.
    %\end{scproof}

    %\begin{lemma}
    %    If $X$ and $Y$ are disjoint $\G$-flows, then at least one of $X$, $Y$ is minimal.
    %\end{lemma}
    %\begin{scproof}
    %    If $X' \subset X$, $Y' \subset Y$ are proper subflows, then \[Z \,\defeq\, (X \times Y') \cup (X' \times Y)\] is a joining distinct from $X \times Y$.
    %\end{scproof}

    %Here we prove Theorems~\ref{theo:disjoint_one} and \ref{theo:disjoint_many}.
    Let $X$ be a $\G$-flow. For a point $x \in X$ and a set $U \subseteq X$, the set of \emphd{visiting times} of $x$ to $U$ is
    \[
        \Vis(x,U) \,\defeq\, \set{\gamma \in \G \,:\, \gamma \cdot x \in U}.
    \]
    Let $\mathcal{F}(U)$ denote the family of all finite sets $F \subset \G$ such that $F \subseteq \Vis(x,U)$ for some $x \in U$. The bridge between disjointness and minimal subdynamics is given by the following proposition:

    \begin{prop}\label{prop:characterize}
        A minimal $\G$-flow $Y$ is disjoint from a $\G$-flow $X$ if and only if $Y$ is $\mathcal{F}(U)$-minimal for all nonempty open sets $U \subseteq X$.
    \end{prop}
    \begin{scproof}
        Suppose first that $Y$ is $\mathcal{F}(U)$-minimal for all nonempty open sets $U \subseteq X$. Let $Z \subseteq X \times Y$ be a joining and consider arbitrary nonempty open sets $U \subseteq X$ and $V \subseteq Y$. Since $Y$ is $\mathcal{F}(U)$-minimal, there is a set $F \in \mathcal{F}(U)$ such that $V$ is $F$-syndetic. Let $x \in U$ be a point satisfying $F \subseteq \Vis(x,U)$. Since $Z$ projects onto $X$, we can find $y \in Y$ such that $(x,y) \in Z$. As $V$ is $F$-syndetic, there is $\sigma \in F$ with $\sigma \cdot y \in V$. Since $F \subseteq \Vis(x,U)$, we have $\sigma \cdot x \in U$, and hence
        \[
            \sigma \cdot (x,y) \,=\, (\sigma \cdot x, \sigma \cdot y) \,\in\, (U \times V) \cap Z.
        \]
        It follows that $Z$ is dense in $X \times Y$. Since $Z$ is closed, we conclude that $Z = X \times Y$, as desired.

        Now suppose $Y$ fails to be $\mathcal{F}(U)$-minimal for some nonempty open set $U \subseteq X$. This means that there is a nonempty open set $V \subseteq Y$ that is not $F$-syndetic for any $F \in \mathcal{F}(U)$. We will construct a joining $Z$ of $X$ and $Y$ such that $Z \cap (U \times V) = \0$, showing that $Y$ is not disjoint from $X$.

        %For each $x \in X$, let $S_x \defeq \Vis(x,U)$.
        For each point $x\in U$, we have $\Vis(x,U)^{-1} \cdot V \neq Y$, for otherwise, by the compactness of $X$, there would exist a finite set $F \subseteq \Vis(x,U)$ such that $V$ is $F$-syndetic. Therefore, we may pick a point \[y_x \,\in\, Y \setminus \big(\Vis(x,U)^{-1} \cdot V\big).\] Consider the following $\G$-invariant (but not necessarily closed) subset $A \subseteq X \times Y$:
        \[
            A \,\defeq\, \set{(\gamma \cdot x, \gamma \cdot y_x) \,:\, x \in U, \, \gamma \in \G}.
        \]
        We claim that $A \cap (U \times V) = \0$. Indeed, if $(x,y) \in A$, then there exist $x' \in U$ and $\gamma \in \G$ such that
        \[x \,=\, \gamma \cdot x' \qquad \text{and} \qquad y \,=\, \gamma \cdot y_{x'}.\] If $x \in U$, then  $\gamma \in \Vis(x', U)$, and hence $\gamma \cdot y_{x'} \notin V$, as claimed. It follows that
        \[
            Z \,\defeq\, \overline{A} \,\cup\, \big((X \setminus (\G \cdot U)) \times Y\big)
        \]
        is a subflow of $X \times Y$ avoiding $U \times V$. The first coordinate projection of $A$ includes $\G \cdot U$, which shows that $Z$ projects onto $X$. The second coordinate projection of $Z$ is a subflow of $Y$, which must be equal to $Y$ since $Y$ is minimal. Therefore, $Z$ is a joining of $X$ and $Y$, and we are done.
    \end{scproof}

    Now we proceed to prove Theorem~\ref{theo:disjoint_many}. %First, a lemma:

    \begin{lemma}\label{lemma:recurrent}
        If $X$ is a %Polish
        $\G$-flow with no wandering points, then for every nonempty open set $U \subseteq X$, there exists a point $x \in U$ such that the set $\Vis(x,U)$ is infinite.
    \end{lemma}
    \begin{scproof}
        Suppose not. Then we have $U = \bigcup_{F} C_F$, where the union is over all finite sets $F \subset \G$ and $C_F$ is the set of all points $x \in U$ with $\Vis(x,U) \subseteq F$. Each set $C_F$ is relatively closed in $U$, so, by the Baire category theorem, at least one of them has nonempty interior. In other words, there exist a nonempty open set $W \subseteq U$ and a finite subset $F \subset \G$ such that $\Vis(x,U) \subseteq F$ for all $x \in W$. But then every point $x \in W$ satisfies $\Vis(x,W) \subseteq \Vis(x,U) \subseteq F$, which implies that $W \cap (\gamma \cdot W) = \0$ for all $\gamma \in \G \setminus F$. Hence, $W$ is a wandering nonempty open set, which is a contradiction.
    \end{scproof}

    %With Lemma~\ref{lemma:recurrent} in hand, we can prove Theorem~\ref{theo:disjoint_many}:

    \begin{scproof}[ of Theorem~\ref{theo:disjoint_many}]
        Let $(X_n)_{n \in \N}$ be a countable family of Polish $\G$-flows with no wandering points. We seek a free minimal subflow of $2^\G$ that is disjoint from every $X_n$. For each $n \in \N$, we fix a countable basis $(U_{n,m})_{m \in \N}$ for the topology on $X_n$ consisting of nonempty open sets. Lemma~\ref{lemma:recurrent} implies that each family $\mathcal{F}(U_{n,m})$ is unbounded. By Corollary~\ref{corl:subshift}, there is a free subflow $Y \subset 2^\G$ that is $\mathcal{F}(U_{n,m})$-minimal for all $n$, $m \in \N$; in particular, $Y$ is minimal. Note that if $U \subseteq X_n$ is a nonempty open set, then $Y$ is $\mathcal{F}(U)$-minimal because $\mathcal{F}(U) \supseteq \mathcal{F}(U_{n,m})$ for any basic open set $U_{n,m} \subseteq U$. Therefore, Proposition~\ref{prop:characterize} implies that $X_n \perp Y$ for all $n \in \N$, as desired.
    \end{scproof}

    
    
    
    Next we prove Theorem~\ref{theo:disjoint_one}, restated here for ease of reference:

    \begin{theocopy}{theo:disjoint_one}
        The following statements are equivalent for a Polish $\G$-flow $X$:
        \begin{enumerate}[label=\ep{\normalfont\arabic*}]
            \item\label{item:infinite} $X$ is disjoint from some infinite $\G$-flow,
            \item\label{item:free} $X$ is disjoint from some free minimal subflow of $2^\G$,
            \item\label{item:non-wandering} $X$ has no wandering points.
        \end{enumerate}
    \end{theocopy}
    \begin{scproof}
        Implication \ref{item:free} $\Longrightarrow$ \ref{item:infinite} is obvious, while \ref{item:non-wandering} $\Longrightarrow$ \ref{item:free} is given by Theorem~\ref{theo:disjoint_many}. It remains to verify \ref{item:infinite} $\Longrightarrow$ \ref{item:non-wandering}. Assume $X$ has a wandering nonempty open subset $U \subseteq X$ and let $F \subset \G$ be a finite set such that $U \cap (\gamma \cdot U) = \0$ for all $\gamma \in \G \setminus F$. Suppose $X$ is disjoint from a $\G$-flow $Y$. We need to show that $Y$ is finite.
        
        Since $X \perp Y$, at least one of $X$ and $Y$ is minimal, and $X$ cannot be minimal because minimal $\G$-flows have no wandering points. %\cite[Lem.~2.1]{LLLDyn3}.
        Therefore, $Y$ is minimal. By Proposition~\ref{prop:characterize}, it follows that $Y$ is $\mathcal{F}(U)$-minimal. Note that each set in $\mathcal{F}(U)$ is a subset of $F$ since $\Vis(x,U) \subseteq F$ for all $x \in U$. It follows that for any point $y \in Y$, the set $F \cdot y$ is dense in $Y$. Since finite sets in Hausdorff spaces are closed, this means that $Y = F \cdot y$, so it is finite.
    \end{scproof}

    To finish this section, we prove Proposition~\ref{prop:cosets}.

    \begin{scproof}[ of Proposition~\ref{prop:cosets}]
        Let $\Delta \leq \G$ be a subgroup and let $X$ be a $\G$-flow. We need to show that $X$ is disjoint from $(\G/\Delta)^*$ if and only if $X$ is $\Delta$-minimal. Note that $(\G/\Delta)^*$ is not minimal (the point at infinity forms a proper subflow), so if $X \perp (\G/\Delta)^*$, then $X$ must be minimal. Hence, by Proposition~\ref{prop:characterize}, $X \perp (\G/\Delta)^*$ if and only if $X$ is $\mathcal{F}(U)$-minimal for all nonempty open $U \subseteq (\G/\Delta)^*$.

        The set $\set{\Delta}$ is open and $\mathcal{F}(\set{\Delta}) = \finset{\Delta}$. Thus, if $X \perp (\G/\Delta)^*$, then $X$ is $\Delta$-minimal.

        Conversely, suppose $X$ is $\Delta$-minimal and consider any nonempty open set $U \subseteq (\G/\Delta)^*$. Pick an arbitrary point $y \in U$ and let $\Delta'$ be the stabilizer of $y$. Then $\Delta'$ is either a subgroup of $\G$ conjugate to $\Delta$ or, if $y$ is the point at infinity, $\G$ itself. In either case, $X$ is $\Delta'$-minimal (it is easy to see that a $\Delta$-minimal $\G$-flow is also minimal with respect to every conjugate of $\Delta$). Since $\Delta' \subseteq \Vis(y, U)$, we conclude that $X$ is $\mathcal{F}(U)$-minimal, as desired.
    \end{scproof}

    \section{Borel complete sections}\label{subsec:subshift_appl}

    %We shall prove the following common generalization of Theorems~\ref{theo:DST1} and \ref{theo:DST}. Say that a set $X \subseteq \Free(2^\G)$ is \emphd{wide} if for all $n \in \N$, there exist a finite set $F \subset \G$ of size $|F| = n$ and a Borel complete section $B \subseteq \Free(2^\G)$ such that $F \cdot B \subseteq X$. Let $\mathfrak{F}$ be the $\sigma$-filter generated by the wide sets; explicitly, a set $X \subseteq \Free(2^\G)$ belongs to $\mathfrak{F}$ if and only if there exist wide sets $(X_n)_{n \in \N}$ such that $X \supseteq \bigcap_{n \in \N} X_n$. 
    
    %\todo{...}
    
    Here we use Corollary~\ref{corl:subshift_prime} to prove Theorems~\ref{theo:DST1} and \ref{theo:DST}. We begin with a couple lemmas.

    \begin{lemma}\label{lemma:one_set}
        Let $\mathcal{F}$ be a family of finite subsets of $\G$ and let $X$ be an $\mathcal{F}$-minimal %Polish
        $\G$-flow. If $A \subseteq X$ is a non-meager Baire-measurable set, then for some $F  \in \mathcal{F}$, $F^{-1} \cdot A$ is comeager in $X$.
    \end{lemma}
    \begin{scproof}
        Let $U \subseteq X$ be the nonempty open set such that $U \symdif A$ is meager. By the $\mathcal{F}$-minimality of $X$, there is $F \in \mathcal{F}$ such that $F^{-1} \cdot U = X$, and hence $F^{-1} \cdot A$ is comeager, as desired.
    \end{scproof}

    \begin{lemma}\label{lemma:many_sets}
        Let $\mathcal{F} = \set{F_n}_{n \in \N}$ be a family of finite subsets of $\G$ and let $X$ be an $\mathcal{F}$-minimal %Polish
        $\G$-flow. If $(A_n)_{n \in \N}$ is a sequence of non-meager Baire-measurable subsets of $X$, then the set $\bigcup_{n \in \N} (F_n \cdot A_n)$ is comeager in $X$.
    \end{lemma}
    \begin{scproof}
        Suppose for contradiction that the set $M \defeq X \setminus \left(\bigcup_{n \in \N} (F_n \cdot A_n)\right)$ is non-meager. Then, by Lemma~\ref{lemma:one_set}, there is $n \in \N$ such that $F_n^{-1} \cdot M$ is comeager in $X$. This implies that $(F_n^{-1} \cdot M) \cap A_n \neq \0$, i.e., $M \cap (F_n \cdot A_n) \neq \0$, which contradicts the definition of $M$. 
    \end{scproof}

    %Now we can prove the following common generalization of Theorems~\ref{theo:DST1} and \ref{theo:DST}:

    %\begin{prop}\label{prop:DST}
    %    Suppose $(B_{n,m})_{n,m \in \N}$ is a family of Borel complete sections in $\Free(2^\G)$ and $(F_{n,m})_{n,m \in \N}$ is a collection of finite subsets of $\G$. Consider the set
    %\[
    %    B^* \,\defeq\, \bigcap_{m \in \N} \bigcup_{n \in \N} (F_{n,m} \cdot B_{n,m})
    %.\] If $\sup_{n \in \N} |F_{n,m}| = \infty$ for all $m \in \N$, then $B^*$ traps a point.
    %\end{prop}
    %\begin{scproof}
    %    For each $m \in \N$, let $\mathcal{F}_m \defeq \set{F_{n,m}}_{n \in \N}$. Since the families $\mathcal{F}_m$ are unbounded, Corollary~\ref{corl:subshift_prime} yields a free subshift $X \subset \Free(2^\G)$ that is $\mathcal{F}_m$-minimal for all $m \in \N$. Let $A_{n,m} \defeq B_{n,m} \cap X$. Since $B_{n,m}$ is a complete section, we have $\G \cdot A_{n,m} = X$, so $A_{n,m}$ is non-meager in $X$. By Lemma~\ref{lemma:many_sets}, %we conclude that for each $m \in \N$,
    %    the set $\bigcup_{n \in \N} (F_{n,m} \cdot A_{n,m})$ is comeager in $X$ for all $m \in \N$, and hence $B^*$ is comeager in $X$ as well. Then the set $\bigcap_{\gamma \in \G} (\gamma \cdot B^*)$ of all points $x \in X$ trapped by $B^*$ is also comeager in $X$; in particular, it is nonempty.
    %\end{scproof}

    %We are now ready to deduce Theorems~\ref{theo:DST1} and \ref{theo:DST}, restated here for ease of reference.

    %\begin{theocopy}{theo:DST1}\label{theo:DST1_copy}
    %    Suppose $(B_n)_{n \in \N}$ is a sequence of Borel complete sections in $\Free(2^\G)$ and $(F_n)_{n \in \N}$ is a sequence of finite subsets of $\G$ such that $\sup_{n \in \N}|F_n| = \infty$. Then there exists a point $x \in \Free(2^\G)$ satisfying $x \in F_n \cdot B_n$ for infinitely many $n \in \N$.
    %\end{theocopy}
    %\begin{scproof}
    %    Follows by applying Proposition~\ref{prop:DST} with $B_{n,m} \defeq B_{n+m}$ and $F_{n,m} \defeq F_{n+m}$. %For each $m \in \N$, let $\mathcal{F}_m \defeq \set{F_n}_{n \geq m}$. Since the families $\mathcal{F}_m$ are unbounded, Corollary~\ref{corl:subshift_prime} yields a free subshift $X \subset \Free(2^\G)$ that is $\mathcal{F}_m$-minimal for all $m \in \N$. Let $A_n \defeq B_n \cap X$. Since $B_n$ is a complete section, we have $\G \cdot A_n = X$, so $A_n$ is non-meager in $X$. By Lemma~\ref{lemma:many_sets}, we conclude that for each $m \in \N$, the set $\bigcup_{n \geq m} (F_n \cdot A_n)$ is comeager in $X$. Thus, the set $\bigcap_{m \in \N} \bigcup_{n \geq m} (F_n \cdot A_n)$ of all points $x \in X$ satisfying $x \in F_n \cdot B_n$ for infinitely many $n \in \N$ is also comeager in $X$. In particular, this set is nonempty. 
    %\end{scproof}

    %\begin{theocopy}{theo:DST}
    %    If $B \subseteq \Free(2^\G)$ is a Borel complete section, then there exists $n \in \N$ such that for every finite set $F \subset \G$ of size at least $n$, the set $F \cdot B$ traps a point.
    %\end{theocopy}
    %\begin{scproof}
    %    Suppose for contradiction that for each $n \in \N$, there exists a set $F_n \subset \G$ of size $|F_n| = n$ such that $F_n \cdot B$ does not trap a point, i.e., the set $B_n \defeq \Free(2^\G) \setminus (F_n \cdot B)$ is a complete section. %Note that $F_n^{-1} \cdot B_n \subseteq \Free(2^\G) \setminus B$.
    %    By Proposition~\ref{prop:DST}, the set $B^* \defeq \bigcup_{n \in \N} (F_n^{-1} \cdot B_n)$ traps a point. But $B \cap B^* = \0$, contradicting the fact that $B$ is a complete section.
    %\end{scproof}

    %\todo{...}

    We are now ready to deduce Theorems~\ref{theo:DST1} and \ref{theo:DST}, restated here for ease of reference.

    \begin{theocopy}{theo:DST1}\label{theo:DST1_copy}
        Suppose $(B_n)_{n \in \N}$ is a sequence of Borel complete sections in $\Free(2^\G)$ and $(F_n)_{n \in \N}$ is a sequence of finite subsets of $\G$ such that $\sup_{n \in \N}|F_n| = \infty$. Then there exists a point $x \in \Free(2^\G)$ satisfying $x \in F_n \cdot B_n$ for infinitely many $n \in \N$.
    \end{theocopy}
    \begin{scproof}
        For each $m \in \N$, let $\mathcal{F}_m \defeq \set{F_n}_{n \geq m}$. Since the families $\mathcal{F}_m$ are unbounded, Corollary~\ref{corl:subshift_prime} yields a free subshift $X \subset \Free(2^\G)$ that is $\mathcal{F}_m$-minimal for all $m \in \N$. Let $A_n \defeq B_n \cap X$. Since $B_n$ is a complete section, we have $\G \cdot A_n = X$, so $A_n$ is non-meager in $X$. By Lemma~\ref{lemma:many_sets}, we conclude that for each $m \in \N$, the set $\bigcup_{n \geq m} (F_n \cdot A_n)$ is comeager in $X$. Thus, the set $\bigcap_{m \in \N} \bigcup_{n \geq m} (F_n \cdot A_n)$ of all points $x \in X$ satisfying $x \in F_n \cdot B_n$ for infinitely many $n \in \N$ is also comeager in $X$. In particular, this set is nonempty. 
    \end{scproof}

    \begin{theocopy}{theo:DST}
        If $B \subseteq \Free(2^\G)$ is a Borel complete section, then there exists $n \in \N$ such that for every finite set $F \subset \G$ of size at least $n$, the set $F \cdot B$ traps a point.
    \end{theocopy}
    \begin{scproof}
        Suppose for contradiction that for each $n \in \N$, there exists a set $F_n \subset \G$ of size $|F_n| = n$ such that $F_n \cdot B$ does not trap a point. Let $\mathcal{F} \defeq \set{F_n^{-1}}_{n \in \N}$ and consider a free $\mathcal{F}$-minimal subshift $X \subset \Free(2^\G)$. As $B$ is a complete section, the set $A \defeq B \cap X$ is non-meager in $X$, so, by Lemma~\ref{lemma:one_set}, $F_n \cdot A$ is comeager in $X$ for some $n \in \N$. Since $F_n \cdot A$ does not trap a point by assumption, the set $M \defeq X \setminus (F_n \cdot A)$ must meet every orbit in $X$, which is impossible as $M$ is meager in $X$.
    \end{scproof}

    Note that our proof of Theorem~\ref{theo:DST1_copy} actually yields the following somewhat stronger statement: For any family $(B_{n,m})_{n,m \in \N}$ of Borel complete sections in $\Free(2^\G)$ and any collection $(F_{n,m})_{n,m \in \N}$ of finite subsets of $\G$ satisfying $\sup_{n \in \N} |F_{n,m}| = \infty$ for all $m \in \N$, the following set traps a point:
    \[
        \bigcap_{m \in \N} \bigcup_{n \in \N} (F_{n,m} \cdot B_{n,m}).
    \]


    

    

    \subsection*{Acknowledgments}

    We are grateful to Be'eri Greenfeld, Gil Goffer, and Tianyi Zheng for organizing the AMS Special Session on \emph{Combinatorial and Probabilistic Methods in Group Theory} at the 2025 Joint Mathematics Meetings (January 8--11, 2025) in Seattle, where this project began, and to Lewis Bowen for providing a stimulating environment during the UT Austin Mini-school in \emph{Analytic Group Theory} (March 17--21, 2025), where part of our work on this project took place. We also thank Andy Zucker for insightful conversations. This material is based upon work partially supported by the Alfred P.~Sloan Foundation and the National Science Foundation under grants DMS-2528522 and DMS-2348981. Any opinions, findings, and conclusions or recommendations expressed in this material are those of the authors and do not necessarily reflect the views of the National Science Foundation.


\printbibliography



\end{document}
